\chapter{Subintegral normalization in characteristic zero}
\index{subintegral extension!envelope theorem}\index{normalization!subintegral}
\label{chap:bj-subintegral-normalization}

\section{Orientation}
\label{sec:ch44-orientation}

\subsection{What is proved here, and out of what}
\label{sec:ch44-what}

\Cref{chap:bj-universal-homeomorphisms} proved two things about a finite
birational universal homeomorphism \(\nu:Y\to X\) with \(Y\) normal.  The
first is topological, holds in every characteristic, and is already
available here: \(\nu^{\An}\) is a homeomorphism, so a function on
\(X^{\An}\) and its pullback to \(Y^{\An}\) are the same function of the
same points, at every semivaluation, including those of proper support and
including where the value is \(-\infty\)
(\Cref{lem:bj-universal-analytic-homeomorphism}).  The second is algebraic
and was proved only for \(\operatorname{char}k=p>0\): an upstairs
generically finite Fubini--Study presentation can be converted into a
downstairs one without changing the function, by raising every local
coefficient to a single Frobenius power
(\Cref{cor:ch43-frobenius-all-affines},
\Cref{prop:bj-frobenius-gf-descent}).

The present chapter supplies the missing algebraic half.  We retain the
setup
\[
                         \nu:Y\longrightarrow X
\tag{44.1}\label{eq:bj-char-zero-universal-map}
\]
of \Cref{chap:bj-universal-homeomorphisms}: \(\nu\) is a finite birational
universal homeomorphism of integral projective varieties over the
algebraically closed trivially valued field \(k\), and \(Y\) is normal;
but now \(\operatorname{char}k=0\).  The analytic spaces have already been
identified by \Cref{lem:bj-universal-analytic-homeomorphism}, and nothing in
that identification used the characteristic, so the whole of the topological
apparatus of \Cref{sec:ch43-analytic} is at our disposal without repetition.

The obstruction is stated in one line.  In characteristic \(p\) the set
\(\{b\in B:b^q\in A\}\) is an \(A\)-subalgebra of \(B\), because
\(x\mapsto x^q\) is a ring homomorphism; in characteristic zero it is not
even closed under addition, and there is no exponent \(q\) whatsoever with
\(b^q\in A\) for all \(b\in B\), already for
\(A=\mathbf Q[t^2,t^3]\subset B=\mathbf Q[t]\)
(\Cref{rem:ch43-charzero-contrast}).  Raising to a power is simply not an
available operation.  What replaces it is the additive-power theorem of
Roberts and Singh: for a subintegral extension of \(\mathbf Q\)-algebras,
the failure of \(b^m\) to lie in \(A\) can be corrected by lower-order terms
whose coefficients \(c_1,\ldots,c_p\in B\) are \emph{fixed}, independent of
the exponent \(m\) (\Cref{thm:appA-roberts-singh}).

That single change of algebraic input propagates through everything
downstream, and the reader should have the propagation in view before the
first proof.

\subsection{The structural contrast with the positive-characteristic proof}
\label{sec:ch44-contrast}

The two chapters prove the same theorem by arguments whose shapes differ at
every joint.  We tabulate the difference once, because each row is a
separate technical obligation discharged below.

\begin{center}
\begin{tabular}{p{.26\textwidth}p{.32\textwidth}p{.32\textwidth}}
\toprule
 & \(\operatorname{char}k=p>0\)
   (\Cref{chap:bj-universal-homeomorphisms})
 & \(\operatorname{char}k=0\) (this chapter)\\
\midrule
algebraic input &
one exponent \(q=p^e\) with \(b^q\in A\) for every \(b\) on every affine
open (\Cref{cor:ch43-frobenius-all-affines}) &
for each fixed \(b\), correction terms \(c_1,\ldots,c_p\in B\) independent
of \(m\), with \(b^m+\sum_i\binom mic_ib^{m-i}\in A\) for \(m\gg0\)
(\Cref{thm:appA-roberts-singh})\\
what is descended &
each section \(s_j\) separately, as \(s_j^{[q]}\) with
\(\nu^*s_j^{[q]}=s_j^{\otimes q}\) &
no single section; a finite \emph{family} \(\mathfrak K_n\) of weighted
sections attached to the whole system at once\\
form of the conclusion &
an \emph{identity} of functions,
\(\psi=\frac1{qm}\max_j\{\log|s_j^{[q]}|+q\lambda_j\}\) &
a two-sided \emph{estimate},
\(nW(\mathfrak a)\le W(\mathfrak K_n\cO_Y)\le(n+C)W(\mathfrak a)\), with
\(C\) independent of \(n\)\\
number of stages &
one; the descended function is already the given one &
infinitely many; the given function is a limit as \(n\to\infty\)\\
monotonicity &
automatic: there is nothing to monotonize &
manufactured, by rescaling \(F_{n_j}\mapsto a_jF_{n_j}\) with
\(a_j\uparrow1\) chosen against \(r_{n_j}=n_j/(n_j+C)\)\\
class of the approximants &
exactly \(L\) &
\(a_jL+\frac{a_j\rho}{n_j+C}H\), converging to \(L\); a fixed twist \(H\)
divided by a growing integer\\
closure property used &
none beyond \Cref{def:bj-psh} &
\Cref{thm:bj-basic-closure}\,(5), decreasing nets with \emph{varying}
classes\\
what would break it &
loss of the subalgebra property, i.e.\ characteristic zero &
loss of uniformity of \(C\), or of the single twist \(H\), in \(n\)\\
\bottomrule
\end{tabular}
\end{center}

Two entries deserve emphasis.  First, the passage from an identity to an
estimate is not a weakening that one then has to repair: a two-sided
estimate with \emph{additive} loss is exactly as good as an identity for the
purpose at hand, because the ratio \(r_n=n/(n+C)\) tends to \(1\).  Had the
loss been multiplicative --- had the sandwich read
\(nW(\mathfrak a)\le W(\mathfrak K_n\cO_Y)\le \kappa n\,W(\mathfrak a)\)
with \(\kappa>1\) --- the limit of the approximants would have been
\(\kappa^{-1}\psi_0\) and not \(\psi_0\), and the argument would prove
nothing.  The whole of \Cref{sec:ch44-sandwich} exists to make the loss
additive, and \Cref{ex:ch44-splitting-fails} exhibits a natural-looking
variant of the construction in which it is multiplicative instead.

Second, the manufactured monotonicity is cheap but not free.  It costs a
factor \(a_j<1\) on the class, which is why the classes of the approximants
move; and that is affordable only because \Cref{def:bj-psh} permits the
classes of a regularizing net to vary.  A reader who has internalized a
fixed-class definition of \(\PSH\) will find this step illegal.  It is not;
see \Cref{sec:ch44-hypotheses}.

\subsection{The two points that require care}
\label{sec:ch44-care}

Two steps of the construction are easy to get wrong in ways that leave the
final statement looking correct.  Both are isolated below as displayed
steps, not as remarks, and each is accompanied by an explicit
counterexample to the wrong version.

\medskip\noindent\textbf{(C1) The additive-power identity is applied to the
full generic linear combination, before any weight component is split off.}
The Fubini--Study weights are encoded by an auxiliary variable \(u\), so the
generators of the upstairs system are
\[
 g_j=u^{d_j}f_j,
 \qquad
 d_j\in\mathbf N,
\]
which are individually \(u\)-homogeneous but of \emph{different} degrees.
The element to which \Cref{thm:appA-roberts-singh} is applied is the full
generic combination
\[
 b=\sum_{j=1}^sz_jg_j\in B[u,\mathbf z],
\]
which is therefore \(u\)-inhomogeneous.  This is legitimate because
subintegrality is preserved by polynomial base change,
\[
 A\subset B
 \quad\Longrightarrow\quad
 A[u,\mathbf z]\subset B[u,\mathbf z]
 \quad\text{subintegral},
\]
which is Tags 0CND and 0EUH of the Stacks Project, audited in
\Cref{sec:stacks-universal-homeomorphisms} and recorded here as
\Cref{lem:ch44-polynomial-base-change}.  Only after the identity has been
formed, and only on its \(\mathbf z\)-coefficients, are \(u\)-homogeneous
components extracted.  \Cref{ex:ch44-splitting-fails} shows what happens
otherwise: splitting the correction coefficients \(c_i\) by \(u\)-degree
destroys membership in \(A[u,\mathbf z]\) at once, and applying the identity
separately to the weight components loses every mixed monomial from the
leading-power relation, degrading the sandwich exponent from \(n+2p\) to
\(s(n+2p)\).

\medskip\noindent\textbf{(C2) All local data are globalized with a single
twist, independent of \(n\).}  Three separate extension problems occur:
extending the fixed polynomials \(c_{\alpha i}\) from a chart to all of
\(Y\); killing the class of a block in the coherent quotient
\(\mathcal Q=\nu_*\cO_Y/\cO_X\); and aligning blocks of different degrees to
a common degree.  Each is solved by twisting.  The three twists are
\[
 R\,A_0,
 \qquad
 S\,A_0,
 \qquad
 H_1,
\]
and \emph{all three are chosen before \(n\) is chosen}: \(R\) because only
finitely many elements of \(B_\alpha\) occur as coefficients of the fixed
\(c_{\alpha i}\); \(S\) because a coherent subsheaf of \(\mathcal Q\) has a
single annihilating power, unchanged by twisting by an invertible sheaf; and
\(H_1\) because only the \(C+1\) residual degrees \(0\le j\le C\) occur, and
\(C=2P\) is determined by the local systems.  The resulting class error is
\[
 \frac{\rho}{n+C}\,c_1(H),
 \qquad
 H:=(R+S)A_0+H_1,
\]
which tends to \(0\) precisely because \(H\) does not move with \(n\).
\Cref{warn:ch44-moving-twist} computes what a twist growing linearly in
\(n\) would do: the classes of the approximants would converge to
\(c_1(L)+\rho\,c_1(A_0)\ne c_1(L)\), and the construction would produce a
psh function in the wrong class.

\subsection{The firewall}
\label{sec:ch44-firewall}

\begin{warning}{The envelope theorem is not available in this chapter}
\label{warn:ch44-firewall}
Nothing in \Cref{chap:bj-subintegral-normalization} may use, directly or
transitively:
\begin{enumerate}
\item \Cref{thm:bj-final-envelope}, in any of its formulations, nor
      \Cref{def:bj-envelope-property-final}, nor
      \Cref{cor:bj-historical-conjecture};
\item the envelope property of \Cref{def:bj-envelope-property} as a
      \emph{theorem} for any class on any variety; the symbol does not
      occur below outside \Cref{sec:ch44-unibranch}, where it appears in a
      statement about what \emph{remains} to be assembled in
      \Cref{chap:bj-envelope-consequences};
\item closure of any \(\PSH_X(\theta)\) under suprema, regularized or not;
\item any statement of \Cref{chap:bj-envelope-consequences}, and none of
      \Cref{cor:ch42-birational-bijection},
      \Cref{cor:ch42-finite-bijection}, \Cref{cor:ch41-finite-bijection},
      which are consequences of \Cref{thm:bj-normal-envelope}.
\end{enumerate}
\Cref{thm:bj-char-zero-universal-reflection} is an \emph{input} to
\Cref{thm:bj-final-envelope}, not a consequence of it, and the same holds
for every numbered statement below.
\end{warning}

\begin{remark}[Where the firewall bites]
\label{rem:ch44-firewall-not-vacuous}
The temptations are the same two as in
\Cref{rem:ch43-firewall-not-vacuous}, and they are worth re-examining
because the characteristic-zero proof passes closer to both.

First, surjectivity of \(\nu^*\) on psh functions.  For a birational
morphism of \emph{normal} projective varieties this is
\Cref{cor:ch42-birational-bijection}; but \(X\) is not normal here --- if it
were, \(\nu\) would be an isomorphism --- and in any case that corollary is
a consequence of \Cref{thm:bj-normal-envelope} and may not be used to
produce an input to the unibranch theorem.  Surjectivity is obtained below
from scratch, by exhibiting for each upstairs GF stage an explicit
decreasing sequence of downstairs GF functions converging to it.

Second, and this is the delicate one, the last step of
\Cref{prop:bj-char-zero-stage-reflection} takes a \emph{decreasing} sequence
\(a_jF_{n_j}\) of downstairs GF functions with classes converging to
\(c_1(L)\) and concludes that the limit is \(c_1(L)\)-psh.  That is
\Cref{thm:bj-basic-closure}\,(5), a definitional closure property of
\(\PSH\) proved in \Cref{chap:bj-fs-pl-psh}, entirely independent of any
envelope statement.  It is emphatically \emph{not} closure under increasing
families.  The place where care is needed is a different one.  The sandwich
\eqref{eq:bj-char-zero-function-sandwich} gives
\(\psi_0\le F_n\le r_n\psi_0\) with \(\psi_0\le0\) and \(r_n\uparrow1\), so
\(r_n\psi_0\downarrow\psi_0\) and therefore
\[
 \psi_0=\lim_nF_n=\inf_nF_n .
\]
It is tempting to stop here and quote a closure property for infima.  But
\Cref{thm:bj-basic-closure}\,(5) is a statement about \emph{decreasing
nets}, and the raw family \((F_n)_n\) is not monotone in either direction:
the sandwich controls each \(F_n\) between two moving barriers but says
nothing about \(F_n\) versus \(F_{n+1}\).  An infimum of a non-monotone
family of psh functions need not be psh, and no statement of this book
asserts that it is.  The monotonization
\eqref{eq:bj-char-zero-monotone-chain} is performed by hand, on a
subsequence, and produces a genuinely \emph{decreasing} family before
\Cref{thm:bj-basic-closure}\,(5) is applied.  The alternative
\(\psi\equiv-\infty\) in that statement is excluded each time by recording
the value at \(v_\triv\).
\end{remark}

\subsection{The five tasks, and where each is discharged}
\label{sec:ch44-tasks}

\begin{enumerate}
\item\label{itm:ch44-task:subintegral}
      \emph{Subintegral and seminormal extensions.}  Identify
      \(\cO_X(U)\subset\cO_Y(\nu^{-1}U)\) as a finite subintegral extension
      in characteristic zero, record the square--cube criterion and its
      stability under polynomial base change, and separate subintegrality
      from seminormality --- the two are complementary, and their
      interaction explains why the hypothesis ``\(X\) unibranch and not
      normal'' is not vacuous.  This is \Cref{sec:ch44-subintegral}, with
      \Cref{lem:ch44-char-zero-subintegral},
      \Cref{lem:ch44-polynomial-base-change} and
      \Cref{prop:ch44-seminormal-unibranch-normal}.
\item\label{itm:ch44-task:rs}
      \emph{The Roberts--Singh additive-power sandwich.}  State the
      elementwise exponent-zero theorem, compute its data explicitly on the
      cusp and on a curve with two gaps, and prove the local sandwich
      \eqref{eq:bj-local-additive-sandwich}.  This is
      \Cref{sec:ch44-roberts-singh} and \Cref{sec:ch44-sandwich}, with
      \Cref{lem:ch44-finite-differences} and
      \Cref{lem:bj-local-additive-sandwich}.
\item\label{itm:ch44-task:uniform}
      \emph{Uniform control of homogeneous functions.}  Encode a weighted
      GF function as the Gauss weight of a homogeneous ideal, globalize the
      local systems with a single twist, and prove the global two-sided
      estimate \eqref{eq:bj-char-zero-global-weight-sandwich}.  This is
      \Cref{sec:ch44-encoding}, \Cref{sec:ch44-globalization} and
      \Cref{sec:ch44-global-sandwich}, with
      \Cref{lem:ch44-gauss-valuation}, \Cref{lem:ch44-weight-encoding},
      \Cref{lem:ch44-uniform-cutoffs} and
      \Cref{prop:ch44-global-sandwich}.
\item\label{itm:ch44-task:reflect}
      \emph{Psh reflection in characteristic zero.}  Convert the estimate
      into a decreasing net and prove
      \Cref{prop:bj-char-zero-stage-reflection} and
      \Cref{thm:bj-char-zero-universal-reflection}.  This is
      \Cref{sec:ch44-decreasing} and \Cref{sec:ch44-reflection}.
\item\label{itm:ch44-task:unibranch}
      \emph{Completion of the unibranch normalization step.}  Combine with
      the positive-characteristic theorem to obtain
      \Cref{cor:bj-universal-reflection-all-characteristics}, apply it to
      the normalization of a unibranch variety, and state exactly what
      remains for \Cref{chap:bj-envelope-consequences}.  This is
      \Cref{sec:ch44-unibranch}.
\end{enumerate}

\subsection{Standing hypotheses and notation}
\label{sec:ch44-hypotheses}

Throughout, \(k\) is an algebraically closed field of characteristic zero
carrying the trivial absolute value, and \(\nu:Y\to X\) is as in
\eqref{eq:bj-char-zero-universal-map}.  The hypothesis
\(\operatorname{char}k=0\) is used in exactly three places, and it is worth
naming them now: to know that the residue extensions of \(\nu\) are trivial,
hence that \(\nu\) is subintegral and not merely radicial
(\Cref{lem:ch44-char-zero-subintegral}); to apply
\Cref{thm:appA-roberts-singh}, which requires \(\mathbf Q\)-algebras; and to
invert the multinomial coefficients \(\binom{n+2p}{\beta}\) at the end of
the proof of \Cref{lem:bj-local-additive-sandwich}.  Nothing else in the
chapter is characteristic-sensitive.

Numerical classes are real unless stated otherwise, and
\(N^1(X)=N^1(X)_{\mathbf R}\); a rational line bundle is an element of
\(\Pic(X)_{\mathbf Q}\).  The analytic vocabulary is that of
\Cref{chap:trivial-analytification}: for \(v\in X^{\An}\) we write \(s(v)\)
for its support (\Cref{def:triv-support}) and \(c(v)\) for its center, and a
point with \(s(v)\subsetneq X\) has proper support
(\Cref{def:ch32-proper-support}).  Sections and their logarithms are
normalized as in \Cref{chap:bj-fs-pl-psh}: \(\log|s|(v)=-v(s)\), computed in
any local frame at \(c(v)\), with the value \(-\infty\) permitted and never
discarded (\Cref{lem:ch34-frame-independence}).  The sign conventions of
\Cref{app:normalizations} are in force without exception.

We write \(\cH^{\mathrm{gf}}_{\mathbf Q}(L)\) for the rational generically
finite Fubini--Study functions attached to \(L\in\Pic(X)_{\mathbf Q}\)
(\Cref{def:gf-fs-functions}) and \(\PSH_X(\theta)\) for the \(\theta\)-psh
functions of \Cref{def:bj-psh}.  We recall once, because the chapter turns
on it, that the approximating net in that definition consists of GF
functions \(\phi_i\in\cH^{\mathrm{gf}}_{\mathbf R}(L_i)\) with
\emph{varying} rational classes \(c_1(L_i)\to\theta\), and that the
consequence \(\PSH_X(\theta)\ne\emptyset\Rightarrow\theta\) pseudoeffective
--- with nefness required only when a \emph{bounded} member is asserted to
exist --- is \Cref{lem:ch34-psh-nonempty}.  No statement below asserts that
any class occurring in it is nef.

For an affine open \(U\subseteq X\) we abbreviate
\[
 A_U:=\cO_X(U),
 \qquad
 B_U:=\cO_Y(\nu^{-1}U),
\]
noting that \(\nu^{-1}U\) is affine, \(\nu\) being finite, and that
\(A_U\to B_U\) is an injective finite ring map, \(\nu\) being finite,
birational and dominant between integral schemes.  Formal variables are
written \(u\) (a single variable, carrying the weights) and
\(\mathbf z=(z_1,\ldots,z_s)\) (carrying the coefficient extraction); the
letter \(p\) denotes the Roberts--Singh length, never a characteristic,
which is zero throughout.

\subsection{What is used from earlier chapters}
\label{sec:ch44-dependencies}

The imports are listed so that \Cref{warn:ch44-firewall} can be checked by
inspection.

From \Cref{chap:bj-universal-homeomorphisms}: the analytic homeomorphism
\Cref{lem:bj-universal-analytic-homeomorphism} and its dictionary
\Cref{lem:ch43-identification-dictionary}; the class-promotion lemma
\Cref{lem:bj-normalization-class-promotion}; the identification of the
normalization of a unibranch variety
\Cref{lem:bj-unibranch-normalization-universal}; and, for the assembly of
\Cref{cor:bj-universal-reflection-all-characteristics}, the
positive-characteristic reflection theorem
\Cref{thm:bj-positive-universal-reflection}.  All four are proved there
without any characteristic hypothesis except the last, and none of them uses
an envelope statement.

From \Cref{chap:bj-fs-pl-psh}: the definition of
\(\cH^{\mathrm{gf}}_\Lambda(L)\) and its elementary calculus
(\Cref{def:gf-fs-functions}, \Cref{prop:ch34-fs-operations},
\Cref{lem:gf-fs-operations}); the definition of \(\PSH_X(\theta)\)
(\Cref{def:bj-psh}); the existence thresholds
\Cref{lem:ch34-psh-nonempty}; and the closure package
\Cref{thm:bj-basic-closure}, of which only items~(1) and~(5) are used.

From \Cref{chap:bj-finite-descent}: nothing.  In particular there is no
determinant norm, no flattening, and no division by a fibre degree anywhere
below; \(\nu\) has singleton fibres and degree one, so there is nothing to
average (\Cref{rem:ch43-no-averaging}).

From \Cref{chap:bj-alterations}: nothing is used in the proofs.  That
chapter supplies \Cref{thm:bj-normal-envelope}, which will be combined with
the present theorem in \Cref{chap:bj-envelope-consequences}; the combination
is performed there.

From \Cref{app:valuation-algebra}: the definition of a subintegral extension
(\Cref{def:appA-subintegral}), of seminormality
(\Cref{def:appA-seminormal}), the square--cube criterion
(\Cref{thm:appA-square-cube}), the universal-homeomorphism criterion
(\Cref{thm:appA-universal-homeomorphism-criterion}), and the Roberts--Singh
theorem in the form \Cref{thm:appA-roberts-singh} together with its generic
corollary \Cref{cor:appA-RS-generic}.  From \Cref{app:stacks}: the tags
0CND and 0EUH, audited in \Cref{sec:stacks-universal-homeomorphisms} with
their exact statements and hypotheses.  No tag is quoted below that is
absent from that audit, and none is invented.

\subsection{Roadmap}
\label{sec:ch44-roadmap}

\Cref{sec:ch44-subintegral} sets up subintegrality and seminormality and
proves the structural facts the chapter needs.
\Cref{sec:ch44-roberts-singh} states the additive-power theorem and computes
its data on two curves.  \Cref{sec:ch44-sandwich} proves the local sandwich;
this is the technical heart.  \Cref{sec:ch44-encoding} encodes a weighted GF
function as a Gauss weight, \Cref{sec:ch44-globalization} globalizes the
local systems with a single twist, and \Cref{sec:ch44-global-sandwich}
assembles the global two-sided estimate.
\Cref{sec:ch44-decreasing} manufactures the decreasing net,
\Cref{sec:ch44-reflection} proves the reflection theorem, and
\Cref{sec:ch44-unibranch} completes the normalization step.
\Cref{sec:ch44-computations} carries the theory to explicit numbers on the
cuspidal cubic and on a nonnormal unibranch surface,
\Cref{sec:ch44-audit} performs the extended-real and circularity audit,
\Cref{sec:ch44-failure} maps the failure boundaries, and
\Cref{sec:ch44-dictionary} and \Cref{sec:ch44-source-note} close the
chapter.

\section{Subintegral extensions and seminormality}
\label{sec:ch44-subintegral}

Subintegrality is the algebraic shadow of ``bijective with trivial residue
extensions''.  It is strictly stronger than radiciality, which allows purely
inseparable residue extensions and is therefore the correct notion in
characteristic \(p\); in characteristic zero the two coincide for finite
maps, and that coincidence is what makes the present chapter possible at
all.

\subsection{Definitions and the square--cube criterion}
\label{sec:ch44-defs}

\begin{definition}[Subintegral extension]
\label{def:bj-subintegral-extension}
An integral extension \(A\subset B\) is \emph{subintegral} if
\(\Spec B\to\Spec A\) is bijective and every induced residue-field extension
is an isomorphism.
\end{definition}

\begin{definition}[Seminormality in an overring]
\label{def:bj-seminormal-overring}
The ring \(A\) is \emph{seminormal in \(B\)} if every \(b\in B\) satisfying
\(b^2,b^3\in A\) already belongs to \(A\).
The maximal subintegral intermediate algebra, when viewed inside a fixed
integral overring, is the subintegral closure or seminormalization of \(A\)
in that overring.
\end{definition}

\Cref{def:bj-subintegral-extension} agrees with \Cref{def:appA-subintegral}
and \Cref{def:bj-seminormal-overring} with \Cref{def:appA-seminormal}; they
are restated here because the two notions are used against each other
repeatedly below and it is important that the reader has both in view at
once.  A ring is \emph{seminormal} without qualification if it is seminormal
in its normalization.

The square--cube criterion makes subintegrality constructive.
For an extension \(A\subset B\), the condition that
\(\Spec B\to\Spec A\) be a universal homeomorphism inducing residue-field
isomorphisms is equivalent to saying that every finite subset of \(B\) lies
in a finite tower obtained by adjoining elements whose squares and cubes
belong to the preceding ring.
This is
\cite[\href{https://stacks.math.columbia.edu/tag/0CND}{Tag 0CND}]
{stacks-project},
imported into the book as \Cref{thm:appA-square-cube} and audited, with the
exact statement and the hypotheses to be checked, in
\Cref{sec:stacks-universal-homeomorphisms}.
Such elementary extensions and their residue-field property persist under
base change by
\cite[\href{https://stacks.math.columbia.edu/tag/0EUH}{Tag 0EUH}]
{stacks-project}.

The consequence that the chapter actually consumes is the following; it is
point~\textbf{(C1)} of \Cref{sec:ch44-care}, and it is stated as a lemma so
that its use can be pointed at.

\begin{lemma}
\label{lem:ch44-polynomial-base-change}
Let \(A\subset B\) be a subintegral extension of rings and let
\(\mathbf w=(w_1,\ldots,w_r)\) be finitely many indeterminates.  Then
\[
 A[\mathbf w]\subset B[\mathbf w]
\]
is again a subintegral extension.  If \(A\subset B\) is moreover finite,
then so is \(A[\mathbf w]\subset B[\mathbf w]\).
\end{lemma}

\begin{proof}
By \Cref{thm:appA-square-cube} it suffices to verify the finite-stage
square--cube condition for \(A[\mathbf w]\subset B[\mathbf w]\).  Let
\(F_1,\ldots,F_M\in B[\mathbf w]\) be finitely many polynomials and let
\(b_1,\ldots,b_M'\in B\) be the finitely many coefficients occurring in
them.  By \Cref{thm:appA-square-cube} applied to \(A\subset B\) there is a
tower
\[
 A=C_0\subset C_1\subset\cdots\subset C_r\subset B,
 \qquad
 C_l=C_{l-1}[x_l],
 \qquad
 x_l^2,x_l^3\in C_{l-1},
\]
containing all the \(b_\bullet\).  Adjoining the same elements over
\(A[\mathbf w]\) produces the tower
\[
 A[\mathbf w]=C_0[\mathbf w]\subset C_1[\mathbf w]\subset\cdots
 \subset C_r[\mathbf w]\subset B[\mathbf w],
\]
in which \(C_l[\mathbf w]=C_{l-1}[\mathbf w][x_l]\) and
\(x_l^2,x_l^3\in C_{l-1}[\mathbf w]\).  It contains every \(F_i\), since it
contains all their coefficients and all the \(w_\bullet\).  This is exactly
the finite-stage condition, so \(A[\mathbf w]\subset B[\mathbf w]\) is
subintegral.  Alternatively, and this is the formulation audited at Tag
0EUH, the square--cube generation of \(B\) over \(A\) is preserved by the
base change \(A\to A[\mathbf w]\), which is what that tag asserts; the
displayed tower is the same statement made explicit.

Finiteness is clear: if \(b_1,\ldots,b_M\) generate \(B\) as an
\(A\)-module, they generate \(B[\mathbf w]\) as an
\(A[\mathbf w]\)-module.
\end{proof}

\begin{warning}{Tag 0EUH does not supply the additive-power theorem}
\label{warn:ch44-tags-scope}
\Cref{lem:ch44-polynomial-base-change} is a statement about square--cube
towers and nothing more.  It says that the \emph{hypothesis} of
\Cref{thm:appA-roberts-singh} survives the passage from \(A\subset B\) to
\(A[u,\mathbf z]\subset B[u,\mathbf z]\).  It says nothing about how the
data \(p\), \(N\), \(c_1,\ldots,c_p\) behave, and in particular it does not
assert that they may be chosen over \(A\) and then extended: they are
produced afresh, over \(A[u,\mathbf z]\), by
\Cref{cor:appA-RS-generic}, and they genuinely involve \(u\) and
\(\mathbf z\).  This is recorded because it is exactly the point at which a
reader may be tempted to shortcut the argument by applying the additive-power
theorem coefficientwise; \Cref{ex:ch44-splitting-fails} shows that this
does not work.
\end{warning}

\subsection{Why the extension at hand is subintegral}
\label{sec:ch44-why-subintegral}

On an affine chart \(U\subset X\), put
\[
                         A=\cO_X(U),
             \qquad
                         B=\cO_Y(\nu^{-1}U).
\tag{44.2}\label{eq:bj-char-zero-affine-extension}
\]
\begin{lemma}
\label{lem:ch44-char-zero-subintegral}
Assume \(\operatorname{char}k=0\) and let \(\nu:Y\to X\) be as in
\eqref{eq:bj-char-zero-universal-map}.  Then for every affine open
\(U\subseteq X\) the extension \eqref{eq:bj-char-zero-affine-extension} is
finite and subintegral, and \Cref{lem:ch44-polynomial-base-change} applies
to it.
\end{lemma}

\begin{proof}
Finiteness: \(\nu\) is finite, so \(\nu^{-1}U\) is affine and \(B_U\) is a
finite \(A_U\)-module; injectivity of \(A_U\to B_U\) holds because \(\nu\)
is dominant between integral schemes.

Bijectivity of \(\Spec B_U\to\Spec A_U\): \(\nu\) is a universal
homeomorphism, hence integral, surjective and universally injective by
\Cref{thm:appA-universal-homeomorphism-criterion}, and universal
injectivity gives at most one point of \(Y\) over each point of \(X\).
Restricting to \(U\) preserves all of this.

Triviality of the residue extensions: universal injectivity says the residue
extensions \(\kappa(x)\subset\kappa(y)\) are purely inseparable.  In
characteristic zero a purely inseparable algebraic field extension is
trivial: if \(\alpha\in\kappa(y)\) then its minimal polynomial over
\(\kappa(x)\) has a single root, and in characteristic zero a minimal
polynomial is separable, so its degree is one.  Hence
\(\kappa(x)=\kappa(y)\).

The two conditions together are \Cref{def:bj-subintegral-extension}.
This is the point recorded in \Cref{rem:appA-subintegral-stronger}: in
positive characteristic the residue extensions may be nontrivial, and
Frobenius supplies a different mechanism; in characteristic zero
radiciality upgrades to subintegrality for free, and this upgrade is the
only thing the characteristic hypothesis buys at the level of the morphism.
\end{proof}

\begin{example}[The cuspidal normalization in characteristic zero]
\label{ex:ch44-cusp-basic}
On the affine cuspidal chart take
\[
              A=\mathbf C[t^2,t^3]\subset B=\mathbf C[t].
\]
The extension is finite and subintegral, but \(t\notin A\), even though
\(t^2,t^3\in A\).  If \(v\) is a semivaluation of the common function field
centered at the cusp and its lift, then
\[
       v(t)=\frac12v(t^2)=\frac13v(t^3),
\]
or, after changing signs,
\[
       \log|t|=\frac12\log|t^2|
              =\frac13\log|t^3|.
\]
Thus the coefficient \(t\) itself does not descend, while its valuative
function is recovered from downstairs powers.  This monomial is unusually
simple: general sums do not enjoy a Frobenius identity in characteristic
zero, which is why a uniform additive-power replacement is needed below.
\end{example}

\begin{remark}[The \(A\)-module structure of the cusp, used repeatedly]
\label{rem:ch44-cusp-decomposition}
For \(A=k[t^2,t^3]\subset B=k[t]\) with \(k\) any field, \(A\) is the set of
polynomials whose coefficient in degree one vanishes, so
\[
 B=A\oplus k\cdot t
\]
as an \(A\)-module, and the projection \(B\to k\) onto the second summand is
\(b\mapsto b'(0)\), the coefficient of \(t\).  Every computation with the
cusp below is an application of this one decomposition; in particular
\(\mathfrak c=t^2B\) is the conductor, and \(B/A\cong k\) is the coherent
quotient \(\mathcal Q\) of
\eqref{eq:bj-char-zero-quotient} in this example.  We record the analogous
description for \(A=k[t^2,t^5]\subset B=k[t]\), which will be used in
\Cref{ex:ch44-two-gaps}: there \(A\) is the set of polynomials whose
coefficients in degrees \(1\) and \(3\) vanish, so
\(B=A\oplus k\cdot t\oplus k\cdot t^3\) and \(\dim_kB/A=2\).
\end{remark}

\subsection{Seminormality, and why it is the complementary condition}
\index{seminormal ring!versus subintegral}
\label{sec:ch44-seminormality}

The reader may reasonably ask why the hypotheses of this chapter are not
vacuous: if \(A\subset B\) is subintegral and \(A\) is seminormal in \(B\),
then \(A=B\).  The resolution is that our \(X\) is precisely \emph{not}
required to be seminormal, and in the situation of interest --- \(X\)
unibranch and nonnormal --- it cannot be.

\begin{proposition}
\label{prop:ch44-seminormal-unibranch-normal}
Let \(k\) be algebraically closed of characteristic zero and let \(X\) be an
integral projective \(k\)-variety with normalization \(\nu:X^\nu\to X\).
\begin{enumerate}
\item If \(A\subset B\) is a subintegral extension and \(A\) is seminormal
      in \(B\), then \(A=B\).
\item If \(X\) is unibranch and seminormal, then \(X\) is normal.
\item Consequently, in \eqref{eq:bj-char-zero-universal-map} with
      \(\nu\) not an isomorphism, \(X\) is unibranch and \emph{not}
      seminormal, and the coherent quotient
      \(\mathcal Q=\nu_*\cO_Y/\cO_X\) is nonzero.
\end{enumerate}
\end{proposition}

\begin{proof}
(1)  Let \(b\in B\).  By \Cref{thm:appA-square-cube} the singleton
\(\{b\}\) lies in a tower
\(A=C_0\subset C_1\subset\cdots\subset C_r\subset B\) with
\(C_l=C_{l-1}[x_l]\) and \(x_l^2,x_l^3\in C_{l-1}\).  We show
\(C_l\subseteq A\) by induction on \(l\).  For \(l=0\) this is trivial.  If
\(C_{l-1}\subseteq A\), then \(x_l^2,x_l^3\in A\), and \(x_l\in B\), so
seminormality of \(A\) in \(B\) gives \(x_l\in A\), whence
\(C_l=C_{l-1}[x_l]\subseteq A\).  Thus \(b\in C_r\subseteq A\).  As \(b\)
was arbitrary, \(A=B\).

(2)  Let \(X\) be unibranch and seminormal, and let \(\nu:X^\nu\to X\) be
the normalization, which is finite and birational.  By
\Cref{lem:bj-unibranch-normalization-universal} it is a universal
homeomorphism, so by \Cref{lem:ch44-char-zero-subintegral} the extension
\(\cO_X(U)\subset\cO_{X^\nu}(\nu^{-1}U)\) is subintegral for every affine
\(U\).  Seminormality of \(X\) means that \(\cO_X(U)\) is seminormal in
\(\cO_{X^\nu}(\nu^{-1}U)\) for every affine \(U\); by~(1) these two rings
coincide, so \(\nu\) is an isomorphism and \(X\) is normal.

(3)  If \(\nu\) is not an isomorphism then \(X\) is not normal, so by~(2)
it is not seminormal.  Nonvanishing of \(\mathcal Q\) is the same statement:
\(\mathcal Q=0\) on an affine \(U\) says \(A_U=B_U\), and if that held for
all \(U\) then \(\nu\) would be an isomorphism.
\end{proof}

\begin{remark}[Why seminormality does not itself descend coefficients]
\label{rem:ch44-nonseminormal-shape}
If \(A\) were seminormal in \(B\), subintegrality would force \(A=B\), by
\Cref{prop:ch44-seminormal-unibranch-normal}\,(1).
Our target \(X\) is allowed to be nonseminormal, so individual coefficients
of an upstairs section need not lie in \(A\).
The proof descends carefully chosen additive-power combinations instead.

Said in another way,
\Cref{prop:ch44-seminormal-unibranch-normal}\,(3) is the precise reason why
no argument of this chapter can descend an individual coefficient.  If
\(b\in B_U\setminus A_U\) --- and by~(3) such \(b\) exist --- then \(b\) is
simply not a function on \(X\), and no amount of manipulation will make it
one.  What the additive-power theorem provides is not a way to move \(b\)
downstairs, but a way to move a \emph{corrected power} of \(b\) downstairs,
at the cost of losing the exact value of the associated Fubini--Study
function and retaining only a two-sided estimate for it.  The chapter is the
bookkeeping needed to show that the estimate is good enough.
\end{remark}

\begin{example}[A seminormal, nonnormal variety, where the descent behaves differently]
\label{ex:ch44-node-seminormal}
Let \(k\) be algebraically closed of characteristic zero and let
\(X\subset\mathbf P^2\) be the nodal cubic \(y^2z=x^3+x^2z\), with affine
chart
\[
 A=k[x,y]/(y^2-x^3-x^2)\cong k[t^2-1,\,t^3-t]\subset B=k[t],
\]
the isomorphism being \(t\mapsto y/x\).  Here \(\Spec B\to\Spec A\) has two
points over the node, namely \(t=\pm1\), so \(X\) is \emph{not} unibranch
and the extension is \emph{not} subintegral.  It is, however, the standard
example of a seminormal nonnormal ring: if \(b\in B\) satisfies
\(b^2,b^3\in A\) then \(b\in A\).  Indeed \(A=\{f\in B:f(1)=f(-1)\}\), and
if \(f(1)\ne f(-1)\) then \(f^2\) and \(f^3\) cannot both take equal values
at \(\pm1\), since \(f(1)^2=f(-1)^2\) and \(f(1)^3=f(-1)^3\) force
\(f(1)=f(-1)\) after dividing.

Three things behave differently, and it is worth listing them because each
is a hypothesis of the present chapter shown to be doing work.
\begin{enumerate}
\item \(\nu^{\An}\) is not injective.  The two points \(t=\pm1\) of
      \(Y=\mathbf P^1\) map to the node, so already the topological input
      \Cref{lem:bj-universal-analytic-homeomorphism} fails, and the phrase
      ``the same function of the same points'' has no meaning.
\item There is no additive-power identity to be had, and none is needed:
      \(A\) is seminormal in \(B\), so by
      \Cref{prop:ch44-seminormal-unibranch-normal}\,(1) the extension is
      subintegral only if \(A=B\), which it is not.  The Roberts--Singh
      theorem simply does not apply, its hypothesis being subintegrality
      and not integrality.
\item The conclusion itself fails.  For \(X\) nodal and \(\theta\) ample,
      \(\PSH_X(\theta)\) is not closed under regularized suprema; this is
      \Cref{ex:bj-node-cusp}, and the necessity of unibranchness in the
      ample case is \Cref{thm:ep-implies-unibranch}.  So the failure at
      item~(1) is not a defect of the method but a reflection of a genuine
      obstruction.
\end{enumerate}
For contrast, replace the node by the cusp \(y^2z=x^3\), i.e.\
\(A=k[t^2,t^3]\).  Now \(A\) is unibranch, \emph{not} seminormal
(\Cref{ex:appA-cusp-seminormalization}), the extension \emph{is}
subintegral, and everything in this chapter applies; the descent is carried
out explicitly in \Cref{ex:ch44-cusp-end-to-end}.  The two curves are
projectively indistinguishable at the level of degree, arithmetic genus, and
even of the homeomorphism type of the underlying topological space of
\(X\); they are distinguished exactly by whether the normalization is
injective.
\end{example}

\section{The Roberts--Singh additive-power theorem}
\index{additive power!Roberts--Singh theorem}\index{subintegral extension!additive powers}
\label{sec:ch44-roberts-singh}

\subsection{Statement and the meaning of exponent zero}
\label{sec:ch44-rs-statement}

Let \(A\subset B\) be a subintegral extension of \(\mathbf Q\)-algebras.
Roberts--Singh Theorem~4.17 says that every element \(b\in B\) is
subintegral over \(A\) in their elementwise sense.
Proposition~4.2 and Corollary~4.3 of \cite{RobertsSingh} give an
exponent-zero system:
there are \(p,N\ge0\) and fixed coefficients
\(c_1,\ldots,c_p\in B\) such that
\[
 a_m:=b^m+\sum_{i=1}^p\binom mi c_i b^{m-i}\in A
                         \qquad(m\ge N).
\tag{44.3}\label{eq:bj-rs-additive-powers}
\]
The coefficients \(c_i\) do not depend on \(m\).
That independence is the feature needed by finite differences.
The hypothesis that \(A\) contain \(\mathbf Q\) is also essential:
multinomial integers will have to be units.

The book's import of this statement is \Cref{thm:appA-roberts-singh}, with
its generic form \Cref{cor:appA-RS-generic}; the reader should treat
\eqref{eq:bj-rs-additive-powers} as a quotation of that theorem, and the
whole of this chapter as a study of what can be built from it.  Three
features of the statement carry the argument, and each is used at a
specific place.

\begin{enumerate}
\item \emph{The length \(p\) and the threshold \(N\) depend on \(b\), but
      not on \(m\).}  In the application \(b\) is the generic combination
      attached to a \emph{fixed} Fubini--Study presentation, so \(p\) and
      \(N\) are fixed once the presentation is fixed, and the exponent
      \(m\) --- which is what tends to infinity --- ranges freely
      afterwards.  This is the source of the constant \(C=2P\) in
      \eqref{eq:bj-char-zero-global-constants}, and of its independence
      of \(n\).
\item \emph{The coefficients \(c_i\) lie in \(B\), not in \(A\).}  They are
      upstairs data.  Globalizing them is the first of the three twisting
      problems of point~\textbf{(C2)}, solved by
      \Cref{lem:ch44-uniform-cutoffs}\,(1).
\item \emph{The binomial coefficients are those of a shifted power.}  This
      is what makes the finite-difference operator
      \(\Delta_b\alpha_m=\alpha_{m+1}-b\,\alpha_m\) act on the family
      \((a_m)_m\) by shifting the index of \(c_i\), and is the mechanism by
      which the individual \(c_ib^{m+p-i}\) are isolated from the \(a_m\).
      It is proved as \Cref{lem:ch44-finite-differences}.
\end{enumerate}

\begin{warning}{The exponent \(p\) is not a uniform datum of the extension}
\label{warn:ch44-p-not-uniform}
\eqref{eq:bj-rs-additive-powers} is quantified as: for every \(b\) there
exist \(p\), \(N\), \(c_\bullet\).  It is \emph{not} asserted, and is in
general false, that one \(p\) works for all \(b\in B\) at once; nor is any
such uniformity needed.  What is needed, and what
\Cref{lem:bj-local-additive-sandwich} extracts, is one \(p\) valid for the
single generic element \(b=\sum_jz_jg_j\) attached to a given finite family
\(g_1,\ldots,g_s\) --- from which a statement uniform in \(n\) about the
\emph{powers} of the ideal \((g_1,\ldots,g_s)\) follows.  The distinction is
exactly parallel to the elementwise-versus-uniform distinction of
\Cref{sec:ch43-uniformity}, and it is resolved here by genericity rather
than, as there, by a subalgebra argument.
\end{warning}

\subsection{The data, computed}
\label{sec:ch44-rs-computed}

Quoted theorems are more useful when one has seen their output.  Both
computations below are complete, and the second shows that the length \(p\)
genuinely exceeds one.

\begin{example}[The cusp: \(p=1\), with \(c_1\) exhibited]
\label{ex:ch44-cusp-rs}
Let \(R\) be any commutative \(\mathbf Q\)-algebra, and let
\[
 A:=\{f\in R[t]:[t^1]f=0\},
 \qquad
 B:=R[t],
\]
where \([t^k]f\) denotes the coefficient of \(t^k\).  For \(R=k\) this is
the cusp \(k[t^2,t^3]\subset k[t]\) of \Cref{rem:ch44-cusp-decomposition};
we allow a general \(R\) because the application takes
\(R=k[u,\mathbf z]\).  The extension is subintegral: \(t^2,t^3\in A\), and
\(B=A[t]\), so the one-step tower of \Cref{thm:appA-square-cube} suffices.

Let \(b=\sum_{k\ge0}\beta_kt^k\in B\).  We claim that
\eqref{eq:bj-rs-additive-powers} holds with
\[
 p=1,
 \qquad
 N=0,
 \qquad
 c_1=-\beta_1t=-\bigl([t^1]b\bigr)\,t.
\]
Indeed, write \(b=\beta_0+\beta_1t+t^2h\) with \(h\in B\).  Expanding,
\[
 [t^1]\bigl(b^m\bigr)
 =m\,\beta_1\beta_0^{\,m-1},
\]
because a monomial in the expansion of \(b^m\) contributes to degree one
only if it uses the factor \(\beta_1t\) exactly once and the factor
\(\beta_0\) the remaining \(m-1\) times; the term \(t^2h\) contributes only
in degrees \(\ge2\).  On the other hand
\[
 [t^1]\bigl(m\,\beta_1t\,b^{m-1}\bigr)
 =m\,\beta_1\,[t^0]\bigl(b^{m-1}\bigr)
 =m\,\beta_1\beta_0^{\,m-1}.
\]
Subtracting,
\[
 a_m=b^m-m\,\beta_1t\,b^{m-1}
 \quad\text{has}\quad
 [t^1]a_m=0,
 \quad\text{i.e.}\quad
 a_m\in A
 \qquad(m\ge0),
\]
with the convention that the second term is absent for \(m=0\).  So one
correction term suffices, and it is \emph{explicit}: the correction
coefficient is minus the degree-one part of \(b\) itself.

Two sanity checks.  For \(b=1+t\) one gets \(c_1=-t\) and
\(a_m=(1+t)^m-mt(1+t)^{m-1}=(1+t)^{m-1}\bigl(1-(m-1)t\bigr)\), whose
degree-one coefficient is \((m-1)-(m-1)=0\); this recovers the computation
quoted for contrast in \Cref{rem:ch43-charzero-contrast}.  For
\(b=t\) one gets \(\beta_0=0\), \(\beta_1=1\), \(c_1=-t\) and
\(a_m=t^m-mt\cdot t^{m-1}=(1-m)t^m\), which lies in \(A\) for \(m\ne1\) and
vanishes for \(m=1\); consistent, since \(t^m\in A\) for \(m\ne1\) anyway.
\end{example}

\begin{example}[Two gaps: \(p=3\), and the correction terms are not monomials in \(b\)]
\label{ex:ch44-two-gaps}
Let \(A=\mathbf Q[t^2,t^5]\subset B=\mathbf Q[t]\), so that
\(A=\{f:[t^1]f=[t^3]f=0\}\) and \(B=A\oplus\mathbf Qt\oplus\mathbf Qt^3\)
(\Cref{rem:ch44-cusp-decomposition}).  Take \(b=1+t\).  We determine the
smallest admissible length \(p\) and the coefficients.

Seek \(c_i=\alpha_it+\gamma_it^3\).  The two conditions are
\([t^1]a_m=0\) and \([t^3]a_m=0\) for all \(m\ge N\).

Degree one.  Since \([t^1]\bigl(t^3f\bigr)=0\) and
\([t^1]\bigl(tf\bigr)=[t^0]f=1\) for \(f=b^{m-i}\), we get
\[
 [t^1]a_m
 =\binom m1+\sum_{i=1}^p\binom mi\alpha_i
 =m(1+\alpha_1)+\sum_{i=2}^p\binom mi\alpha_i .
\]
The functions \(m\mapsto\binom mi\) are linearly independent polynomials in
\(m\) of degree \(i\), so the vanishing of this expression for all large
\(m\) forces
\[
 \alpha_1=-1,
 \qquad
 \alpha_i=0\quad(i\ge2).
\]
Degree three.  With \(c_1=-t+\gamma_1t^3\) and \(c_i=\gamma_it^3\) for
\(i\ge2\), and using \([t^3]b^r=\binom r3\), \([t^3](tb^r)=\binom r2\),
\([t^3](t^3b^r)=1\),
\[
 [t^3]a_m
 =\binom m3-\binom m1\binom{m-1}2+\sum_{i=1}^p\binom mi\gamma_i .
\]
Now
\[
 \begin{aligned}
 \binom m3-m\binom{m-1}2
 &=\frac{m(m-1)(m-2)}6
   -\frac{m(m-1)(m-2)}2\\
 &=-\frac{m(m-1)(m-2)}3
  =-2\binom m3 .
 \end{aligned}
\]
So the condition is \(\sum_{i=1}^p\binom mi\gamma_i=2\binom m3\) for all
large \(m\), which by the same linear independence forces
\[
 \gamma_3=2,
 \qquad
 \gamma_i=0\quad(i\ne3),
 \qquad
 p\ge3 .
\]
Hence the minimal data are
\[
 p=3,
 \qquad
 c_1=-t,
 \qquad
 c_2=0,
 \qquad
 c_3=2t^3,
 \qquad
 N=0 .
\]
Let us verify \(m=3\) and \(m=4\) by hand.  For \(m=3\),
\[
 (1+t)^3-3t(1+t)^2+2t^3
 =(1+3t+3t^2+t^3)-(3t+6t^2+3t^3)+2t^3
 =1-3t^2,
\]
which lies in \(A\).  For \(m=4\),
\[
 \begin{aligned}
 &(1+t)^4-4t(1+t)^3+4\cdot2t^3(1+t)\\
 &\quad=(1+4t+6t^2+4t^3+t^4)
 -(4t+12t^2+12t^3+4t^4)+(8t^3+8t^4)\\
 &\quad=1-6t^2+5t^4,
 \end{aligned}
\]
again in \(A\).  Two features are worth naming.  First, \(p=3\) exceeds the
number \(\dim_{\mathbf Q}B/A=2\) of gaps, so \(p\) is not simply the
colength of the conductor; it is governed by the degrees in \(m\) of the
functionals to be cancelled.  Second, \(c_2=0\) while \(c_3\ne0\): the
correction terms need not be supported in an initial segment of indices, so
the sum in \eqref{eq:bj-rs-additive-powers} cannot be truncated below \(p\).
\end{example}

\begin{remark}[What the examples do and do not show]
\label{rem:ch44-examples-scope}
Both examples exhibit the data of \eqref{eq:bj-rs-additive-powers} for a
particular \(b\) in a particular extension, by direct computation; neither
proves the theorem, whose content is the existence of such data for
\emph{every} \(b\) in \emph{every} subintegral extension of
\(\mathbf Q\)-algebras, and which we quote from \cite{RobertsSingh} through
\Cref{thm:appA-roberts-singh}.  What they do show is that the data are
computable, that \(p\) is a genuine parameter, and --- in
\Cref{ex:ch44-cusp-rs}, which is stated over an arbitrary
\(\mathbf Q\)-algebra \(R\) --- that in the case relevant to the worked
descent of \Cref{sec:ch44-computations} the generic version required by
\Cref{cor:appA-RS-generic} may be written down in closed form.  That closed
form is used in \Cref{ex:ch44-splitting-fails} and in
\Cref{ex:ch44-cusp-end-to-end}.
\end{remark}

\section{The local additive-power sandwich}
\label{sec:ch44-sandwich}

The GF weights require one formal variable, and coefficient extraction
requires another family of formal variables.
Keeping those two gradings separate prevents a false homogeneous splitting.

\subsection{The two gradings}
\label{sec:ch44-two-gradings}

Fix a finite subintegral extension \(A\subset B\) of
\(\mathbf Q\)-algebras.  We work in the polynomial ring
\(B[u,\mathbf z]=B[u,z_1,\ldots,z_s]\), equipped with two independent
gradings:
\begin{itemize}
\item the \(u\)-grading, in which \(\deg u=1\) and \(\deg z_j=\deg b=0\)
      for \(b\in B\).  It records Fubini--Study weights: a \(u^d\)-component
      will represent the weight \(-d\).  The subring \(A[u]\subset B[u]\)
      is \(u\)-graded, and so is every ideal generated by
      \(u\)-homogeneous elements.
\item the \(\mathbf z\)-grading, in which \(\deg z_j=1\) and
      \(\deg u=\deg b=0\).  It records the power of the ideal: the
      \(\mathbf z\)-degree-\(r\) part of an expression in
      \(b=\sum_jz_jg_j\) collects the monomials of degree \(r\) in the
      \(g_j\).
\end{itemize}
Both gradings are preserved by \(A[u,\mathbf z]\subset B[u,\mathbf z]\), and
both are used; the entire content of point~\textbf{(C1)} of
\Cref{sec:ch44-care} is the \emph{order} in which their homogeneous
components are extracted.  The order is:
\[
 \boxed{\ \text{form }a_m\ \longrightarrow\
 \text{take }\mathbf z\text{-components}\ \longrightarrow\
 \text{take }u\text{-components}.\ }
\]
The first arrow is legitimate because \(A[u,\mathbf z]\) is
\(\mathbf z\)-graded; the second because \(A[u]\) is \(u\)-graded.  Neither
arrow may be applied to \(b\) or to the \(c_i\) before \(a_m\) has been
formed, because \(a_m\in A[u,\mathbf z]\) is a statement about the total
expression and not about its pieces.

\begin{warning}{Do not split the correction coefficients}
\label{warn:ch44-no-splitting}
Let \(c_i=\sum_ec_i^{(e)}u^e\) be the \(u\)-decomposition of a
Roberts--Singh coefficient, and fix one index \(e\).  Although
\[
 b^m+\sum_{i=1}^p\binom mi\,c_i\,b^{m-i}\in A[u,\mathbf z]
 \qquad(m\ge N)
\]
by \eqref{eq:bj-rs-additive-powers}, it is \emph{false} in general that the
truncated expression
\[
 b^m+\sum_{i=1}^p\binom mi\,c_i^{(e)}u^e\,b^{m-i}
\]
lies in \(A[u,\mathbf z]\).  The reason is that \(b\) is \(u\)-inhomogeneous
whenever the weights \(d_j\) are not all equal, so the \(u\)-components of
\(a_m\) are convolutions, not termwise products.  An explicit failure, with
\(m=1\) and \(m=2\), is \Cref{ex:ch44-splitting-fails}\,(1).  The cost of
this error is not a worse constant but a fatal one: the elements one would
declare to be members of \(K_n\) do not lie in \(A[u]\) at all, so they are
not sections of any sheaf on \(X\), and there is nothing to descend.
\end{warning}

\subsection{The finite-difference calculus}
\label{sec:ch44-finite-differences}

The mechanism converting \eqref{eq:bj-rs-additive-powers} into an ideal
statement is elementary and purely formal.  We isolate it, because it is
used twice and because its hypotheses are weaker than those of the
surrounding lemma: no subintegrality, no \(\mathbf Q\)-algebra hypothesis,
and no finiteness are needed.

\begin{lemma}
\label{lem:ch44-finite-differences}
Let \(S\) be a commutative ring, \(b,c_1,\ldots,c_p\in S\), and let
\((\alpha_m)_{m\ge N}\) be the family
\[
 \alpha_m:=b^m+\sum_{i=1}^p\binom mic_ib^{m-i}
 \qquad(m\ge N).
\]
Define \(\Delta_b\alpha_m:=\alpha_{m+1}-b\,\alpha_m\) and let
\(\Delta_b^r\) denote its \(r\)-fold iterate.  Then, for \(0\le r\le p\) and
\(m\ge N\),
\[
 \Delta_b^r\alpha_m
 =\sum_{i=r}^p\binom m{i-r}c_ib^{m+r-i},
\]
and consequently
\[
 c_rb^{m+p-r}\in(\alpha_m,\ldots,\alpha_{m+p})S
 \qquad(1\le r\le p),
 \qquad
 b^{m+p}\in(\alpha_m,\ldots,\alpha_{m+p})S .
\]
\end{lemma}

\begin{proof}
The displayed formula for \(\Delta_b^r\) is proved by induction on \(r\).
For \(r=0\) it is the definition of \(\alpha_m\), the term \(i=0\) being
read as \(b^m\) with \(\binom m0=1\); we therefore adopt the convention
\(c_0:=1\) and write \(\alpha_m=\sum_{i=0}^p\binom mic_ib^{m-i}\).

Assume the formula for \(r\), that is
\(\Delta_b^r\alpha_m=\sum_{i=r}^p\binom m{i-r}c_ib^{m+r-i}\).  Then
\[
 \Delta_b^{r+1}\alpha_m
 =\Delta_b^r\alpha_{m+1}-b\,\Delta_b^r\alpha_m
 =\sum_{i=r}^p\binom{m+1}{i-r}c_ib^{m+1+r-i}
  -\sum_{i=r}^p\binom m{i-r}c_ib^{m+1+r-i},
\]
where we used that \(\Delta_b^r\) applied to the family shifts the index and
that the same expression is valid at \(m+1\).  By Pascal's identity
\(\binom{m+1}{j}-\binom mj=\binom m{j-1}\), the difference of the two sums
is
\[
 \sum_{i=r}^p\binom m{i-r-1}c_ib^{m+1+r-i}
 =\sum_{i=r+1}^p\binom m{i-(r+1)}c_ib^{m+(r+1)-i},
\]
the term \(i=r\) dropping out because \(\binom m{-1}=0\).  This is the
formula for \(r+1\).

Each \(\Delta_b^r\alpha_m\) is by construction an \(S\)-linear combination
of \(\alpha_m,\ldots,\alpha_{m+r}\), with coefficients in
\(\mathbf Z[b]\subseteq S\); explicitly
\(\Delta_b^r\alpha_m=\sum_{l=0}^r(-1)^{r-l}\binom rl b^{r-l}\alpha_{m+l}\).
Hence \(\Delta_b^r\alpha_m\in(\alpha_m,\ldots,\alpha_{m+p})S\) for
\(r\le p\).

We now prove \(c_rb^{m+p-r}\in(\alpha_m,\ldots,\alpha_{m+p})S\) by downward
induction on \(r\), starting at \(r=p\).  For \(r=p\) the formula reads
\(\Delta_b^p\alpha_m=\binom m0c_pb^m=c_pb^m\), which is the assertion.
Suppose the assertion known for all indices \(>r\).  Multiply the formula
for \(\Delta_b^r\alpha_m\) by \(b^{p-r}\):
\[
 b^{p-r}\Delta_b^r\alpha_m
 =c_rb^{m+p-r}+\sum_{i>r}\binom m{i-r}c_ib^{m+p-i}.
\]
The left-hand side lies in \((\alpha_m,\ldots,\alpha_{m+p})S\), and by the
inductive hypothesis so does every summand on the right with \(i>r\).
Hence so does \(c_rb^{m+p-r}\).

Finally, multiplying \(\alpha_m\) itself by \(b^p\) gives
\[
 b^p\alpha_m=b^{m+p}+\sum_{i=1}^p\binom mic_ib^{m+p-i},
\]
whose left-hand side and whose sum both lie in
\((\alpha_m,\ldots,\alpha_{m+p})S\) by what precedes; hence so does
\(b^{m+p}\).
\end{proof}

\begin{remark}[Why exponent zero is the operative hypothesis]
\label{rem:ch44-exponent-zero}
\Cref{lem:ch44-finite-differences} uses the \(c_i\) only through the fact
that they do not depend on \(m\).  Had \eqref{eq:bj-rs-additive-powers} been
available only in the weaker form with \(m\)-dependent coefficients
\(c_i^{(m)}\), the operator \(\Delta_b\) would not act on the family by
shifting binomial indices, the telescoping above would fail at the first
step, and no ideal-theoretic conclusion would be available.  This is the
precise sense in which the \emph{exponent-zero} refinement
\cite[Proposition~4.2 and Corollary~4.3]{RobertsSingh} --- and not merely
the elementwise subintegrality of \cite[Theorem~4.17]{RobertsSingh} --- is
what the chapter consumes.  It is recorded in \Cref{thm:appA-roberts-singh}
in exactly that form.
\end{remark}

\subsection{The sandwich}
\label{sec:ch44-sandwich-proof}

\begin{lemma}\label{lem:bj-local-additive-sandwich}
Let \(A\subset B\) be a finite subintegral extension of
\(\mathbf Q\)-algebras.
Let \(g_1,\ldots,g_s\in B[u]\) be homogeneous elements for the
\(u\)-grading, and set
\[
                         \mathfrak a=(g_1,\ldots,g_s)\subset B[u].
\]
There are integers \(p,N\ge0\) such that for every
\(n\ge\max\{0,N-p\}\) one can construct a homogeneous ideal
\(K_n\subset A[u]\) satisfying
\[
             \mathfrak a^{\,n+2p}
               \subset K_nB[u]
               \subset\mathfrak a^n.
\tag{44.4}\label{eq:bj-local-additive-sandwich}
\]
The generators are covariant under a simultaneous rescaling
\(g_j\mapsto t^{-1}g_j\) with \(t\in A^\times\).
\end{lemma}

\begin{proof}
Introduce independent variables \(\mathbf z=(z_1,\ldots,z_s)\) and form the
generic element
\[
                         b=\sum_j z_jg_j
                         \in B[u,\mathbf z].
\tag{44.5}\label{eq:bj-char-zero-generic-element}
\]
By Tags 0CND and 0EUH, the polynomial extension
\[
              A[u,\mathbf z]\subset B[u,\mathbf z]
\]
is subintegral.
Apply \eqref{eq:bj-rs-additive-powers} to \(b\).
There are fixed \(c_1,\ldots,c_p\in B[u,\mathbf z]\) and \(N\) such that
\[
 a_m=b^m+\sum_{i=1}^p\binom mi c_i b^{m-i}
                    \in A[u,\mathbf z]
                    \qquad(m\ge N).
\tag{44.6}\label{eq:bj-char-zero-am}
\]

\emph{Step 1: normalizing the \(\mathbf z\)-degree.}
Give every \(z_j\) degree one.
Take the \(\mathbf z\)-degree-\(m\) component of
\eqref{eq:bj-char-zero-am}.
This replaces \(c_i\) by its \(\mathbf z\)-degree-\(i\) component and makes
\(a_m\) homogeneous of \(\mathbf z\)-degree \(m\).
We keep the notation \(c_i,a_m\).

Two verifications are needed here and are worth spelling out.  First, the
operation is legitimate: \(A[u,\mathbf z]\) is a \(\mathbf z\)-graded
subring of \(B[u,\mathbf z]\), so the \(\mathbf z\)-homogeneous components
of an element of \(A[u,\mathbf z]\) again lie in \(A[u,\mathbf z]\).
Second, the replacement is consistent across all \(m\) simultaneously:
since \(b\) is \(\mathbf z\)-homogeneous of degree one, the summand
\(\binom mi\tilde c_ib^{m-i}\), where \(\tilde c_i\) denotes the original
Roberts--Singh coefficient, has \(\mathbf z\)-degree-\(m\) part equal to
\(\binom mic_ib^{m-i}\) with \(c_i\) the \(\mathbf z\)-degree-\(i\)
component of \(\tilde c_i\); the component extracted from \(\tilde c_i\) is
the \emph{same} for every \(m\), namely the one in degree \(i\).  So the new family \((a_m)\) is again of the form treated by
\Cref{lem:ch44-finite-differences}, with the same \(p\) and \(N\) and with
coefficients still independent of \(m\).  Had the extracted component
depended on \(m\), \Cref{rem:ch44-exponent-zero} would apply and the
argument would stop here.

\emph{Step 2: no premature splitting.}
At this point we do not split \(c_i\) into \(u\)-homogeneous components.
Although every \(g_j\) is \(u\)-homogeneous, the degrees may differ, so the
generic sum \(b\) is usually \(u\)-inhomogeneous.
Splitting before the additive-power identities would destroy the algebraic
relation between its terms.
We first form the complete \(a_m\), then extract \(\mathbf z\)-coefficients,
and only then split those coefficients by \(u\)-degree.

Explicitly, the order of operations is the boxed one of
\Cref{sec:ch44-two-gradings}, and the two prohibited shortcuts are those of
\Cref{warn:ch44-no-splitting} and of
\Cref{ex:ch44-splitting-fails}\,(2).  Neither is used below.

\emph{Step 3: isolating the leading power.}
Define the finite-difference operator
\[
                         \Delta_ba_m:=a_{m+1}-ba_m.
\]
Pascal's identity gives, by induction on \(r\),
\[
 \Delta_b^r a_m
       =\sum_{i=r}^p\binom m{i-r}c_i b^{m+r-i}.
\tag{44.7}\label{eq:bj-char-zero-finite-differences}
\]
Every left-hand side belongs to the ideal
\((a_m,\ldots,a_{m+p})B[u,\mathbf z]\).

For \(r=p\), \eqref{eq:bj-char-zero-finite-differences} gives \(c_pb^m\).
Suppose, descending in \(r\), that the terms with index greater than \(r\)
have been isolated.
Multiply the identity for \(r\) by \(b^{p-r}\).
After subtracting the already isolated higher-index terms, the only remaining
term is \(c_rb^{m+p-r}\).
Thus every such correction term lies in
\((a_m,\ldots,a_{m+p})B[u,\mathbf z]\).

Now
\[
 b^pa_m=b^{m+p}
       +\sum_{i=1}^p\binom mi c_i b^{m+p-i}.
\]
Subtracting the isolated corrections yields
\[
 b^{m+p}\in(a_m,\ldots,a_{m+p})B[u,\mathbf z]
                         \qquad(m\ge N).
\tag{44.8}\label{eq:bj-char-zero-leading-power}
\]
This whole step is \Cref{lem:ch44-finite-differences} applied in the ring
\(S=B[u,\mathbf z]\); we have repeated the computation in the notation at
hand rather than merely citing it, because the indices matter for what
follows.

\emph{Step 4: the upper inclusion.}
Fix \(n\ge\max\{0,N-p\}\) and take the block
\[
                         a_{n+p},\ldots,a_{n+2p}.
\]
Every index in the block is at least \(N\): indeed
\(n+p\ge\max\{0,N-p\}+p\ge N\).
Every \(\mathbf z\)-coefficient of every block belongs to
\(\mathfrak a^n\).
Indeed, the correction involving \(c_i\) retains at least \(n\) factors from
the generic element, and the leading term has more.

In detail: for \(n+p\le m\le n+2p\) the summand \(\binom mic_ib^{m-i}\)
has \(m-i\ge(n+p)-p=n\), and the \(\mathbf z\)-coefficients of \(b^{m-i}\)
are, up to multinomial integers, the monomials
\(g^\gamma=g_1^{\gamma_1}\cdots g_s^{\gamma_s}\) with
\(|\gamma|=m-i\ge n\); each such monomial lies in \(\mathfrak a^n\), and
\(c_i\in B[u,\mathbf z]\) multiplies it back into \(\mathfrak a^n\), which is an
ideal of \(B[u]\).  The leading term \(b^m\) contributes monomials of degree
\(m\ge n\).  Hence every \(\mathbf z\)-coefficient of every \(a_m\) in the
block lies in \(\mathfrak a^n\).

Let \(K_n\subset A[u]\) be generated by every \(u\)-homogeneous component of
all those \(\mathbf z\)-coefficients.
These components do lie in \(A[u]\): the \(\mathbf z\)-coefficients of
\(a_m\) lie in \(A[u]\) by \eqref{eq:bj-char-zero-am} and Step~1, and
\(A[u]\) is \(u\)-graded.
The ideal \(\mathfrak a^n\) is \(u\)-homogeneous, so component extraction
preserves the inclusion
\[
                         K_nB[u]\subset\mathfrak a^n.
\]
That last implication is the only place where \(u\)-homogeneity of the
\(g_j\) is used, and it is used essentially: if \(\mathfrak a\) were not a
\(u\)-homogeneous ideal, an element of \(\mathfrak a^n\) could perfectly
well have \(u\)-components outside \(\mathfrak a^n\), and \(K_n\) would fail
the upper inclusion.

\emph{Step 5: the lower inclusion, and where \(\mathbf Q\) is used.}
Use \eqref{eq:bj-char-zero-leading-power} with \(m=n+p\).
Every coefficient of \(b^{n+2p}\) belongs to \(K_nB[u]\).
Such a coefficient is a nonzero multinomial integer times a monomial of total
degree \(n+2p\) in the \(g_j\).
Because \(A\) is a \(\mathbf Q\)-algebra, that integer is a unit.
All generators of \(\mathfrak a^{n+2p}\) therefore belong to \(K_nB[u]\).
This proves \eqref{eq:bj-local-additive-sandwich}.

Once more in detail: \eqref{eq:bj-char-zero-leading-power} places
\(b^{n+2p}\) in the ideal generated by the block; the block members are
\(\mathbf z\)-homogeneous of degrees \(n+p,\ldots,n+2p\), and the ideal they
generate in \(B[u,\mathbf z]\) is \(\mathbf z\)-graded, so extracting the
\(\mathbf z\)-coefficient of a fixed monomial \(\mathbf z^\beta\) with
\(|\beta|=n+2p\) from both sides is legitimate and gives
\[
 \binom{n+2p}{\beta}\,g^\beta\in K_nB[u],
 \qquad
 \binom{n+2p}{\beta}=\frac{(n+2p)!}{\beta_1!\cdots\beta_s!},
\]
by the multinomial theorem applied to \(b^{n+2p}=(\sum_jz_jg_j)^{n+2p}\).
The multinomial coefficient is a positive integer, hence invertible in the
\(\mathbf Q\)-algebra \(A\), so \(g^\beta\in K_nB[u]\).  The monomials
\(g^\beta\) with \(|\beta|=n+2p\) generate \(\mathfrak a^{n+2p}\), whence
the lower inclusion.  This is the third and last use of the characteristic
hypothesis listed in \Cref{sec:ch44-hypotheses}, and it is the reason the
lemma cannot simply be reduced modulo \(p\) to give a second proof of the
positive-characteristic case.

\emph{Step 6: covariance under a change of frame.}
Finally change a local frame by \(e'=te\), \(t\in A^\times\).
Then \(g_j'=t^{-1}g_j\) and \(b'=t^{-1}b\).
Transport the Roberts--Singh choices by
\[
                         c_i'=t^{-i}c_i,
             \qquad
                         a_m'=t^{-m}a_m.
\tag{44.9}\label{eq:bj-char-zero-frame-law}
\]
These choices preserve \eqref{eq:bj-char-zero-am}.
Thus the extracted coefficients, after multiplication by the appropriate
power of the frame, are intrinsic local sections.

Concretely: substituting \(b'=t^{-1}b\), \(c_i'=t^{-i}c_i\) into the
left-hand side of \eqref{eq:bj-char-zero-am} gives
\[
 a_m'=(t^{-1}b)^m+\sum_{i=1}^p\binom mi t^{-i}c_i\,(t^{-1}b)^{m-i}
 =t^{-m}\Bigl(b^m+\sum_{i=1}^p\binom mic_ib^{m-i}\Bigr)
 =t^{-m}a_m,
\]
so \(a_m'\in A[u,\mathbf z]\) because \(t^{-1}\in A^\times\).  Hence the
element \(a_m\,e^{\otimes m}\), formed from a section-valued datum and a
frame \(e\), is independent of the frame; and the same holds for
\(c_i\,e^{\otimes i}\).  It is this observation, and not any global
vanishing theorem, that makes the local systems of the next section fit
together into sections of a line bundle.
\end{proof}

\begin{remark}[The two grading rule]
\label{rem:ch44-two-grading-rule}
The \(\mathbf z\)-grading counts powers of the ideal and isolates monomial
coefficients.
The \(u\)-grading stores the FS weights.
The \(\mathbf z\)-homogeneous component is taken before the \(u\)-components.
Interchanging this order is not justified because \(b\) is not
\(u\)-homogeneous.
\end{remark}

\begin{remark}[What the sandwich is, dimensionally]
\label{rem:ch44-sandwich-shape}
It is worth pausing on the shape of \eqref{eq:bj-local-additive-sandwich}.
The ideal \(K_n\) lives downstairs and is squeezed between two powers of an
upstairs ideal whose exponents differ by the \emph{constant} \(2p\).  Read
valuatively --- which is how it will be read in
\Cref{sec:ch44-global-sandwich} --- this says that
\(W(K_nB[u])\) agrees with \(n\,W(\mathfrak a)\) to within an additive error
\(2p\,W(\mathfrak a)\), and after division by \(n+2p\) the error is
\(O(1/n)\).  A statement with \(n+2p\) replaced by \(\kappa n\) would give
relative error \(1-\kappa^{-1}\), which does not tend to zero and is
useless.  Thus the entire force of the lemma is in the word ``\(2p\)'':
\emph{additive}, and independent of \(n\).  This is the sense in which
\Cref{lem:bj-local-additive-sandwich} is the exact
characteristic-zero substitute for the identity \(f^q\in A\) of
\Cref{cor:ch43-frobenius-all-affines}, which has additive error zero.
\end{remark}

\begin{example}[Splitting the weight components before applying the identity fails]
\label{ex:ch44-splitting-fails}
We exhibit both prohibited shortcuts of \Cref{sec:ch44-care}\,\textbf{(C1)}
on one example, with explicit numbers.

Let \(A=\mathbf Q[t^2,t^3]\subset B=\mathbf Q[t]\), and let
\(s=2\) with
\[
 g_1:=1+t
 \quad(u\text{-degree }0),
 \qquad
 g_2:=u\,(2+t)
 \quad(u\text{-degree }1).
\]
Both are \(u\)-homogeneous, so the hypotheses of
\Cref{lem:bj-local-additive-sandwich} hold, and
\(\mathfrak a=(g_1,g_2)\subset B[u]\).  The generic element is
\[
 b=z_1(1+t)+z_2u(2+t)
 =\underbrace{(z_1+2uz_2)}_{\beta_0}
  +\underbrace{(z_1+uz_2)}_{\beta_1}t,
\]
and by \Cref{ex:ch44-cusp-rs} applied over
\(R=\mathbf Q[u,z_1,z_2]\) the Roberts--Singh data may be taken to be
\[
 p=1,
 \qquad
 N=0,
 \qquad
 c_1=-\beta_1t=-(z_1+uz_2)\,t,
 \qquad
 a_m=b^m-m(z_1+uz_2)t\,b^{m-1}.
\]

\smallskip
\emph{(1) Splitting the correction coefficient by \(u\)-degree.}
The \(u\)-components of \(c_1\) are
\[
 c_1^{(0)}=-z_1t,
 \qquad
 c_1^{(1)}u=-z_2ut .
\]
Neither is an admissible correction term.  For \(m=1\),
\[
 b+c_1^{(0)}
 =(z_1+2uz_2)+(z_1+uz_2)t-z_1t
 =(z_1+2uz_2)+uz_2\,t,
\]
whose coefficient in degree one is \(uz_2\ne0\); so
\(b+c_1^{(0)}\notin A[u,\mathbf z]\).  For \(m=2\),
\[
 [t^1]\bigl(b^2-2z_1t\,b\bigr)
 =2\beta_0\beta_1-2z_1\beta_0
 =2\beta_0(\beta_1-z_1)
 =2uz_2(z_1+2uz_2)\ne0,
\]
again a failure.  By contrast the unsplit correction works at every \(m\):
\(a_1=b-\beta_1t=\beta_0\in A[u,\mathbf z]\), and
\(a_2=b^2-2\beta_1tb=(\beta_0+\beta_1t)(\beta_0-\beta_1t)
=\beta_0^2-\beta_1^2t^2\in A[u,\mathbf z]\) since \(t^2\in A\).
The failure is fatal rather than quantitative: an ideal generated by the
\(u\)-components of \(b+c_1^{(0)}\) is not contained in \(A[u]\), so its
members are not functions on \(X\) and there is nothing to descend.

\smallskip
\emph{(2) Applying the identity to the weight components separately.}
Now suppose one groups the generators by \(u\)-degree and applies
\eqref{eq:bj-rs-additive-powers} to each group.  Here the groups are
singletons, with generic elements
\[
 b^{(0)}=z_1g_1,
 \qquad
 b^{(1)}=z_2g_2 .
\]
By \Cref{ex:ch44-cusp-rs} again, \(p=1\) for each, and the resulting blocks
\(a^{(0)}_m,a^{(1)}_m\) for \(m\in\{n+1,n+2\}\) have
\(\mathbf z\)-coefficients
\[
 G_m:=(1+t)^m-mt(1+t)^{m-1}=(1+t)^{m-1}\bigl(1-(m-1)t\bigr),
\]
\[
 H_m:=u^m\Bigl((2+t)^m-mt(2+t)^{m-1}\Bigr)
      =u^m(2+t)^{m-1}\bigl(2-(m-1)t\bigr),
\]
respectively, each already \(u\)-homogeneous, of \(u\)-degrees \(0\) and
\(m\).  Let \(K_n^{\mathrm{grp}}\subset A[u]\) be the ideal they generate,
and consider the monomial
\[
 g_1^{2}g_2^{\,n}=(1+t)^2u^n(2+t)^n
 \in\mathfrak a^{\,n+2}=\mathfrak a^{\,n+2p}.
\]
We claim \(g_1^2g_2^n\notin K_n^{\mathrm{grp}}B[u]\) as soon as \(n\ge3\).
Indeed \(K_n^{\mathrm{grp}}B[u]\) is generated by \(u\)-homogeneous elements
of \(u\)-degrees \(0\) (namely \(G_{n+1},G_{n+2}\)) and \(n+1,n+2\) (namely
\(H_{n+1},H_{n+2}\)), hence is a \(u\)-homogeneous ideal; and
\(g_1^2g_2^n\) is \(u\)-homogeneous of degree \(n\).  Since \(B[u]\) has no
elements of negative \(u\)-degree, the generators \(H_{n+1},H_{n+2}\) cannot
contribute to \(u\)-degree \(n\).  So the claim reduces to
\[
 (1+t)^2(2+t)^n\in\bigl(G_{n+1},G_{n+2}\bigr)B .
\]
Now \(B=\mathbf Q[t]\) is a principal ideal domain, the polynomials
\(1-nt\), \(1-(n+1)t\) and \(1+t\) are pairwise coprime for \(n\ge1\), and
\[
 \gcd\bigl(G_{n+1},G_{n+2}\bigr)
 =\gcd\Bigl((1+t)^n(1-nt),\ (1+t)^{n+1}\bigl(1-(n+1)t\bigr)\Bigr)
 =(1+t)^n .
\]
So \((G_{n+1},G_{n+2})B=\bigl((1+t)^n\bigr)\), and
\((1+t)^2(2+t)^n\) lies in it if and only if \((1+t)^n\) divides
\((1+t)^2(2+t)^n\), i.e.\ if and only if \(n\le2\).  This proves the claim.

The correct construction, by contrast, does contain this monomial: taking
the \(\mathbf z^{(2,n)}\)-coefficient in
\eqref{eq:bj-char-zero-leading-power} with \(m=n+p=n+1\) gives
\[
 \binom{n+2}{2}\,g_1^2g_2^{\,n}\in K_nB[u],
\]
and \(\binom{n+2}2\) is a unit in \(\mathbf Q\).

\smallskip
\emph{(3) The size of the damage.}
The grouped construction is not merely different, it is quantitatively
worse.  What survives it is only the pigeonhole statement: any monomial of
degree \(s(M-1)+1\) in \(s\) generators has some \(g_j\) to a power at least
\(M\), so
\[
 \mathfrak a^{\,s(n+2p)}\subset K_n^{\mathrm{grp}}B[u]\subset\mathfrak a^n
\]
is available while \eqref{eq:bj-local-additive-sandwich} is not.  Feeding
this into \Cref{sec:ch44-decreasing} in place of
\eqref{eq:bj-char-zero-global-weight-sandwich} would give
\(r_n=n/(sn+2sp)\to1/s\), and the construction would converge to
\(s^{-1}\psi_0\) instead of \(\psi_0\).  The loss would be
\emph{multiplicative}, and the theorem would be false as stated; cf.\
\Cref{rem:ch44-sandwich-shape}.
\end{example}

\section{Encoding a weighted GF function}
\label{sec:ch44-encoding}

The local sandwich is a statement about ideals.  To use it we must translate
a Fubini--Study function into an ideal and back.  The translation is the
content of this section, and it is worth stating at the outset what it has
to achieve: a GF function carries two kinds of data --- sections, and
\emph{weights} --- and the ideal must remember both, because the weights are
what distinguishes \(\max_j\{\log|s_j|+\lambda_j\}\) from
\(\max_j\log|s_j|\).  The device is a formal variable \(u\) of Gauss weight
one; a factor \(u^d\) attached to a section records the weight \(-d\).

\subsection{The Gauss weight}
\index{Gauss weight!of a homogeneous ideal}
\label{sec:ch44-gauss}

Let \(v\) be a semivaluation with center \(c(v)\), and let \(B\) be a local
ring at that center (or any ring on which \(v\) is defined and nonnegative).
The definition \eqref{eq:bj-char-zero-gauss-weight} below is the standard
Gauss extension of \(v\) to \(B[u]\) in which \(u\) is given value one.  We
record its formal properties once, because they are used at every subsequent
step and because the second of them --- that \(W_v\) is a \emph{valuation},
not merely a filtration --- is what converts the ideal inclusions of
\Cref{lem:bj-local-additive-sandwich} into the numerical inequalities of
\Cref{prop:ch44-global-sandwich}.

\begin{lemma}
\label{lem:ch44-gauss-valuation}
Let \(v:B\to[0,+\infty]\) be a semivaluation and define, for
\(f=\sum_db_du^d\in B[u]\),
\[
 W_v(f):=\min_d\{v(b_d)+d\}.
\]
Then:
\begin{enumerate}
\item\label{itm:ch44-gauss:well}
      \(W_v\) is well defined, \(W_v(u)=1\), \(W_v(b)=v(b)\) for \(b\in B\),
      and \(W_v(f)=+\infty\) if and only if every \(b_d\) has
      \(v(b_d)=+\infty\);
\item\label{itm:ch44-gauss:val}
      \(W_v(fg)=W_v(f)+W_v(g)\) and
      \(W_v(f+g)\ge\min\{W_v(f),W_v(g)\}\); that is, \(W_v\) is a
      semivaluation on \(B[u]\);
\item\label{itm:ch44-gauss:ideal}
      if \(\mathfrak b\subseteq B[u]\) is an ideal and
      \(W_v(\mathfrak b):=\min_{f\in\mathfrak b}W_v(f)\), then
      \(W_v(\mathfrak b)=\min_lW_v(h_l)\) for any finite generating set
      \((h_l)\), and
      \(W_v(\mathfrak b\mathfrak b')=W_v(\mathfrak b)+W_v(\mathfrak b')\);
      in particular \(W_v(\mathfrak b^n)=n\,W_v(\mathfrak b)\);
\item\label{itm:ch44-gauss:mono}
      \(\mathfrak b\subseteq\mathfrak b'\) implies
      \(W_v(\mathfrak b)\ge W_v(\mathfrak b')\).
\end{enumerate}
\end{lemma}

\begin{proof}
Item~\cref{itm:ch44-gauss:well} is immediate from the definition, the minimum
being taken over the finitely many \(d\) with \(b_d\neq0\) and being
\(+\infty\) when all terms are.

Item~\cref{itm:ch44-gauss:val}: the ultrametric inequality for a sum is clear,
since the coefficient of \(u^d\) in \(f+g\) is a sum of the corresponding
coefficients.  For the product, write \(W_v(f)=v(b_{d_0})+d_0\) with \(d_0\)
minimal among the indices realizing the minimum, and similarly
\(W_v(g)=v(c_{e_0})+e_0\) with \(e_0\) minimal.  The coefficient of
\(u^{d_0+e_0}\) in \(fg\) is \(\sum_{d+e=d_0+e_0}b_dc_e\).  For
\((d,e)\ne(d_0,e_0)\) with \(d+e=d_0+e_0\) one has either \(d<d_0\) or
\(e<e_0\); in the first case \(v(b_d)+d>W_v(f)\) by minimality of \(d_0\),
and \(v(c_e)+e\ge W_v(g)\), so \(v(b_dc_e)+d+e>W_v(f)+W_v(g)\); the second
case is symmetric.  Hence the term \(b_{d_0}c_{e_0}\) strictly dominates and
\(v\bigl(\sum_{d+e=d_0+e_0}b_dc_e\bigr)=v(b_{d_0})+v(c_{e_0})\), giving
\(W_v(fg)\le W_v(f)+W_v(g)\).  The reverse inequality holds coefficientwise.
(If \(W_v(f)=+\infty\) or \(W_v(g)=+\infty\) both sides are \(+\infty\), by
item~\cref{itm:ch44-gauss:well} and multiplicativity of \(v\) on the support
prime.)

Item~\cref{itm:ch44-gauss:ideal}: any \(f\in\mathfrak b\) is
\(\sum_l\gamma_lh_l\) with \(\gamma_l\in B[u]\), so
\(W_v(f)\ge\min_l\{W_v(\gamma_l)+W_v(h_l)\}\ge\min_lW_v(h_l)\), because
\(W_v\ge0\) on \(B[u]\) when \(v\ge0\) on \(B\).  The reverse is clear.  For
products, \(\mathfrak b\mathfrak b'\) is generated by the \(h_lh'_{l'}\), so
\(W_v(\mathfrak b\mathfrak b')=\min_{l,l'}\{W_v(h_l)+W_v(h'_{l'})\}
=W_v(\mathfrak b)+W_v(\mathfrak b')\) by
item~\cref{itm:ch44-gauss:val}.

Item~\cref{itm:ch44-gauss:mono} is immediate, a minimum over a larger set
being no larger.
\end{proof}

\begin{remark}[The two arithmetics in \(\Rbar\)]
\label{rem:ch44-gauss-infinity}
\Cref{lem:ch44-gauss-valuation} is stated with values in
\([0,+\infty]\), and the value \(+\infty\) occurs exactly at points where
all the relevant sections vanish.  The additions performed above are of the
form \(a+b\) with \(a,b\in[0,+\infty]\), which is unambiguous; the
conventions of \Cref{app:normalizations} forbid only
\((+\infty)-(+\infty)\), and no such expression occurs anywhere in this
chapter.  The one place where the value \(+\infty\) requires an explicit
argument rather than an arithmetic convention is the common-zero clause of
\Cref{prop:ch44-global-sandwich}, and it is discharged there.
\end{remark}

\subsection{From a GF function to a weighted ideal}
\label{sec:ch44-encoding-ideal}

Let \(L\in\Pic(X)_{\mathbf Q}\) and let
\[
 \psi=\frac1m\max_j\{\log|s_j|+\lambda_j\}
          \in\cH^{\mathrm{gf}}_{\mathbf Q}(\nu^*L),
 \qquad
 s_j\in H^0(Y,\nu^*(mL)).
\tag{44.10}\label{eq:bj-char-zero-gf}
\]
Discard zero sections.
Put
\[
                         c=\frac1m\max_j\lambda_j
\]
and \(\psi_0=\psi-c\).
Choose \(q>0\) such that
\[
                 d_j:=-q(\lambda_j-\max_\ell\lambda_\ell)
                         \in\mathbf N.
\tag{44.11}\label{eq:bj-char-zero-weight-exponents}
\]
Set
\[
            M:=qmL,
            \qquad
            \rho:=\frac1{qm},
            \qquad
            \mathfrak a=(u^{d_j}s_j^q)_j.
\tag{44.12}\label{eq:bj-char-zero-weighted-ideal}
\]
Here \(\mathfrak a\) is an upstairs family of weighted sections of
\(\nu^*M\).

Three bookkeeping points.  First, \(q\) exists because the \(\lambda_j\) are
rational: take \(q\) a common denominator of the finitely many differences
\(\lambda_j-\max_\ell\lambda_\ell\).  Second, \(d_j\ge0\) by construction,
with \(d_{j_0}=0\) for any index \(j_0\) attaining
\(\max_\ell\lambda_\ell\); this single fact is what will force
\(W(\mathfrak a)(v_\triv)=0\) and hence exclude the identically
\(-\infty\) alternative in \Cref{sec:ch44-decreasing}.  Third, \(M=qmL\) is
an honest line bundle on \(X\), and \(\mathfrak a\) is a family of sections
of \(\nu^*M\) tensored with powers of \(u\); it is \emph{not} a family of
sections of a bundle on \(X\), and the whole of
\Cref{sec:ch44-globalization} exists to manufacture one.

For a centered semivaluation \(v\), define the Gauss weight
\[
 W_v\left(\sum_d b_du^d\right)
       :=\min_d\{v(b_d)+d\},
 \qquad
 W_v(\mathfrak b):=\min_{f\in\mathfrak b}W_v(f).
\tag{44.13}\label{eq:bj-char-zero-gauss-weight}
\]
Regular local coefficients have nonnegative value.
The normalized constants in \(\psi_0\) are nonpositive, so
\(\psi_0\le0\).
A direct calculation gives
\[
                         \psi_0=-\rho W(\mathfrak a).
\tag{44.14}\label{eq:bj-char-zero-weighted-function}
\]
If all \(s_j\) vanish at a point, then \(W(\mathfrak a)=+\infty\), and
\eqref{eq:bj-char-zero-weighted-function} reads \(\psi_0=-\infty\).

The direct calculation is the following, which we record as a lemma because
its sign is the one place where a slip would invert the whole construction.

\begin{lemma}
\label{lem:ch44-weight-encoding}
With the notation of
\eqref{eq:bj-char-zero-gf}--\eqref{eq:bj-char-zero-weighted-ideal}, and with
\(W(\mathfrak a)(v):=W_v(\mathfrak a)\) computed in any local frame of
\(\nu^*M\) at the center of \(v\), one has
\[
 \psi_0=-\rho\,W(\mathfrak a)
 \qquad\text{at every point of }Y^{\An}=X^{\An},
\]
with both sides in \([-\infty,0]\), and with both sides equal to
\(-\infty\) exactly on the common zero locus \(\bigcap_j\{s_j=0\}\).
Moreover \(W(\mathfrak a)\) is independent of the chosen frame.
\end{lemma}

\begin{proof}
Frame independence: two frames of \(\nu^*M\) at \(c(v)\) differ by a unit
\(g\) of the local ring, and \(v(g)=0\), so each \(v(f_{j})\) is unchanged;
this is \Cref{lem:ch34-frame-independence} in the present notation.

Fix \(v\) and a frame, and write \(s_j=f_j\cdot e\) locally.  Then
\(\log|s_j|(v)=-v(f_j)\), by the normalization recalled in
\Cref{sec:ch44-hypotheses}.  The generators of \(\mathfrak a\) are
\(u^{d_j}f_j^{\,q}\) in this frame, and by
\Cref{lem:ch44-gauss-valuation}\,\cref{itm:ch44-gauss:ideal},
\[
 W_v(\mathfrak a)
 =\min_j\bigl\{v(f_j^{\,q})+d_j\bigr\}
 =\min_j\bigl\{q\,v(f_j)+d_j\bigr\}.
\]
Substituting \(d_j=-q(\lambda_j-\max_\ell\lambda_\ell)\) and
\(v(f_j)=-\log|s_j|(v)\),
\[
 W_v(\mathfrak a)
 =q\min_j\Bigl\{-\log|s_j|(v)-\lambda_j+\max_\ell\lambda_\ell\Bigr\}
 =-q\Bigl(\max_j\bigl\{\log|s_j|(v)+\lambda_j\bigr\}-\max_\ell\lambda_\ell\Bigr),
\]
using \(\min_j(-x_j)=-\max_jx_j\).  Multiplying by
\(-\rho=-1/(qm)\) and recalling \(\psi=\frac1m\max_j\{\log|s_j|+\lambda_j\}\)
and \(c=\frac1m\max_\ell\lambda_\ell\),
\[
 -\rho\,W_v(\mathfrak a)
 =\frac1m\max_j\{\log|s_j|(v)+\lambda_j\}-\frac1m\max_\ell\lambda_\ell
 =\psi(v)-c=\psi_0(v).
\]
All the manipulations are between elements of \([-\infty,+\infty)\) and
involve no forbidden difference: the only subtraction is of the finite
rational constant \(\max_\ell\lambda_\ell\).

Signs and extreme values.  Since \(d_j\ge0\) and \(v(f_j)\ge0\) for a
centered \(v\) --- the coefficients being regular at the center --- we have
\(W_v(\mathfrak a)\ge0\), so \(\psi_0\le0\); this is the assertion that the
normalization by \(c\) makes \(\psi_0\) nonpositive.  Finally
\(W_v(\mathfrak a)=+\infty\) if and only if \(v(f_j)=+\infty\) for every
\(j\), i.e.\ if and only if every \(s_j\) lies in the support prime of
\(v\), i.e.\ if and only if \(v\) lies in the analytic common zero locus of
the \(s_j\) in the sense of \Cref{def:ch34-analytic-closed-subset}; and
there \(\psi=-\infty\) by \Cref{prop:ch34-generic-finiteness}.
\end{proof}

\begin{example}[The encoding, on two weights]
\label{ex:ch44-encoding-numbers}
Let \(m=1\), \(s=2\), \(\lambda_1=0\), \(\lambda_2=-\tfrac13\).  Then
\(c=\frac1m\max_j\lambda_j=0\), \(\psi_0=\psi\); the smallest admissible
\(q\) is \(3\), giving
\[
 d_1=-3(0-0)=0,
 \qquad
 d_2=-3\bigl(-\tfrac13-0\bigr)=1,
 \qquad
 \rho=\tfrac13,
 \qquad
 M=3L .
\]
The weighted system is \(\mathfrak a=(s_1^3,\ us_2^3)\).  At a semivaluation
\(v\) with \(v(f_1)=2\) and \(v(f_2)=\tfrac12\) one gets
\[
 W_v(\mathfrak a)=\min\{3\cdot2+0,\ 3\cdot\tfrac12+1\}=\min\{6,\tfrac52\}
 =\tfrac52,
 \qquad
 \psi_0(v)=-\tfrac13\cdot\tfrac52=-\tfrac56,
\]
in agreement with
\(\max\{\log|s_1|,\log|s_2|-\tfrac13\}=\max\{-2,-\tfrac12-\tfrac13\}
=-\tfrac56\).  The example is worth carrying because it shows the two roles
of \(q\): it clears the denominators of the weights, and it raises the
sections to a power so that the weights become integral \emph{shifts} rather
than fractional ones.  Nothing about \(\psi\) changes: \(q\) rescales
\(W\) and \(\rho\) inversely.
\end{example}

\section{Globalizing the local sandwiches}
\label{sec:ch44-globalization}

\Cref{lem:bj-local-additive-sandwich} is local: it produces, on each affine
chart, an ideal of the downstairs ring squeezed between two powers of the
upstairs ideal.  Three obstacles separate this from a family of global
sections of a line bundle on \(X\).

\begin{enumerate}
\item The Roberts--Singh coefficients \(c_{\alpha i}\) are local sections
      upstairs, defined only on \(\nu^{-1}U_\alpha\), and must be extended
      to \(Y\).
\item The elements of \(K_{n,\alpha}\) lie in \(A_\alpha[u]\), that is,
      they are functions on \(U_\alpha\); to be sections of a sheaf on \(X\)
      they must extend across \(X\setminus U_\alpha\).
\item The blocks from a single chart have different degrees
      \(m\in[n+p_\alpha,n+2p_\alpha]\), and the charts have different
      \(p_\alpha\); a single \emph{family} of sections of a single line
      bundle is wanted.
\end{enumerate}

Each obstacle is removed by twisting, and point~\textbf{(C2)} of
\Cref{sec:ch44-care} is the assertion that the three twists can be chosen
\emph{before} \(n\).  We first fix the geometric data, then record the three
uniformities as a single lemma, then carry out the globalization.

\subsection{The geometric data}
\label{sec:ch44-global-data}

Choose a finite affine cover
\[
                         U_\alpha=D(\tau_\alpha)
\tag{44.15}\label{eq:bj-char-zero-affine-cover}
\]
that trivializes \(M\), where all \(\tau_\alpha\) are sections of one ample
line bundle \(A_0\) on \(X\).
Let \(e_\alpha\) be a frame of \(M\) on \(U_\alpha\), and write
\[
                         s_j^q=f_{\alpha j}e_\alpha.
\]
Apply \Cref{lem:bj-local-additive-sandwich} to
\(g_{\alpha j}=u^{d_j}f_{\alpha j}\).
Write \(p_\alpha,N_\alpha\) for its constants.
The systems on different charts need not agree.
Each will be globalized separately and then placed into one finite family.

Such a cover exists: \(A_0\) ample means that for every closed point
\(x\in X\) some power \(rA_0\) has a section not vanishing at \(x\) whose
nonvanishing locus is affine and trivializes \(M\); replacing \(A_0\) by
\(rA_0\) and using quasi-compactness of \(X\) gives a finite such cover.  The
extension \(A_\alpha:=\cO_X(U_\alpha)\subset B_\alpha:=\cO_Y(\nu^{-1}U_\alpha)\)
is finite subintegral by \Cref{lem:ch44-char-zero-subintegral}, so
\Cref{lem:bj-local-additive-sandwich} does apply.

The generic element attached to chart \(\alpha\) is
\(b_\alpha=\sum_jz_jg_{\alpha j}\), and here the first piece of good luck of
the construction appears: although \(b_\alpha\) is a function on
\(\nu^{-1}U_\alpha\) depending on the frame, the \emph{section}
\[
 b_\alpha\,\nu^*e_\alpha
 =\sum_jz_ju^{d_j}f_{\alpha j}\nu^*e_\alpha
 =\sum_jz_ju^{d_j}s_j^{\otimes q}
 =:\mathbf s
\]
does not depend on \(\alpha\) at all.  It is one global section of
\(\nu^*M\) with coefficients in \(\mathbf Z[u,\mathbf z]\).  All the
chartwise systems are therefore attached to the same global object, and
differ only in the choice of Roberts--Singh data.  This is exactly the
covariance recorded in \eqref{eq:bj-char-zero-frame-law}, and it is what
allows the leading terms of the blocks to be written globally with no work
at all; only the correction terms need extending.

\subsection{Three uniform choices}
\label{sec:ch44-uniform-choices}

\begin{lemma}
\label{lem:ch44-uniform-cutoffs}
With the notation above there exist integers \(R,S\ge0\) and a line bundle
\(H_1\) on \(X\), all three depending only on \(\nu\), on the cover
\eqref{eq:bj-char-zero-affine-cover}, on \(M\), and on the finitely many
Roberts--Singh data \(\{c_{\alpha i}\}\) --- and in particular
\emph{not} on \(n\), nor on the block index \(m\) --- such that:
\begin{enumerate}
\item\label{itm:ch44-cut:R}
      for every \(\alpha\) and every \(1\le i\le p_\alpha\), each coefficient
      of \(\tau_\alpha^Rc_{\alpha i}e_\alpha^{\otimes i}\), taken separately
      in each \(\mathbf z\)-degree and each \(u\)-weight, extends to a
      global section of \(\nu^*(iM+RA_0)\) over \(Y\);
\item\label{itm:ch44-cut:S}
      for every \(\alpha\), the subsheaf
      \(\Gamma_{V(\tau_\alpha)}\mathcal Q\subseteq\mathcal Q\) of sections
      supported on \(V(\tau_\alpha)\) is annihilated by \(\tau_\alpha^S\),
      and the same \(S\) works after twisting \(\mathcal Q\) by any
      invertible sheaf;
\item\label{itm:ch44-cut:H1}
      each of the finitely many line bundles \(H_1+jM\), \(0\le j\le C\),
      with \(C=2\max_\alpha p_\alpha\), is globally generated.
\end{enumerate}
\end{lemma}

\begin{proof}
Item~\cref{itm:ch44-cut:R}.  Fix \(\alpha\) and \(i\).  The element
\(c_{\alpha i}\in B_\alpha[u,\mathbf z]\) is a polynomial in \(u\) and
\(\mathbf z\) with finitely many coefficients in \(B_\alpha\); multiplying by
the frame gives finitely many elements of
\(\cO_Y\bigl(\nu^{-1}U_\alpha,\nu^*(iM)\bigr)\).  Since
\(\nu^{-1}U_\alpha=\nu^{-1}D(\tau_\alpha)=D(\nu^*\tau_\alpha)\) and
\(\nu^*A_0\) is ample on \(Y\) (\(\nu\) being finite,
\Cref{lem:ch43-ample-pullback}), each such local section is killed into
global sections by a sufficiently high power of \(\nu^*\tau_\alpha\): this
is the standard statement that
\(H^0\bigl(Y,\nu^*(iM+rA_0)\bigr)\to
 H^0\bigl(\nu^{-1}U_\alpha,\nu^*(iM)\bigr)\),
with source mapped by division by \(\tau_\alpha^r\), has image exhausting
the target as \(r\to\infty\).  Take \(R\) to be the maximum of the finitely
many exponents so obtained, over the finitely many \(\alpha\), the finitely
many \(i\le p_\alpha\), and the finitely many coefficients.  Nothing here
involves \(m\) or \(n\): the \(c_{\alpha i}\) are the \emph{fixed} data of
\eqref{eq:bj-rs-additive-powers}, and their independence of \(m\) is
precisely the exponent-zero property (\Cref{rem:ch44-exponent-zero}).

Item~\cref{itm:ch44-cut:S}.  \(\mathcal Q=\nu_*\cO_Y/\cO_X\) is coherent,
being the cokernel of a map of coherent sheaves, and its subsheaf
\(\Gamma_{V(\tau_\alpha)}\mathcal Q\) of sections killed by some power of
\(\tau_\alpha\) is coherent, since \(X\) is noetherian; a coherent sheaf all
of whose local sections are killed by powers of \(\tau_\alpha\) is killed by
one power, by quasi-compactness and finite generation.  Let \(S_\alpha\) be
such a power and \(S:=\max_\alpha S_\alpha\).  For the last clause, let
\(\mathcal P\) be invertible.  Multiplication by \(\tau_\alpha\) on
\(\mathcal Q\otimes\mathcal P\) is the tensor of multiplication by
\(\tau_\alpha\) on \(\mathcal Q\) with the identity of \(\mathcal P\); since
\(\mathcal P\) is locally free of rank one, the kernel of the \(l\)-th power
is \(\bigl(\ker\tau_\alpha^l|_{\mathcal Q}\bigr)\otimes\mathcal P\), and the
annihilator of a sheaf is unchanged by tensoring with an invertible sheaf.
So the same \(S\) works for every twist, in particular for the twists
\(mM+RA_0\) that vary with the block index.  This is the crux of
item~\cref{itm:ch44-cut:S}: the exponent is uniform in the twist, hence
in \(m\) and \(n\).

Item~\cref{itm:ch44-cut:H1}.  The set of residual degrees is
\(\{0,1,\ldots,C\}\), a fixed finite set, because \(C=2\max_\alpha p_\alpha\)
depends only on the local Roberts--Singh lengths, which depend only on the
fixed generic elements \(b_\alpha\).  Choose an ample \(A_1\) on \(X\).  By
Serre's theorem, for each of the finitely many \(j\) with \(0\le j\le C\)
there is \(r_j\) such that the coherent sheaf \(\cO_X(jM)\otimes
A_1^{\otimes r}\) is globally generated for every \(r\ge r_j\).  Put
\(H_1:=r\,A_1\) with \(r:=\max_{0\le j\le C}r_j\); then \(H_1+jM\) is
globally generated for every \(0\le j\le C\), as required.  Again, no
dependence on \(n\): the index set \(\{0,\ldots,C\}\) is fixed once \(C\)
is.
\end{proof}

\begin{warning}{A twist that moves with \(n\) destroys the theorem}
\label{warn:ch44-moving-twist}
Suppose one of \(R,S\) were allowed to grow, say \(R=R(n)\).  The aligned
degree \eqref{eq:bj-char-zero-aligned-degree} would become
\(D_n=(n+C)M+H(n)\) with \(H(n)=(R(n)+S)A_0+H_1\), and the class of the
approximant \(F_n\) of \eqref{eq:bj-char-zero-gf-approximants} would be
\[
 c_1(L)+\frac{\rho}{n+C}\,c_1\bigl(H(n)\bigr).
\]
If \(R(n)\sim\lambda n\) with \(\lambda>0\), then, since
\(\rho(n+C)^{-1}R(n)\to\rho\lambda\), the classes converge to
\[
 c_1(L)+\rho\lambda\,c_1(A_0)\ \ne\ c_1(L),
\]
because \(A_0\) is ample and hence nonzero in \(N^1(X)\).  The construction
would then exhibit \(\psi_0\) as a decreasing limit of GF functions whose
classes converge to a \emph{strictly larger} class, proving only
\(\psi_0\in\PSH_X\bigl(c_1(L)+\rho\lambda\,c_1(A_0)\bigr)\) --- a true but
useless statement, since every \(\theta\)-psh function is
\(\theta'\)-psh for \(\theta'-\theta\) semiample and the assertion would
degenerate.  Nothing in the subsequent argument would visibly break; the
error would surface only in the class, which is why
\Cref{lem:ch44-uniform-cutoffs} is stated as a lemma with its uniformities
spelled out rather than left implicit.  The same remark applies to the
alignment bundle \(H_1\): if the set of residual degrees were unbounded, no
single \(H_1\) would serve, and this is why
\eqref{eq:bj-char-zero-block-degrees} bounds \(n+C-m\) by the constant
\(C\) and not by anything growing.
\end{warning}

\subsection{The blocks, globalized}
\label{sec:ch44-blocks}

Only finitely many coefficients occur in the fixed polynomials
\(c_{\alpha i}(u,\mathbf z)\).
Choose \(R\ge0\), independent of \(m\) and \(n\), so that every coefficient
of
\[
                         \tau_\alpha^R c_{\alpha i}e_\alpha^i
\tag{44.16}\label{eq:bj-char-zero-first-cutoff}
\]
extends to a global upstairs section of
\(\nu^*(iM+RA_0)\), separately in its \(\mathbf z\)-degree and \(u\)-weight.
Use these fixed extensions in the formula for
\(\tau_\alpha^Ra_{\alpha m}e_\alpha^m\).

These are structured extensions, not arbitrary extensions of the local sums.
A correction term is a fixed extension of \(c_{\alpha i}\) multiplied by
\(m-i\) original weighted sections.
The leading term is a product of \(m\) original sections.
For
\[
                 n+p_\alpha\le m\le n+2p_\alpha,
\]
every global block therefore remains in \(\mathfrak a^n\) everywhere on
\(Y\).

Written out, the global block is
\[
 \mathbf a_{\alpha m}
 :=\tau_\alpha^R\,\mathbf s^{\,m}
   +\sum_{i=1}^{p_\alpha}\binom mi
     \bigl[\tau_\alpha^Rc_{\alpha i}e_\alpha^{\otimes i}\bigr]^{\sim}
     \cdot\mathbf s^{\,m-i}
 \ \in\ H^0\bigl(Y,\nu^*(mM+RA_0)\bigr)[u,\mathbf z],
\]
where \([\,\cdot\,]^{\sim}\) denotes the extension fixed once and for all by
\Cref{lem:ch44-uniform-cutoffs}\,\cref{itm:ch44-cut:R}, and
\(\mathbf s=\sum_jz_ju^{d_j}s_j^{\otimes q}\) is the global section
constructed in \Cref{sec:ch44-global-data}.  Two properties of this formula
carry the argument, and they should be separated.

\emph{It restricts correctly.}  On \(\nu^{-1}U_\alpha\) the chosen
extensions restrict to \(\tau_\alpha^Rc_{\alpha i}e_\alpha^{\otimes i}\) by
construction, and \(\mathbf s=b_\alpha\nu^*e_\alpha\); hence
\(\mathbf a_{\alpha m}=\tau_\alpha^R\,a_{\alpha m}\,\nu^*e_\alpha^{\otimes m}\)
there, with \(a_{\alpha m}\) the local block of
\eqref{eq:bj-char-zero-am}.  In particular the \(\mathbf z\)-coefficients of
\(\mathbf a_{\alpha m}\), read in the frame \(e_\alpha\) and divided by
\(\tau_\alpha^R\), are exactly the local data whose \(u\)-components
generate \(K_{n,\alpha}\).

\emph{It stays in \(\mathfrak a^n\) globally.}  This is the point at which
the ``structured'' nature of the extension matters, and it deserves the
emphasis it is given above.  We did \emph{not} extend the local sum
\(a_{\alpha m}\); we extended the finitely many \(c_{\alpha i}\), once, and
then rebuilt the sum out of them and the global \(\mathbf s\).  Consequently
each summand of \(\mathbf a_{\alpha m}\) visibly carries a factor
\(\mathbf s^{\,m-i}\) with \(m-i\ge n\) \emph{on all of \(Y\)}, and the
\(\mathbf z\)-coefficients of \(\mathbf s^{\,m-i}\) are integer multiples of
monomials of degree \(m-i\ge n\) in the weighted generators
\(u^{d_j}s_j^{\otimes q}\) of \(\mathfrak a\).  Had we extended
\(a_{\alpha m}\) as an undifferentiated whole, the extension would agree
with a member of \(\mathfrak a^n\) only over \(U_\alpha\), and the upper
inclusion would be lost off that chart.

\begin{warning}{Extend the coefficients, not the blocks}
\label{warn:ch44-extend-coefficients}
The two operations
\[
 \text{extend }c_{\alpha i},\text{ then form the block}
 \qquad\text{versus}\qquad
 \text{form the block, then extend it}
\]
give sections with the same restriction to \(\nu^{-1}U_\alpha\) but
different behaviour elsewhere, and only the first is admissible.  There is
also an efficiency reason: the \(c_{\alpha i}\) are finitely many and fixed,
so one exponent \(R\) suffices for all of them and for all \(m\)
simultaneously, whereas the blocks \(a_{\alpha m}\) form an infinite family
indexed by \(m\), and the exponents needed to extend them could well grow
with \(m\).  The reader should compare
\Cref{warn:ch44-moving-twist}: extending the blocks is exactly the route by
which a twist depending on \(n\) would enter.
\end{warning}

To descend the blocks, consider the coherent quotient
\[
            \mathcal Q:=\nu_*\cO_Y/\cO_X.
\tag{44.17}\label{eq:bj-char-zero-quotient}
\]
The class in \(\mathcal Q\) of a block built from \(U_\alpha\) vanishes after
localizing at \(\tau_\alpha\).
The kernels of
\[
 \mathcal Q\xrightarrow{\ \tau_\alpha^s\ }
       \mathcal Q\otimes\cO_X(sA_0)
\tag{44.18}\label{eq:bj-char-zero-stable-kernels}
\]
form an ascending chain of coherent subsheaves and hence stabilize.
Tensoring by an invertible sheaf does not change the stabilized-kernel
calculation.
One exponent \(S\), independent of the varying powers \(m,n\), therefore
kills the quotient classes of all these blocks.

In more detail, and this is
\Cref{lem:ch44-uniform-cutoffs}\,\cref{itm:ch44-cut:S} in use.  Each
\(\mathbf z\)- and \(u\)-coefficient \(\sigma\) of \(\mathbf a_{\alpha m}\)
is a global section of \(\nu^*(mM+RA_0)\) over \(Y\), that is, a global
section of \(\nu_*\cO_Y\otimes\mathcal P\) with
\(\mathcal P:=\cO_X(mM+RA_0)\) --- here the projection formula for the
finite map \(\nu\) is used, in the innocuous form
\(\nu_*\nu^*\mathcal P\cong\nu_*\cO_Y\otimes\mathcal P\).  Its image
\(\bar\sigma\in H^0(X,\mathcal Q\otimes\mathcal P)\) restricts to zero on
\(U_\alpha\), because there \(\sigma\) equals
\(\tau_\alpha^Ra_{\alpha m}e_\alpha^{\otimes m}\) with
\(a_{\alpha m}\in A_\alpha[u,\mathbf z]\), a downstairs function.  Hence
\(\bar\sigma\) is a section of \(\mathcal Q\otimes\mathcal P\) supported on
\(V(\tau_\alpha)\), so \(\tau_\alpha^S\bar\sigma=0\).

After multiplication by \(\tau_\alpha^S\), the exact sequence
\[
 0\longrightarrow\cO_X\longrightarrow\nu_*\cO_Y
       \longrightarrow\mathcal Q\longrightarrow0
\tag{44.19}\label{eq:bj-char-zero-quotient-exact}
\]
shows that each coefficient is pulled back from a downstairs global section.
There is no \(H^1\)-vanishing assertion here.
After tensoring \eqref{eq:bj-char-zero-quotient-exact} by the relevant line
bundle, exactness at the middle term says precisely that an upstairs global
section with zero quotient class is in the image of the downstairs global
sections.

Spelled out: \eqref{eq:bj-char-zero-quotient-exact} tensored with the
invertible \(\mathcal P':=\cO_X\bigl(mM+(R+S)A_0\bigr)\) stays exact,
\(\mathcal P'\) being locally free, and taking global sections gives the
left-exact sequence
\[
 0\to H^0(X,\mathcal P')\to H^0(X,\nu_*\cO_Y\otimes\mathcal P')
   \to H^0(X,\mathcal Q\otimes\mathcal P').
\]
The middle term is \(H^0\bigl(Y,\nu^*\mathcal P'\bigr)\), the element
\(\tau_\alpha^S\sigma\) lies in it, and its image in the right-hand term
vanishes; hence it comes from a \emph{unique} element of
\(H^0(X,\mathcal P')\).  Uniqueness is what makes the descended section
well defined; existence is what makes it exist.  No cohomology beyond
\(H^0\) appears, and in particular no vanishing theorem is invoked.

\begin{remark}[Not a torsion-kernel descent argument]
\label{rem:ch44-not-torsion-kernel}
The quotient \(\mathcal Q\) measures failure of coefficients to lie
downstairs.
The argument kills their quotient classes by fixed cutoffs.
It neither removes torsion from a pulled-back coherent sheaf nor asserts that
a finite pushdown equals the original sheaf.
No unsupported normalization-kernel principle is being used.
\end{remark}

\begin{remark}[\(\mathcal Q\) is exactly what obstructs, and it is not small]
\label{rem:ch44-Q-size}
By \Cref{prop:ch44-seminormal-unibranch-normal}\,(3) the sheaf
\(\mathcal Q\) is nonzero whenever \(\nu\) is not an isomorphism, so the
cutoff by \(\tau_\alpha^S\) is not a formality.  For the cusp
\(A=k[t^2,t^3]\subset B=k[t]\) one has \(\mathcal Q\cong k\), a skyscraper of
length one at the cusp, and \(S=1\) suffices: with
\(\mathfrak m=(t^2,t^3)\) one has
\(\mathfrak m\mathcal Q=0\), since \(t^2\cdot t=t^3\in A\), so
\(\mathcal Q\) is killed by any \(\tau_\alpha\in\mathfrak m\).  For
\(A=k[t^2,t^5]\) the quotient \(\mathcal Q\) is a skyscraper of length two,
spanned by the classes of \(t\) and \(t^3\), and now \(S=1\) does
\emph{not} suffice: \(t^2\cdot t=t^3\notin A\), so
\(\mathfrak m\mathcal Q\ne0\), while
\(\mathfrak m^2\mathcal Q=0\) because \(t^4\cdot t=t^5\in A\) and
\(t^2\cdot t^3=t^5\in A\).  There \(S=2\) is needed.  The example shows that
\(S\) measures the depth of the \(\tau_\alpha\)-power filtration on
\(\mathcal Q\), which is bounded by but not equal to its length.  What
\(S\) genuinely bounds, in general, is the least \(l\) with
\(\tau_\alpha^l\Gamma_{V(\tau_\alpha)}\mathcal Q=0\), and it is finite for
the sole reason that \(\mathcal Q\) is coherent on a noetherian scheme.
\end{remark}

\section{One degree and one global estimate}
\label{sec:ch44-global-sandwich}

Define
\[
 \begin{aligned}
 P&:=\max_\alpha p_\alpha,
 &N_0&:=\max_\alpha\max\{0,N_\alpha-p_\alpha\},\\
 C&:=2P,
 &H_0&:=(R+S)A_0.
 \end{aligned}
\tag{44.20}\label{eq:bj-char-zero-global-constants}
\]
All four quantities are determined by \(\nu\), the cover, \(M\) and the
finitely many local Roberts--Singh systems.  None of them depends on \(n\),
and this is the substance of point~\textbf{(C2)} of \Cref{sec:ch44-care};
\(H_0\) and \(H_1\) are supplied by
\Cref{lem:ch44-uniform-cutoffs}\,\cref{itm:ch44-cut:R},
\cref{itm:ch44-cut:S} and \cref{itm:ch44-cut:H1} respectively.

For \(n\ge N_0\), the local estimate on \(U_\alpha\) includes
\[
 \mathfrak a_\alpha^{\,n+C}
 \subset\mathfrak a_\alpha^{\,n+2p_\alpha}
 \subset K_{n,\alpha}B_\alpha[u]
 \subset\mathfrak a_\alpha^n.
\tag{44.21}\label{eq:bj-char-zero-uniform-local-sandwich}
\]
The first inclusion in \eqref{eq:bj-char-zero-uniform-local-sandwich} holds
because \(C=2P\ge2p_\alpha\) and ideal powers are decreasing; the middle two
are \eqref{eq:bj-local-additive-sandwich}, applicable because
\(n\ge N_0\ge\max\{0,N_\alpha-p_\alpha\}\) for every \(\alpha\).  The
passage from the chart-dependent exponent \(2p_\alpha\) to the uniform
\(C\) is where the finitely many local systems are reconciled, and it costs
nothing: it merely weakens the lower inclusion, which is the direction that
can afford to be weakened.

The descended blocks have line-bundle degrees \(mM+H_0\), where
\[
       n+p_\alpha\le m\le n+2p_\alpha,
       \qquad
       0\le n+C-m\le C.
\tag{44.22}\label{eq:bj-char-zero-block-degrees}
\]

Choose a line bundle \(H_1\) so positive that each of the finitely many line
bundles
\[
                         H_1+jM,\qquad0\le j\le C,
\]
is globally generated.
Multiply every block by a complete generating system in the appropriate
residual degree.
All blocks then have common degree
\[
                         D_n:=(n+C)M+H,
             \qquad
                         H:=H_0+H_1.
\tag{44.23}\label{eq:bj-char-zero-aligned-degree}
\]
Because the aligning sections are regular downstairs sections, the global
upper containment in \(\mathfrak a^n\) is preserved.

Let \(\mathfrak K_n\) be the resulting finite homogeneous family of
downstairs weighted sections of \(D_n\).
Every analytic semivaluation on projective \(X\) is centered.
At its center choose a chart \(U_\alpha\).
Then \(\tau_\alpha\) is a unit and has value zero.
Choose from each complete aligning system a section nonvanishing at that
center; it also has value zero.
Neither cutoff nor alignment changes the local lower estimate.
Therefore, at every analytic semivaluation,
\[
       nW(\mathfrak a)
       \le W(\nu^{-1}\mathfrak K_n)
       \le(n+C)W(\mathfrak a).
\tag{44.24}\label{eq:bj-char-zero-global-weight-sandwich}
\]

If \(W(\mathfrak a)=+\infty\), the inclusion
\(\nu^{-1}\mathfrak K_n\subset\mathfrak a^n\) forces the middle term of
\eqref{eq:bj-char-zero-global-weight-sandwich} to be \(+\infty\).
Thus all three entries are \(+\infty\).
This is the proper-support common-zero clause of the estimate.

We now record the outcome of the construction as a statement, and give its
proof in the order in which the estimate is actually verified.  The reader
who has followed \Cref{sec:ch44-globalization} will recognize every
ingredient; the point of restating is that this proposition, and not the
construction, is what the rest of the chapter uses.

\begin{proposition}
\label{prop:ch44-global-sandwich}
Let \(\operatorname{char}k=0\), let \(\nu:Y\to X\) be as in
\eqref{eq:bj-char-zero-universal-map}, let \(L\in\Pic(X)_{\mathbf Q}\), and
let \(\psi\in\cH^{\mathrm{gf}}_{\mathbf Q}(\nu^*L)\) be presented as in
\eqref{eq:bj-char-zero-gf}, with associated data
\(q,\rho,M,\mathfrak a\) as in
\eqref{eq:bj-char-zero-weight-exponents}--\eqref{eq:bj-char-zero-weighted-ideal}.
Then there exist
\begin{itemize}
\item an integer \(C\ge0\) and an integer \(N_0\ge0\),
\item a line bundle \(H\) on \(X\),
\end{itemize}
all depending only on the presentation \eqref{eq:bj-char-zero-gf} and on
\(\nu\) --- and in particular not on \(n\) --- together with, for every
\(n\ge N_0\), a finite family \(\mathfrak K_n\) of \(u\)-homogeneous
weighted global sections of
\[
 D_n=(n+C)M+H
\]
on \(X\), such that at every semivaluation \(v\in X^{\An}=Y^{\An}\)
\[
 n\,W_v(\mathfrak a)
 \ \le\ W_v\bigl(\nu^{-1}\mathfrak K_n\bigr)
 \ \le\ (n+C)\,W_v(\mathfrak a)
\]
in \([0,+\infty]\).  In particular, all three quantities are \(+\infty\)
simultaneously, and this happens exactly on the common zero locus of the
\(s_j\).
\end{proposition}

\begin{proof}
Take \(C\), \(N_0\), \(H_0\) as in
\eqref{eq:bj-char-zero-global-constants}, \(H_1\) as in
\Cref{lem:ch44-uniform-cutoffs}\,\cref{itm:ch44-cut:H1},
\(H:=H_0+H_1\), and \(\mathfrak K_n\) as constructed above: the family of
all \(u\)-homogeneous components of all \(\mathbf z\)-coefficients of all
aligned descended blocks.  Independence of \(n\) is
\Cref{lem:ch44-uniform-cutoffs} together with the observation that \(C\) and
\(N_0\) are maxima over the finitely many charts of data attached to the
fixed generic elements \(b_\alpha\).

\emph{The lower bound.}  Every member of \(\mathfrak K_n\) pulls back into
the ideal \(\mathfrak a^n\subseteq B[u]\), at every point of \(Y\): this was
verified in \Cref{sec:ch44-blocks} for the blocks themselves, is preserved
by extracting \(\mathbf z\)- and \(u\)-components (because \(\mathfrak a^n\)
is homogeneous for both gradings, the generators \(u^{d_j}s_j^{\otimes q}\)
being \(u\)-homogeneous and \(\mathbf z\)-free), and is preserved by
multiplication by the aligning sections, which are regular sections of
globally generated bundles pulled back from \(X\).  Fix \(v\) and let
\(\sigma\in\mathfrak K_n\).  By
\Cref{lem:ch44-gauss-valuation}\,\cref{itm:ch44-gauss:ideal} and
\cref{itm:ch44-gauss:mono},
\[
 W_v(\nu^*\sigma)\ \ge\ W_v(\mathfrak a^n)=n\,W_v(\mathfrak a),
\]
the last equality being multiplicativity of the Gauss weight on ideal
powers.  Taking the minimum over \(\sigma\) gives the lower bound.  Note
that this half needs no information about \(v\) beyond \(W_v\ge0\).

\emph{The upper bound.}  Here the center is used.  Since \(X\) is
projective, every \(v\in X^{\An}\) is centered
(\Cref{thm:ch32-support-center-maps}); let \(x:=c(v)\) and choose
\(\alpha\) with \(x\in U_\alpha\).  Then \(\tau_\alpha\) does not vanish at
\(x\), so \(\tau_\alpha\) is a unit in the local ring at \(x\) and
\(v(\tau_\alpha)=0\); consequently the factor \(\tau_\alpha^{R+S}\)
introduced by the two cutoffs contributes \(0\) to every Gauss weight
computed at \(v\).  Likewise, each aligning system generates a globally
generated bundle, so some member of it is nonvanishing at \(x\) and has
\(v\)-value \(0\); choosing that member, the alignment contributes \(0\) as
well.  For that choice, the members of \(\mathfrak K_n\) coming from chart
\(\alpha\) compute at \(v\) exactly the Gauss weights of the local
generators of \(K_{n,\alpha}\).  Hence, using
\eqref{eq:bj-char-zero-uniform-local-sandwich} and
\Cref{lem:ch44-gauss-valuation}\,\cref{itm:ch44-gauss:mono},
\[
 W_v\bigl(\nu^{-1}\mathfrak K_n\bigr)
 \le W_v\bigl(K_{n,\alpha}B_\alpha[u]\bigr)
 \le W_v\bigl(\mathfrak a^{\,n+2p_\alpha}\bigr)
 =(n+2p_\alpha)\,W_v(\mathfrak a)
 \le(n+C)\,W_v(\mathfrak a),
\]
the last step because \(2p_\alpha\le C\) and \(W_v(\mathfrak a)\ge0\).

\emph{The common-zero clause.}  If \(W_v(\mathfrak a)<+\infty\) the two
bounds are finite and there is nothing to check.  If
\(W_v(\mathfrak a)=+\infty\), the lower bound gives
\(W_v(\nu^{-1}\mathfrak K_n)=+\infty\) as well, and the upper bound reads
\(+\infty\le+\infty\); so all three are \(+\infty\).  By
\Cref{lem:ch44-weight-encoding} this happens exactly where all \(s_j\)
vanish.  No forbidden extended-real operation occurs: the products
\(n\cdot(+\infty)\) and \((n+C)\cdot(+\infty)\) have \(n,n+C\ge1\) --- for
\(n\ge N_0\) and \(n\ge1\), and for \(n=0\) the lower bound is the trivial
\(0\le\,\cdot\) --- and no difference of infinities is formed.
\end{proof}

\begin{remark}[The two halves use different hypotheses]
\label{rem:ch44-two-halves}
The lower bound in \Cref{prop:ch44-global-sandwich} is a formal consequence
of an ideal inclusion valid on all of \(Y\), and would survive almost any
mangling of the construction.  The upper bound is where the work is: it
requires a chart adapted to the point \(v\), and it requires that passing to
that chart cost nothing valuatively, which is why the cutoff was performed
with powers of \(\tau_\alpha\) --- a section of a bundle on \(X\) that is
invertible precisely on \(U_\alpha\) --- and why the alignment was performed
with \emph{complete} generating systems rather than with a single chosen
section.  A reader auditing the proof should check exactly these two
choices; everything else is bookkeeping.
\end{remark}

\begin{remark}[Comparison with the characteristic-\(p\) statement]
\label{rem:ch44-compare-frobenius}
In \Cref{chap:bj-universal-homeomorphisms} the analogous statement is an
equality of functions,
\eqref{eq:bj-positive-gf-identity}, obtained after a \emph{single} step; the
number \(q\) there is an exponent applied once, and there is no family
indexed by \(n\), no twist \(H\), and no constant \(C\).  The reason the
present statement is unavoidably weaker is
\Cref{rem:ch43-charzero-contrast}: there is no \(q\) with
\(b^q\in A\) for all \(b\).  What
\Cref{prop:ch44-global-sandwich} shows is that the corrected powers of
\eqref{eq:bj-rs-additive-powers}, though they do not compute \(\psi_0\),
compute it to within a factor \(r_n=n/(n+C)\to1\), and that the error is
uniform over \(X^{\An}\) --- including at the points where \(\psi_0=-\infty\),
which is the clause that would be lost if one worked with generic points
only.
\end{remark}

\section{A decreasing psh approximation}
\label{sec:ch44-decreasing}

Set
\[
 r_n:=\frac n{n+C},
 \qquad
 F_n:=-\frac{\rho}{n+C}W(\mathfrak K_n)
       \in\cH^{\mathrm{gf}}_{\mathbf Q}
          \left(L+\frac{\rho}{n+C}H\right).
\tag{44.25}\label{eq:bj-char-zero-gf-approximants}
\]
The class follows exactly from
\eqref{eq:bj-char-zero-aligned-degree}, since \(\rho M=L\).
Each \(u^d\)-component contributes coefficient weight \(-d\).
Combining
\eqref{eq:bj-char-zero-weighted-function} and
\eqref{eq:bj-char-zero-global-weight-sandwich} gives
\[
                         \psi_0\le F_n\le r_n\psi_0.
\tag{44.26}\label{eq:bj-char-zero-function-sandwich}
\]

Three verifications are owed here, and we give them.

\emph{\(F_n\) is a GF function of the stated class.}  The family
\(\mathfrak K_n\) consists of \(u\)-homogeneous weighted global sections of
\(D_n=(n+C)M+H\) on \(X\); writing a member as \(u^{d}\sigma\) with
\(\sigma\in H^0(X,D_n)\), the definition
\eqref{eq:bj-char-zero-gauss-weight} gives
\(W_v(u^d\sigma)=v(\sigma)+d=-\log|\sigma|(v)+d\), so
\[
 -\frac{\rho}{n+C}\,W(\mathfrak K_n)
 =\frac{\rho}{n+C}\max_{u^d\sigma\in\mathfrak K_n}
   \bigl\{\log|\sigma|-d\bigr\},
\]
which is exactly the presentation \eqref{eq:gf-fs-presentation} of a
generically finite Fubini--Study function for the \(\mathbf Q\)-line bundle
\(\frac{\rho}{n+C}D_n\), with denominator \((n+C)/\rho\) and rational
weights \(-d\).  Its class is
\[
 \frac{\rho}{n+C}\,c_1(D_n)
 =\rho\,c_1(M)+\frac{\rho}{n+C}c_1(H)
 =c_1(L)+\frac{\rho}{n+C}c_1(H),
\]
since \(\rho M=\frac1{qm}\cdot qmL=L\).  Generic finiteness holds because
\(\mathfrak K_n\) is nonempty and its members are not all identically zero:
the leading term of the block for \(m=n+p_\alpha\) has
\(\mathbf z\)-coefficients that are nonzero multiples of monomials in the
\(s_j\), which are nonzero on \(Y\), hence descend to nonzero sections.

\emph{The sandwich \eqref{eq:bj-char-zero-function-sandwich}.}  Multiply
\Cref{prop:ch44-global-sandwich} by \(-\rho/(n+C)<0\), which reverses the
inequalities:
\[
 -\rho\,W(\mathfrak a)
 \le-\frac{\rho}{n+C}W(\nu^{-1}\mathfrak K_n)
 \le-\frac{n\rho}{n+C}W(\mathfrak a),
\]
that is, \(\psi_0\le F_n\le r_n\psi_0\) by
\Cref{lem:ch44-weight-encoding}.  Here we have used that \(W\) of a family
of downstairs sections, computed upstairs after pullback, is the same number
as computed downstairs --- which is the content of the analytic
identification \Cref{lem:bj-universal-analytic-homeomorphism} together with
\Cref{lem:ch43-identification-dictionary}: the two spaces are the same
space, and a semivaluation takes the same value on a downstairs section and
on its pullback.  This is the only place where the topological half of
\Cref{chap:bj-universal-homeomorphisms} is invoked, and it is invoked
pointwise.

\emph{The inequalities are consistent at \(-\infty\).}  Where
\(\psi_0=-\infty\), all three members of
\eqref{eq:bj-char-zero-function-sandwich} are \(-\infty\), by the
common-zero clause of \Cref{prop:ch44-global-sandwich}; the products
\(r_n\cdot(-\infty)=-\infty\) are legitimate because \(r_n>0\).  Where
\(\psi_0\) is finite, all three are finite.  No comparison of the form
\(-\infty\le-\infty\) is ever converted into an equality of functions except
in this sense.

\subsection{Manufacturing monotonicity}
\label{sec:ch44-monotonization}

The family \(F_n\) need not decrease as \(n\) increases.
Choose rational \(a_j\uparrow1\) and integers \(n_j\uparrow\infty\), with
\(n_j\ge N_0\), such that
\[
                         a_{j+1}r_{n_{j+1}}\ge a_j.
\tag{44.27}\label{eq:bj-char-zero-monotonicity-threshold}
\]
Since \(\psi_0\le0\), the two sides of
\eqref{eq:bj-char-zero-function-sandwich} give
\[
 a_{j+1}F_{n_{j+1}}
 \le a_{j+1}r_{n_{j+1}}\psi_0
 \le a_j\psi_0
 \le a_jF_{n_j}.
\tag{44.28}\label{eq:bj-char-zero-monotone-chain}
\]
Thus \(a_jF_{n_j}\) is decreasing.
The same sandwich shows pointwise convergence to \(\psi_0\).
If \(\psi_0=-\infty\), both bounds force every selected approximant to equal
\(-\infty\) there.

The choice \eqref{eq:bj-char-zero-monotonicity-threshold} is possible, and
it is worth carrying out once so that the reader sees how little room there
is.  Since \(r_n=n/(n+C)\uparrow1\), pick \(a_j:=1-2^{-j}\) --- rational and
increasing to \(1\) --- and then choose \(n_j\) recursively: given \(n_j\),
choose \(n_{j+1}>n_j\) with
\[
 r_{n_{j+1}}\ \ge\ \frac{a_j}{a_{j+1}}
 =\frac{1-2^{-j}}{1-2^{-j-1}},
\]
which is possible because the right-hand side is \(<1\) and
\(r_n\to1\); explicitly \(r_n\ge\lambda\) as soon as
\(n\ge C\lambda/(1-\lambda)\), so
\[
 n_{j+1}
 :=\max\Bigl\{n_j+1,\ N_0,\
  \bigl\lceil C\,(2^{j+1}-2)\bigr\rceil\Bigr\}
\]
does it, since \(\lambda/(1-\lambda)=(2^{j+1}-2)\) for
\(\lambda=a_j/a_{j+1}\).  The essential feature of this choice is that
\(n_j\) grows like \(2^j\): the monotonization forces the approximants to be
taken along a \emph{sparse} subsequence, which is harmless here because
nothing is claimed about the rate.

The chain \eqref{eq:bj-char-zero-monotone-chain} deserves one word.  Its
first inequality is the right half of
\eqref{eq:bj-char-zero-function-sandwich} scaled by \(a_{j+1}>0\); its
second is \eqref{eq:bj-char-zero-monotonicity-threshold} multiplied by
\(\psi_0\le0\), which reverses the inequality; its third is the left half of
\eqref{eq:bj-char-zero-function-sandwich} scaled by \(a_j>0\).  The sign of
\(\psi_0\) is used exactly once, in the middle step, and this is why
\eqref{eq:bj-char-zero-gf} was normalized by subtracting \(c\) in the first
place.  If one forgot to normalize, \(\psi\) could be positive somewhere and
the middle inequality would reverse there, destroying monotonicity.

The classes are
\[
             a_jL+\frac{a_j\rho}{n_j+C}H
                         \longrightarrow L.
\tag{44.29}\label{eq:bj-char-zero-approximating-classes}
\]
At the trivial valuation, some nonzero \(s_j\) has \(d_j=0\).
Its coefficient has value zero, so \(W(\mathfrak a)=0\).
The sandwich forces every normalized approximant to have value zero there.
PSH2 applied to the decreasing sequence and
\eqref{eq:bj-char-zero-approximating-classes} therefore gives
\[
                         \psi_0\in\PSH_X(L),
\]
not the identically \(-\infty\) alternative.
Adding back the rational constant \(c\) proves the following proposition.

Here ``PSH2'' is \Cref{thm:bj-basic-closure}\,(5), closure of \(\PSH\) under
decreasing nets with \emph{varying} classes; the classes
\eqref{eq:bj-char-zero-approximating-classes} do converge to \(c_1(L)\)
because \(a_j\to1\) and \(\rho\,c_1(H)/(n_j+C)\to0\), the latter because
\(H\) is a fixed line bundle (\Cref{lem:ch44-uniform-cutoffs},
\Cref{warn:ch44-moving-twist}).  The exclusion of the identically
\(-\infty\) alternative is the hypothesis of generic finiteness in
\Cref{def:bj-psh}(2), and is verified at the single point \(v_\triv\):
indeed at \(v_\triv\) every nonzero regular function has value zero, so with
\(j_0\) an index attaining \(\max_\ell\lambda_\ell\) --- so that
\(d_{j_0}=0\) --- one gets \(W_{v_\triv}(\mathfrak a)\le
v_\triv(s_{j_0}^{\otimes q})+0=0\), while \(W_{v_\triv}\ge0\) always; hence
\(W_{v_\triv}(\mathfrak a)=0\), and
\eqref{eq:bj-char-zero-function-sandwich} gives
\(0=\psi_0(v_\triv)\le a_jF_{n_j}(v_\triv)\le0\).  Thus every member of the
decreasing sequence takes the value \(0\) at \(v_\triv\), and so does the
limit.  This single-point check is the standard device of the chapter and
recurs in \Cref{thm:bj-char-zero-universal-reflection}; it is also the
device of \Cref{rem:ch43-firewall-not-vacuous}, and it is what keeps the
appeal to \Cref{thm:bj-basic-closure}\,(5) honest.

\begin{proposition}\label{prop:bj-char-zero-stage-reflection}
Under \eqref{eq:bj-char-zero-universal-map}, every
\[
                 \psi\in\cH^{\mathrm{gf}}_{\mathbf Q}(\nu^*L)
\]
is, under the analytic identification, a member of \(\PSH_X(L)\).
The resulting decreasing approximation retains the full
\(-\infty\)-locus of \(\psi\).
\end{proposition}

\begin{proof}
The construction and PSH2 argument above apply to an arbitrary presentation
\eqref{eq:bj-char-zero-gf}.
The conclusion is independent of the chosen presentation because the
pointwise limit is the given function \(\psi\).

For the second sentence, observe that at a point \(v\) with
\(\psi(v)=-\infty\) --- equivalently \(\psi_0(v)=-\infty\), \(c\) being a
finite constant --- the sandwich
\eqref{eq:bj-char-zero-function-sandwich} forces \(F_{n_j}(v)=-\infty\) for
every \(j\), hence \(a_jF_{n_j}(v)+c=-\infty\); and conversely at a point
where \(\psi(v)>-\infty\), the same sandwich gives
\(a_jF_{n_j}(v)\ge a_j\psi_0(v)>-\infty\).  So the approximants have exactly
the same polar locus as \(\psi\), not a larger one.  This is the assertion
that no analytic point is lost, and it is the characteristic-zero analogue
of the corresponding clause in
\Cref{prop:bj-frobenius-gf-descent}; there it was automatic from an
identity, here it is a consequence of the two-sided estimate.
\end{proof}

\begin{remark}[The approximants are downstairs GF functions, not merely psh]
\label{rem:ch44-approximants-are-gf}
\Cref{prop:bj-char-zero-stage-reflection} asserts membership in
\(\PSH_X(L)\), and \emph{not} in
\(\cH^{\mathrm{gf}}_{\mathbf Q}(X,L)\).  The difference is real and is a
genuine divergence from the positive-characteristic picture, where
\Cref{cor:ch43-gf-bijection} gives a bijection at the level of GF functions
themselves.  Here the descended objects are the \(a_jF_{n_j}\), which are
downstairs GF functions but of classes \(a_jL+\frac{a_j\rho}{n_j+C}H\)
strictly larger than \(L\); the function \(\psi_0\) itself is recovered only
in the limit and need not be a downstairs GF function at all.  Consequently
no statement of the form ``\(\nu^*\) is a bijection of GF spaces'' is
available in characteristic zero, and none is claimed;
\Cref{warn:ch44-no-gf-bijection} records what fails.
\end{remark}

\begin{warning}{The characteristic-\(p\) mechanism fails at its first step}
\label{warn:ch44-no-gf-bijection}
We do \emph{not} assert that
\(\nu^*:\cH^{\mathrm{gf}}_{\mathbf Q}(X,L)\to
 \cH^{\mathrm{gf}}_{\mathbf Q}(Y,\nu^*L)\) fails to be surjective, and we
prove nothing about that question; the point is a different and sharper one,
namely that the \emph{route} taken in
\Cref{prop:bj-frobenius-gf-descent} is unavailable, at its very first step,
already in the simplest example.

That route consists in producing, for each upstairs section \(s_j\), a
downstairs section \(s_j^{[q]}\) with
\(\nu^*s_j^{[q]}=s_j^{\otimes q}\): the local coefficients
\(f_{U,j}\in B_U\) satisfy \(f_{U,j}^{\,q}\in A_U\)
(\Cref{cor:ch43-frobenius-all-affines}), so their \(q\)-th powers glue.
Take \(A=\mathbf Q[t^2,t^3]\subset B=\mathbf Q[t]\) and \(f=1+t\).  Then for
every \(q\ge1\),
\[
 [t^1]\bigl((1+t)^q\bigr)=q\neq0
 \quad\text{in }\mathbf Q,
 \qquad\text{so}\qquad
 (1+t)^q\notin A,
\]
by \Cref{rem:ch44-cusp-decomposition}.  Thus \emph{no} power of this one
local coefficient descends, and there is no candidate \(s^{[q]}\) whatever.
Two further remarks.  First, the failure is invisible on monomials: for
\(f=t\) one has \(t^q\in A\) for every \(q\ge2\).  Second, it is not
repaired by enlarging \(q\) or by passing to a Veronese multiple, since the
computation above is valid for every \(q\).  This is exactly the
observation recorded from the other side in
\Cref{rem:ch43-charzero-contrast}, and it is why the present chapter
replaces an identity by a limit.
\end{warning}

\section{Psh reflection in characteristic zero}
\label{sec:ch44-reflection}

\Cref{prop:bj-char-zero-stage-reflection} treats a single GF stage whose
class is already a pullback.  Two things remain: the classes of a
regularizing net are only \emph{numerically} close to \(\nu^*\theta\), and
must be promoted to honest pullbacks; and the stages must be assembled.  The
first is not proved here --- it is characteristic-free and was done in
\Cref{lem:bj-normalization-class-promotion} --- and the second is one
application of \Cref{thm:bj-basic-closure}\,(5).

\begin{theorem}\label{thm:bj-char-zero-universal-reflection}
Assume \(\operatorname{char}k=0\).
For every \(\theta\in N^1(X)_{\mathbf R}\), pullback is a bijection
\[
 \nu^*:\PSH_X(\theta)
       \xrightarrow{\ \sim\ }
       \PSH_Y(\nu^*\theta).
\tag{44.30}\label{eq:bj-char-zero-universal-reflection}
\]
\end{theorem}

\begin{proof}
Pullback preserves psh functions and classes by
\cite[Theorem~4.7]{BJ22}, and it is injective because \(\nu^{\An}\) is a
homeomorphism.
Let \(\psi\in\PSH_Y(\nu^*\theta)\).
Choose a decreasing GF regularization
\[
 \psi_i\in\cH^{\mathrm{gf}}_{\mathbf Q}(M_i),
 \qquad
 \psi_i\downarrow\psi,
 \qquad
 c_1(M_i)\longrightarrow\nu^*\theta.
\tag{44.31}\label{eq:bj-char-zero-regularization}
\]
By \Cref{lem:bj-normalization-class-promotion}, after cofinal refinement
each \(M_i=\nu^*L_i\) exactly and \(c_1(L_i)\to\theta\).
\Cref{prop:bj-char-zero-stage-reflection} regards the same pointwise function
\(\psi_i\) as an \(L_i\)-psh function downstairs.
The sequence remains decreasing under the analytic identification and
converges to \(\psi\) at every point.
The trivial valuation excludes the identically \(-\infty\) alternative.
PSH2 yields a member of \(\PSH_X(\theta)\) whose pullback is \(\psi\).

We expand the last four sentences, since they carry the whole assembly.

\emph{Injectivity and preservation.}  If \(\phi\in\PSH_X(\theta)\), then
\(\nu^*\phi=\phi\circ\nu^{\An}\) is the same function read on the same
space, and it is \(\nu^*\theta\)-psh: an approximating net
\(\phi_i\in\cH^{\mathrm{gf}}_{\mathbf R}(L_i)\) pulls back to
\(\nu^*\phi_i\in\cH^{\mathrm{gf}}_{\mathbf R}(\nu^*L_i)\) --- a GF
presentation pulls back termwise --- and the classes \(\nu^*c_1(L_i)\)
converge to \(\nu^*\theta\) since \(\nu^*\) is linear and continuous on
\(N^1\).  Generic finiteness is preserved because \(\nu^{\An}\) is
surjective.  This is \Cref{prop:ch43-free-half}, and it is
characteristic-free.  Injectivity is immediate from bijectivity of
\(\nu^{\An}\).

\emph{The regularization.}  \eqref{eq:bj-char-zero-regularization} is
\Cref{thm:ch34-fs-regularization} applied on \(Y\), in the form recorded as
\Cref{def:bj-psh}(3) followed by GF refinement; it produces a decreasing net
of GF functions with classes tending to \(\nu^*\theta\).  Nothing about
\(X\) is used, and no envelope statement enters
(\Cref{warn:ch44-firewall}).

\emph{Class promotion.}  \Cref{lem:bj-normalization-class-promotion} replaces
each \(M_i\) by an honest pullback \(\nu^*L_i\) without changing the
functions, at the price of passing to a cofinal subnet and of moving the
class by \(o(1)\); the mechanism is the addition of an ample buffer whose
complete linear system contributes the zero function
(\Cref{lem:ch43-zero-fs}).  That lemma is proved in every characteristic and
is available here.

\emph{Assembly.}  Each \(\psi_i\), now a member of
\(\cH^{\mathrm{gf}}_{\mathbf Q}(Y,\nu^*L_i)\), is by
\Cref{prop:bj-char-zero-stage-reflection} a member of
\(\PSH_X\bigl(c_1(L_i)\bigr)\); it is the \emph{same} function of the same
points, by \Cref{lem:bj-universal-analytic-homeomorphism}.  Hence
\((\psi_i)_i\) is a decreasing net in \(\PSH_X\), with classes
\(c_1(L_i)\to\theta\), converging pointwise to \(\psi\).  Its limit is not
identically \(-\infty\): the hypothesis \(\psi\in\PSH_Y(\nu^*\theta)\)
includes generic finiteness, \(\psi\not\equiv-\infty\) on \(Y^{\An}\) by
\Cref{def:bj-psh}(2), and \(Y\) is irreducible; since \(X\) is irreducible
too and \(X^{\An}=Y^{\An}\) as sets by
\Cref{lem:bj-universal-analytic-homeomorphism}, the same function is
generically finite as a function on \(X^{\An}\).  Concretely, the point
\(v_\triv\) is the same point of the two analytifications by
\Cref{lem:ch43-identification-dictionary}\,\cref{itm:ch43-id:trivial}, and
each \(\psi_i\) is finite there --- the presentations at hand have a
generator of \(v\)-value zero at \(v_\triv\), as computed in
\Cref{sec:ch44-decreasing} --- so the decreasing limit is bounded below by
no constant a priori but is finite wherever \(\psi\) is, which by generic
finiteness is somewhere.  Now
\Cref{thm:bj-basic-closure}\,(5) gives \(\psi\in\PSH_X(\theta)\), and
\(\nu^*\psi=\psi\).  Surjectivity follows.
\end{proof}

\begin{remark}[Which hypotheses were used where]
\label{rem:ch44-reflection-hypotheses}
Of the standing hypotheses on \(\nu:Y\to X\), the proof of
\Cref{thm:bj-char-zero-universal-reflection} used: \emph{finite} (to make
\(\nu^{-1}U\) affine, and to make \(\nu^*\) preserve ampleness, in
\Cref{lem:ch43-ample-pullback} and
\Cref{lem:ch44-uniform-cutoffs}\,\cref{itm:ch44-cut:R}); \emph{universal
homeomorphism} (for the analytic identification and for subintegrality);
\emph{projective} (for centeredness of every semivaluation, in
\Cref{prop:ch44-global-sandwich}, and for the existence of the ample \(A_0\)
and \(A_1\)); \emph{integral} (so that \(A_U\to B_U\) is injective);
\emph{\(\operatorname{char}k=0\)} (three times, as listed in
\Cref{sec:ch44-hypotheses}).  It did \emph{not} use normality of \(Y\), and
it did \emph{not} use birationality of \(\nu\).  We record this as
\Cref{prop:ch44-no-birational} below, in parallel with
\Cref{prop:ch43-frobenius-no-birational}; the hypotheses are carried in
\eqref{eq:bj-char-zero-universal-map} because the application is to a
normalization.
\end{remark}

\begin{proposition}
\label{prop:ch44-no-birational}
Let \(k\) be algebraically closed of characteristic zero and trivially
valued, and let \(\nu:Y\to X\) be \emph{any} finite universal homeomorphism
of integral projective \(k\)-varieties.  Then for every
\(\theta\in N^1(X)_{\mathbf R}\), pullback is a bijection
\[
 \nu^*:\PSH_X(\theta)\xrightarrow{\ \sim\ }\PSH_Y(\nu^*\theta).
\]
\end{proposition}

\begin{proof}
Inspect the proofs.  \Cref{lem:ch44-char-zero-subintegral} used only that
\(\nu\) is a finite universal homeomorphism between integral schemes and
that \(\operatorname{char}k=0\); dominance, which gives injectivity of
\(A_U\to B_U\), follows from surjectivity of \(\nu\) and integrality of
\(X\) and \(Y\), not from birationality.
\Cref{lem:bj-local-additive-sandwich} is a statement about an abstract
finite subintegral extension of \(\mathbf Q\)-algebras.
\Cref{lem:ch44-uniform-cutoffs} used ampleness of \(A_0\) and \(A_1\) on
\(X\), coherence of \(\mathcal Q=\nu_*\cO_Y/\cO_X\) --- which requires
\(\cO_X\hookrightarrow\nu_*\cO_Y\), again dominance --- and finiteness.
\Cref{prop:ch44-global-sandwich} used projectivity for centeredness.
\Cref{prop:bj-char-zero-stage-reflection} and
\Cref{thm:bj-char-zero-universal-reflection} used, in addition,
\Cref{lem:bj-universal-analytic-homeomorphism} and
\Cref{lem:bj-normalization-class-promotion}, both of which are proved in
\Cref{chap:bj-universal-homeomorphisms} for a finite universal homeomorphism
without any birationality or normality hypothesis.  No other property of
\(\nu\) was invoked.
\end{proof}

\begin{remark}[Why the extra generality is recorded but not exploited]
\label{rem:ch44-generality-not-exploited}
\Cref{prop:ch44-no-birational} is stated for the same reason as its
positive-characteristic counterpart
\Cref{prop:ch43-frobenius-no-birational}: it separates the hypothesis that
does the work (finite universal homeomorphism) from the hypothesis that
describes the application (birational, with normal source).  In
characteristic zero a finite universal homeomorphism of \emph{normal}
varieties is automatically an isomorphism --- a normal variety is seminormal
in its normalization, hence in any subintegral overring, so
\Cref{prop:ch44-seminormal-unibranch-normal}\,(1) applies --- so the extra
generality of \Cref{prop:ch44-no-birational} is genuinely about nonnormal
targets, and is exactly the generality needed if one later wants to compose
several such maps.  We record the composability as
\Cref{cor:ch44-reflection-tower}.
\end{remark}

\begin{corollary}
\label{cor:ch44-reflection-tower}
Let \(\operatorname{char}k=0\) and let
\(Z\xrightarrow{\ \mu\ }Y\xrightarrow{\ \nu\ }X\) be finite universal
homeomorphisms of integral projective \(k\)-varieties.  Then \(\nu\mu\) is
one as well, and for every \(\theta\in N^1(X)_{\mathbf R}\) the three
pullback maps
\[
 \nu^*,\qquad\mu^*,\qquad(\nu\mu)^*=\mu^*\nu^*
\]
are bijections, compatibly.  In particular the isomorphism class of
\(\PSH_X(\theta)\) as an ordered set depends only on \(X^{\An}\) and
\(\theta\), not on which member of a tower of universal homeomorphisms one
computes it on.
\end{corollary}

\begin{proof}
Composites of finite universal homeomorphisms are finite universal
homeomorphisms (\Cref{lem:ch43-uh-stability}\,\cref{itm:ch43-stab:comp}, and
finiteness is stable under composition).  Apply
\Cref{prop:ch44-no-birational} three times; compatibility is the equality
\((\nu\mu)^*=\mu^*\nu^*\) of pullback of functions, which holds because
\((\nu\mu)^{\An}=\nu^{\An}\circ\mu^{\An}\) as maps of sets.
\end{proof}

\begin{corollary}
\label{cor:ch44-reflection-structure}
Under the hypotheses of \Cref{thm:bj-char-zero-universal-reflection}, the
bijection \eqref{eq:bj-char-zero-universal-reflection} preserves and
reflects the pointwise order, commutes with the addition of real constants,
with finite maxima, and with decreasing limits that are not identically
\(-\infty\), and satisfies
\(\nu^*\bigl(\PSH_X(\theta)\cap\{\phi\le c\}\bigr)
 =\PSH_Y(\nu^*\theta)\cap\{\psi\le c\}\) for every \(c\in\mathbf R\).
\end{corollary}

\begin{proof}
All the operations named are pointwise operations on functions on the single
space \(X^{\An}=Y^{\An}\), and \(\nu^*\) is the identity on functions under
that identification; the only content is that both sides of each identity
consist of psh functions, which is
\Cref{thm:bj-basic-closure}\,(1),~(2),~(4) upstairs and downstairs together
with the bijection.  The last assertion is the statement that boundedness
above by a constant is a pointwise condition.
\end{proof}

\begin{warning}{Boundedness, nefness and the class}
\label{warn:ch44-nef-not-transported}
\Cref{cor:ch44-reflection-structure} transports bounded members to bounded
members, and by \Cref{lem:ch34-psh-nonempty} the existence of a bounded
member of \(\PSH_X(\theta)\) forces \(\theta\in\Nef(X)\).  It does
\emph{not} follow, and is not asserted, that \(\theta\) is nef whenever
\(\PSH_X(\theta)\ne\emptyset\): the correct threshold there is
pseudoeffectivity.  The distinction matters below because none of the
classes \(a_jL+\frac{a_j\rho}{n_j+C}H\) occurring in
\eqref{eq:bj-char-zero-approximating-classes} is asserted to be nef, and
none needs to be: \Cref{def:bj-psh} imposes no positivity on the classes of
a regularizing net.
\end{warning}

\begin{corollary}\label{cor:bj-universal-reflection-all-characteristics}
Let \(\nu:Y\to X\) be a finite birational universal homeomorphism of integral
projective varieties over an algebraically closed trivially valued field,
with \(Y\) normal.
Then, in every characteristic and for every real class \(\theta\),
\[
 \nu^*:\PSH_X(\theta)
       \xrightarrow{\ \sim\ }
       \PSH_Y(\nu^*\theta).
\tag{44.32}\label{eq:bj-universal-reflection-all-characteristics}
\]
\end{corollary}

\begin{proof}
Use \Cref{thm:bj-positive-universal-reflection} in positive characteristic
and \Cref{thm:bj-char-zero-universal-reflection} in characteristic zero.

The two cases are exhaustive: \(k\) is a field, so its characteristic is
either \(0\) or a prime \(p\).  They are also disjoint, and no statement is
used in both directions: the positive-characteristic proof rests on
\Cref{cor:ch43-frobenius-all-affines}, which is false in characteristic zero
(\Cref{warn:ch44-no-gf-bijection}); the present proof rests on
\Cref{thm:appA-roberts-singh}, whose hypothesis is that \(A\) be a
\(\mathbf Q\)-algebra and which is not available in characteristic \(p\)
(\Cref{sec:ch44-failure-charp}).  What the two arguments share is exactly
the characteristic-free material of
\Cref{chap:bj-universal-homeomorphisms}: the analytic homeomorphism, the
identification dictionary, and the class-promotion lemma.
\end{proof}

\section{Completion of the unibranch normalization step}
\label{sec:ch44-unibranch}

The application is to normalizations.  We record it, in the generality in
which \Cref{chap:bj-envelope-consequences} will consume it, and we then say
precisely what remains to be done there --- without doing it, and without
using it.

\begin{corollary}
\label{cor:ch44-unibranch-normalization-reflection}
Let \(k\) be an algebraically closed trivially valued field of arbitrary
characteristic and let \(X\) be an integral unibranch projective
\(k\)-variety, with normalization \(\nu:X^\nu\to X\).  Then:
\begin{enumerate}
\item\label{itm:ch44-unib:uh}
      \(\nu\) is a finite birational universal homeomorphism and
      \(\nu^{\An}:(X^\nu)^{\An}\to X^{\An}\) is a homeomorphism;
\item\label{itm:ch44-unib:subint}
      if \(\operatorname{char}k=0\), then for every affine open
      \(U\subseteq X\) the extension
      \(\cO_X(U)\subseteq\cO_{X^\nu}(\nu^{-1}U)\) is finite and
      subintegral, and it is an equality for all \(U\) if and only if
      \(X\) is normal;
\item\label{itm:ch44-unib:refl}
      for every \(\theta\in N^1(X)_{\mathbf R}\), pullback is a bijection
      \[
      \nu^*:\PSH_X(\theta)\xrightarrow{\ \sim\ }
             \PSH_{X^\nu}(\nu^*\theta),
      \]
      preserving the pointwise order, finite maxima, addition of constants,
      generically finite decreasing limits, and boundedness.
\end{enumerate}
\end{corollary}

\begin{proof}
\cref{itm:ch44-unib:uh} is
\Cref{lem:bj-unibranch-normalization-universal} together with
\Cref{lem:bj-universal-analytic-homeomorphism}.
\cref{itm:ch44-unib:subint} is \Cref{lem:ch44-char-zero-subintegral}, with
the last clause from
\Cref{prop:ch44-seminormal-unibranch-normal}\,(3): equality for all \(U\)
says \(\nu\) is an isomorphism, i.e.\ \(X\) is normal.
\cref{itm:ch44-unib:refl} is
\Cref{cor:bj-universal-reflection-all-characteristics} applied to \(\nu\),
whose hypotheses hold by \cref{itm:ch44-unib:uh}, \(X^\nu\) being normal by
construction; the structural clauses are
\Cref{cor:ch44-reflection-structure} in characteristic zero and
\Cref{cor:ch43-reflection-structure} in characteristic \(p\).
\end{proof}

\begin{remark}[What has been achieved, and what has not]
\label{rem:ch44-what-remains}
\Cref{cor:ch44-unibranch-normalization-reflection} is the last of the
descent steps.  In the vocabulary of \Cref{sec:ch37-proof-map} the chain is
\[
 \text{smooth}\ \longrightarrow\ \text{normal}\ \longrightarrow\
 \text{unibranch},
\]
whose first arrow is \Cref{thm:bj-normal-envelope}
(\Cref{chap:bj-alterations}) and whose second arrow is supplied by
\Cref{cor:ch44-unibranch-normalization-reflection}.  What remains is the
assembly: for \(X\) integral unibranch and
\((\phi_\alpha)\subseteq\PSH_X(\theta)\) nonempty and bounded above, one
pulls back to \(X^\nu\), forms the regularized supremum there, uses the
envelope property of the \emph{normal} variety \(X^\nu\), and pushes the
result back through the bijection.  That the regularized supremum commutes
with \(\nu^*\) is \Cref{cor:ch43-sup-star}, proved in
\Cref{chap:bj-universal-homeomorphisms} from the analytic homeomorphism
alone.

That assembly is \emph{not} performed here.  It uses
\Cref{thm:bj-normal-envelope} as an input, and combining it with the present
theorem produces \Cref{thm:bj-final-envelope}, which by
\Cref{warn:ch44-firewall} may not appear in this chapter in any role.  The
combination is carried out in \Cref{chap:bj-envelope-consequences}.  The
reader who wishes to see the shape of the final argument before then will
find it in \Cref{sec:ch37-proof-map}; the reader auditing this chapter for
circularity should note that the shape is displayed there and used nowhere
below.
\end{remark}

\begin{proposition}
\label{prop:ch44-ledger}
The following inputs, and no others, are used in the proof of
\Cref{thm:bj-char-zero-universal-reflection}:
\begin{enumerate}
\item the definitions and elementary calculus of
      \(\cH^{\mathrm{gf}}_\Lambda\) and \(\PSH\)
      (\Cref{def:gf-fs-functions}, \Cref{def:bj-psh},
      \Cref{prop:ch34-fs-operations});
\item the closure properties \Cref{thm:bj-basic-closure}\,(1) and~(5), and
      the existence thresholds \Cref{lem:ch34-psh-nonempty};
\item GF regularization of psh functions on the \emph{normal} projective
      variety \(Y\) (\Cref{thm:ch34-fs-regularization});
\item the analytic homeomorphism and its dictionary
      (\Cref{lem:bj-universal-analytic-homeomorphism},
      \Cref{lem:ch43-identification-dictionary}) and class promotion
      (\Cref{lem:bj-normalization-class-promotion});
\item the Roberts--Singh additive-power theorem
      (\Cref{thm:appA-roberts-singh}) and its generic form
      (\Cref{cor:appA-RS-generic});
\item the square--cube criterion and its base change
      (\Cref{thm:appA-square-cube}; Stacks Tags 0CND, 0EUH, audited in
      \Cref{sec:stacks-universal-homeomorphisms});
\item standard coherent-sheaf facts on a projective variety: Serre's
      theorems on global generation after a twist, coherence of
      \(\nu_*\cO_Y\) and of \(\mathcal Q\), and left exactness of
      \(H^0\).
\end{enumerate}
In particular no envelope property, in any formulation and on any variety,
is used.
\end{proposition}

\begin{proof}
This is an inspection of the proofs of
\Cref{lem:ch44-char-zero-subintegral},
\Cref{lem:ch44-polynomial-base-change},
\Cref{lem:ch44-finite-differences},
\Cref{lem:bj-local-additive-sandwich},
\Cref{lem:ch44-gauss-valuation},
\Cref{lem:ch44-weight-encoding},
\Cref{lem:ch44-uniform-cutoffs},
\Cref{prop:ch44-global-sandwich},
\Cref{prop:bj-char-zero-stage-reflection} and
\Cref{thm:bj-char-zero-universal-reflection}, in that order; each cites only
items on the list.  The audit is carried out again, statement by statement
and with the extended-real values tracked, in \Cref{sec:ch44-audit}.
\end{proof}

\section{The construction, carried to numbers}
\label{sec:ch44-computations}

We now run the machine.  The first example is the cuspidal cubic, where
every constant can be named; the second is a nonnormal unibranch surface,
where the point is that the same constants exist and are still finite.

\begin{example}[The cuspidal cubic in characteristic zero, end to end]
\label{ex:ch44-cusp-end-to-end}
Let \(k\) be algebraically closed with \(\operatorname{char}k=0\), let
\[
 X=\{y^2z=x^3\}\subset\mathbf P^2_k,
 \qquad
 \nu:Y=\mathbf P^1\longrightarrow X,
 \quad
 [\,a:b\,]\longmapsto[\,a^2b:a^3:b^3\,],
\]
so that \(\nu\) is the normalization: it is finite, birational, bijective,
and induces isomorphisms on residue fields, hence is a finite birational
universal homeomorphism by
\Cref{lem:bj-unibranch-normalization-universal}.  Indeed
\((a^3)^2b^3=(a^2b)^3\), and on the affine chart \(z\ne0\), with
\(t:=a/b\), one has \(x/z=t^2\) and \(y/z=t^3\), so the affine coordinate
ring of \(X\) is
\(A=k[t^2,t^3]\subset B=k[t]=\cO_Y\), and this is
\Cref{rem:ch44-cusp-decomposition}.  Take
\[
 L:=\cO_X(1),
 \qquad
 \nu^*L=\cO_{\mathbf P^1}(3),
\]
and consider the upstairs GF function of the presentation
\[
 \psi=\max\Bigl\{\log|s_1|,\ \log|s_2|-\tfrac13\Bigr\}
 \in\cH^{\mathrm{gf}}_{\mathbf Q}(\nu^*L),
 \qquad
 m=1,\quad\lambda_1=0,\ \lambda_2=-\tfrac13,
\]
where \(s_1,s_2\in H^0(\mathbf P^1,\cO(3))\) are chosen so that, in the
frame of \(\nu^*L\) over \(\nu^{-1}\{z\ne0\}\), their coefficients are
\[
 f_1=1+t,
 \qquad
 f_2=2+t .
\]
(Any two sections with these affine coefficients will do; they exist because
\(1+t\) and \(2+t\) have degree \(\le3\).)  The function \(\psi\) is not
visibly a downstairs GF function, since neither \(f_1\) nor \(f_2\) lies in
\(A\), and by \Cref{warn:ch44-no-gf-bijection} no power of either does.

\emph{Step 0: normalization and weights.}  Here \(c=\frac1m\max_j\lambda_j=0\),
so \(\psi_0=\psi\), and the smallest \(q\) with
\(d_j=-q(\lambda_j-\max_\ell\lambda_\ell)\in\mathbf N\) is \(q=3\):
\[
 d_1=0,
 \qquad
 d_2=1,
 \qquad
 M=qmL=3L,
 \qquad
 \rho=\frac1{qm}=\frac13,
 \qquad
 \mathfrak a=\bigl(s_1^{\otimes3},\ u\,s_2^{\otimes3}\bigr).
\]
This is the situation of \Cref{ex:ch44-encoding-numbers}.  In the chart
\(U=\{z\ne0\}=D(z)\), with \(A_0=\cO_X(1)\) and \(\tau=z\), the generators
of the local ideal are
\[
 g_1=(1+t)^3,
 \qquad
 g_2=u\,(2+t)^3 .
\]

\emph{Step 1: the local Roberts--Singh data.}  Apply
\Cref{ex:ch44-cusp-rs} over \(R=k[u,z_1,z_2]\) to
\[
 b=z_1(1+t)^3+z_2u(2+t)^3
 =\underbrace{\bigl(z_1+8uz_2\bigr)}_{\beta_0}
  +\underbrace{\bigl(3z_1+12uz_2\bigr)}_{\beta_1}t+t^2(\cdots),
\]
where \(\beta_0=b|_{t=0}\) and \(\beta_1=(\partial b/\partial t)|_{t=0}\)
are read off from \((1+t)^3=1+3t+3t^2+t^3\) and
\((2+t)^3=8+12t+6t^2+t^3\).  Hence
\[
 p=1,
 \qquad
 N=0,
 \qquad
 c_1=-\beta_1t=-\bigl(3z_1+12uz_2\bigr)t,
 \qquad
 a_m=b^m-m\,\beta_1t\,b^{m-1}.
\]
These are the promised explicit additive-power elements: one correction
term, equal to \(-t\) times the \(t\)-derivative of the generic combination
at the cusp.

\emph{Step 2: the constants.}  A single chart does not cover \(X\), so take
the cover \(U_1=D(z)\), \(U_2=D(x)\), \(U_3=D(y)\) by the three coordinate
sections of \(A_0=\cO_X(1)\); on \(U_2\) and \(U_3\) the variety \(X\) is
smooth, so there \(A_\alpha=B_\alpha\), \(p_\alpha=0\), \(N_\alpha=0\), and
\Cref{lem:bj-local-additive-sandwich} holds with \(K_{n,\alpha}\) the
\(n\)-th power itself.  Therefore
\[
 P=\max_\alpha p_\alpha=1,
 \qquad
 C=2P=2,
 \qquad
 N_0=0 .
\]
The correction \(c_1\) has coefficients \(3t,12t\in B_1\); multiplying by
\(\tau^1=z\) already extends them to global sections of
\(\nu^*(M+A_0)\), because \(t\) is the affine coordinate and \(zt\) extends,
so \(R=1\) suffices.  The quotient \(\mathcal Q=\nu_*\cO_Y/\cO_X\) is a
skyscraper of length one at the cusp, killed by the maximal ideal, so
\(S=1\) suffices (\Cref{rem:ch44-Q-size}).  Thus
\[
 H_0=(R+S)A_0=2\,\cO_X(1),
\]
and for \(H_1\) it suffices that \(H_1+jM=H_1+3j\,\cO_X(1)\) be globally
generated for \(j=0,1,2\), which holds with \(H_1=\cO_X\), all the bundles
\(3j\cO_X(1)\) being globally generated.  Hence
\[
 \begin{aligned}
  H&=H_0+H_1=2\,\cO_X(1),\\
  D_n&=(n+2)M+H
      =\bigl(3(n+2)+2\bigr)\cO_X(1)
      =(3n+8)\,\cO_X(1).
 \end{aligned}
\]

\emph{Step 3: the approximants and their classes.}  For \(n\ge0\),
\[
 \begin{aligned}
 r_n&=\frac n{n+2},\\
 F_n&=-\frac{1/3}{n+2}\,W(\mathfrak K_n)\\
 &\in\cH^{\mathrm{gf}}_{\mathbf Q}
   \Bigl(\cO_X(1)+\tfrac{2}{3(n+2)}\cO_X(1)\Bigr)\\
 &=\cH^{\mathrm{gf}}_{\mathbf Q}
   \Bigl(\tfrac{3n+8}{3n+6}\,\cO_X(1)\Bigr).
 \end{aligned}
\]
and \(\psi_0\le F_n\le\frac n{n+2}\psi_0\).  For \(n=4\), say, the class of
\(F_4\) is \(\frac{20}{18}\cO_X(1)=\frac{10}9\cO_X(1)\) and
\(r_4=\frac23\); for \(n=100\) the class is \(\frac{308}{306}\cO_X(1)\) and
\(r_{100}=\frac{50}{51}\).  Monotonizing with \(a_j=1-2^{-j}\) and
\(n_{j+1}=\max\{n_j+1,\lceil2(2^{j+1}-2)\rceil\}=\max\{n_j+1,2^{j+2}-4\}\),
the classes of the decreasing sequence are
\[
 a_j\Bigl(1+\tfrac2{3(n_j+2)}\Bigr)\cO_X(1)
 \longrightarrow\cO_X(1),
\]
for instance \(\tfrac12\cdot\tfrac{10}{9}=\tfrac59\) at \(j=1,n_1=4\), then
\(\tfrac34\cdot\tfrac{40}{36}=\tfrac56\) at \(j=2,n_2=12\), then
\(\tfrac78\cdot\tfrac{88}{84}=\tfrac{11}{12}\) at \(j=3,n_3=28\), and so
on, increasing to \(1\).

\emph{Step 4: the estimate at three test points.}  Let \(v_\triv\) be the
trivial valuation, \(v_0:=\operatorname{ord}_{t=0}\) the divisorial
valuation at the cusp read on \(Y\), and \(v_{-1}:=\operatorname{ord}_{t=-1}\).
Then, in the chart \(U_1\):
\[
 W_{v_\triv}(\mathfrak a)=\min\{0,0+1\}=0,
 \qquad
 W_{v_0}(\mathfrak a)=\min\{3\cdot0,\ 3\cdot0+1\}=0,
\]
\[
 W_{v_{-1}}(\mathfrak a)=\min\{3\cdot1,\ 3\cdot0+1\}=1 .
\]
Correspondingly \(\psi_0(v_\triv)=\psi_0(v_0)=0\) and
\(\psi_0(v_{-1})=-\tfrac13\).  The sandwich predicts
\[
 0\le W_{v}(\nu^{-1}\mathfrak K_n)\le0
 \quad\text{at }v_\triv\text{ and }v_0,
 \qquad
 n\le W_{v_{-1}}(\nu^{-1}\mathfrak K_n)\le n+2,
\]
hence \(F_n(v_\triv)=F_n(v_0)=0\) exactly, and
\[
 -\frac13\ \le\ F_n(v_{-1})\ \le\ -\frac13\cdot\frac n{n+2},
\]
so \(F_4(v_{-1})\in[-\tfrac13,-\tfrac29]\) and
\(F_{100}(v_{-1})\in[-\tfrac13,-\tfrac{50}{153}]\).  The convergence to
\(-\tfrac13\) is visible, and the width of the bracket is exactly the
additive constant \(C=2\) divided by \(n+2\).

\emph{Step 5: what happens at the cusp.}  Note that \(v_0\) is the point at
which \(A\ne B\), and yet the sandwich is exact there: both \(f_1=1+t\) and
\(f_2=2+t\) are units at \(t=0\), so \(W_{v_0}(\mathfrak a)=0\) and the
estimate collapses.  The mechanism is not that the cusp is easy but that
this particular \(\psi\) does not see it.  A presentation that does see it
--- for instance one with \(f_1=t\), whose \(v_0\)-value is \(1\) --- gives
\(W_{v_0}(\mathfrak a)=\min\{3,1\}=1\) and the genuine bracket
\(n\le W_{v_0}(\nu^{-1}\mathfrak K_n)\le n+2\).  In that case the descended
family really does involve the corrected powers: the element
\(a_{n+1}\) contributes, in \(\mathbf z\)-degree \(n+1\) and in the
\(u\)-degree-zero part, the polynomial
\(G_{n+1}=(1+t)^n(1-nt)\) of \Cref{ex:ch44-splitting-fails}, which lies in
\(A\) although \((1+t)^{n+1}\) does not.
\end{example}

\begin{example}[The cone over the cuspidal cubic, end to end]
\label{ex:ch44-surface-end-to-end}
The previous example is a curve, and the analogous surface used in
\Cref{ex:ch43-surface-endtoend} --- the cuspidal cylinder
\(C\times\mathbf P^1\) --- has product structure that trivializes the
globalization.  We therefore take a surface whose nonnormal locus is a
curve and which is not a product, namely the projective cone.

\emph{(a) The two rings.}  Let \(\operatorname{char}k=0\) and put
\[
 \begin{aligned}
  B_\bullet&:=k[a,b]^{(3)}[w]
    =k\bigl[a^3,\,a^2b,\,ab^2,\,b^3\bigr][w],\\
  A_\bullet&:=k\bigl[a^3,\,a^2b,\,b^3\bigr][w]
    \subset B_\bullet,
 \end{aligned}
\]
graded so that \(a^3,a^2b,ab^2,b^3\) and \(w\) all have degree one.  Thus
\(B_\bullet\) is the third Veronese subring of \(k[a,b]\), extended by one
free variable, and \(A_\bullet\) is obtained by discarding the generator
\(ab^2\).  Setting \(x:=a^2b\), \(y:=a^3\), \(z:=b^3\) one has the single
relation \(y^2z=x^3\), so
\[
 A_\bullet\cong k[x,y,z,w]/(y^2z-x^3),
 \qquad
 X:=\Proj A_\bullet=\{y^2z=x^3\}\subset\mathbf P^3
\]
is the projective cone over the cuspidal cubic of
\Cref{ex:ch44-cusp-end-to-end}, with vertex \([\,0:0:0:1\,]\), and
\[
 Y:=\Proj B_\bullet\subset\mathbf P^4
\]
is the projective cone over the rational normal cubic curve.  Both are
integral projective surfaces; \(Y\) is normal, being the \(\Proj\) of the
semigroup ring of the saturated semigroup
\(\{(i,j):3\mid i+j\}\) extended by a free variable, and \(\nu:Y\to X\) is
finite, since \(B_\bullet=A_\bullet+A_\bullet\,ab^2\) is a finite
\(A_\bullet\)-module, and birational.

\emph{(b) Subintegrality.}  Set \(\zeta:=ab^2\).  Then
\[
 \zeta^2=a^2b^4=(a^2b)(b^3)=xz\in A_\bullet,
 \qquad
 \zeta^3=a^3b^6=(a^3)(b^3)^2=yz^2\in A_\bullet,
\]
and \(B_\bullet=A_\bullet[\zeta]\).  So on every chart \(D_+(f)\) with
\(f\in A_1\) the degree-zero extension \(A_{(f)}\subset B_{(f)}\) is
generated by \(\zeta/f\), whose square and cube lie in \(A_{(f)}\); by
\Cref{thm:appA-square-cube} it is subintegral.  Hence \(\nu\) is a finite
birational universal homeomorphism, \(X\) is unibranch and, by
\Cref{prop:ch44-seminormal-unibranch-normal}\,(2), not seminormal.  The
hypotheses \eqref{eq:bj-char-zero-universal-map} hold.

\emph{(c) The gap structure and the length \(P\).}  Grade \(B_\bullet\) by
the exponent of \(a\), writing \(f=\sum_if^{(i)}\) with \(f^{(i)}\) the part
of \(a\)-degree \(i\).  The monomials \(a^ib^j\) with \(3\mid i+j\) that lie
\emph{outside} \(A_\bullet\) are exactly those with \(i=1\): if \(i=0\) then
\(3\mid j\) and \(a^ib^j\in k[b^3]\); if \(i\ge2\) one checks directly, by
writing \(i=3\mu\), \(i=3\mu+1\ (\mu\ge1)\) or \(i=3\mu+2\), that
\[
 (i,j)\in\mathbf N(3,0)+\mathbf N(2,1)+\mathbf N(0,3);
\]
and \((1,j)\) is never of this form, since \(3\mu+2\nu=1\) is insoluble in
\(\mathbf N\).  Therefore
\[
 A_\bullet[w]\text{-membership of }f\in B_\bullet
 \iff f^{(1)}=0 .
\]
Since \(B^{(0)}\subseteq A_\bullet\), \(B^{(\ge2)}\subseteq A_\bullet\), and
the \(a\)-degree is additive, the \(a\)-degree-one part of a product is
computed by choosing the degree-one factor exactly once:
\[
 \bigl(f^m\bigr)^{(1)}=m\,\bigl(f^{(0)}\bigr)^{m-1}f^{(1)},
 \qquad
 \bigl(f^{(1)}f^{m-1}\bigr)^{(1)}=f^{(1)}\bigl(f^{(0)}\bigr)^{m-1}.
\]
Subtracting, \(f^m-m\,f^{(1)}f^{m-1}\) has vanishing \(a\)-degree-one part,
hence lies in \(A_\bullet\).  So the Roberts--Singh data are
\[
 p=1,\qquad N=0,\qquad c_1=-f^{(1)},
\]
exactly as for the cusp (\Cref{ex:ch44-cusp-rs}), with the ``coefficient of
\(t\)'' replaced by the ``\(a\)-degree-one layer''.  On the charts
\(D_+(x)\) and \(D_+(y)\) the surface \(X\) is smooth --- dividing
\(y^2z=x^3\) by \(y^3\) gives \(z/y=(x/y)^3\), and by \(x^3\) gives
\((y/x)^2(z/x)=1\) --- so there \(p_\alpha=0\).  Therefore
\[
 P=1,\qquad C=2P=2,\qquad N_0=0 .
\]

\emph{(d) The GF function.}  Put \(L:=\cO_X(1)\), so \(\nu^*L=\cO_Y(1)\),
and take \(m=1\), \(r=2\), \(\lambda_1=0\), \(\lambda_2=-\tfrac13\),
\[
 \psi:=\max\Bigl\{\log|s_1|,\ \log|s_2|-\tfrac13\Bigr\},
 \qquad
 s_1:=\zeta+z=ab^2+b^3,
 \qquad
 s_2:=w .
\]
The section \(s_1\in H^0(Y,\cO_Y(1))\) is \emph{not} a pullback: the
degree-one part of \(B_\bullet\) is spanned by \(a^3,a^2b,ab^2,b^3,w\) and
that of \(A_\bullet\) by \(a^3,a^2b,b^3,w\), so \(H^0(Y,\nu^*L)\) exceeds
\(H^0(X,L)\) by exactly the line \(k\zeta\), and \(s_1\) meets it.  This is
the reason there is something to prove.  As in
\Cref{ex:ch44-encoding-numbers}, \(c=0\), \(q=3\), \(d_1=0\), \(d_2=1\),
\[
 \begin{gathered}
 M=3L,\qquad \rho=\tfrac13,\\
 \mathfrak a
 =\bigl(s_1^{\otimes3},\ u\,s_2^{\otimes3}\bigr).
 \end{gathered}
\]

\emph{(e) The twists, computed.}  Take \(A_0:=\cO_X(1)\) and use the cover
\[
 D_+(x),\qquad D_+(y),\qquad D_+(z),\qquad D_+(w).
\]
The generic element on the chart
\(D_+(w)\), in the frame \(e_w=w^{\otimes3}\), is
\[
 b_w=z_1\frac{(ab^2+b^3)^3}{w^3}+z_2u ,
 \qquad
 (ab^2+b^3)^3=b^9+3ab^8+3a^2b^7+a^3b^6,
\]
whose \(a\)-degree-one layer is \(3ab^8/w^3\); hence
\(c_{w1}=-3z_1\,ab^8/w^3\) and
\[
 c_{w1}\,e_w^{\otimes1}=-3z_1\,ab^8,
\]
which is \emph{already} a global section of \(\nu^*M=\cO_Y(3)\), since
\(ab^8\) has \(a\)-\(b\)-degree \(9=3\cdot3\).  On the chart \(D_+(z)\),
with \(t:=a/b\) and \(W:=w/z\), the frame is \(e_z=z^{\otimes3}=b^9\), the
local rings are \(A_z=k[t^2,t^3][W]\subset B_z=k[t][W]\), the generic
element is \(z_1(1+t)^3+z_2uW^3\), and \Cref{ex:ch44-cusp-rs} gives
\(c_{z1}=-3z_1t\), so
\(c_{z1}e_z^{\otimes1}=-3z_1t\,b^9=-3z_1ab^8\): the \emph{same} global
section, as it must be by the covariance
\eqref{eq:bj-char-zero-frame-law}.  Hence
\[
 R=0 .
\]
For \(S\), the sheaf \(\mathcal Q=\nu_*\cO_Y/\cO_X\) is generated by the
class of \(\zeta\), and
\[
 x\zeta=a^3b^3\in A_\bullet,
 \quad
 y\zeta=a^4b^2\in A_\bullet,
 \quad
 z\zeta=ab^5\notin A_\bullet,
 \quad
 w\zeta\notin A_\bullet,
\]
so \(\operatorname{Ann}\mathcal Q\supseteq(x,y)\) and \(\mathcal Q\) is
supported on the line \(\ell=\{x=y=0\}\), the cone over the cusp.  On
\(D_+(w)\cap\ell\) one finds \(\mathcal Q\cong\zeta\cdot k[z]\), free of
rank one over \(\cO_\ell\); so \(\mathcal Q\) is an invertible sheaf on
\(\ell\cong\mathbf P^1\), and has no nonzero section supported at a point.
Consequently
\[
 \Gamma_{V(x)}\mathcal Q=\Gamma_{V(y)}\mathcal Q=\mathcal Q,
 \quad\text{killed by }x,y;
 \qquad
 \Gamma_{V(z)}\mathcal Q=\Gamma_{V(w)}\mathcal Q=0,
\]
whence \(S=1\).  Finally \(H_1=\cO_X\) works, since \(jM=\cO_X(3j)\) is
globally generated for \(0\le j\le2\).  Therefore
\[
 H=(R+S)A_0+H_1=\cO_X(1),
 \qquad
 D_n=(n+2)\cdot3L+\cO_X(1)=\cO_X(3n+7).
\]
Every one of \(R=0\), \(S=1\), \(H_1=\cO_X\), \(C=2\) was determined before
\(n\) entered: this is point~\textbf{(C2)} of \Cref{sec:ch44-care}, verified
by hand.

\emph{(f) The approximants.}  For \(n\ge0\),
\[
 r_n=\frac n{n+2},
 \qquad
 F_n\in\cH^{\mathrm{gf}}_{\mathbf Q}
   \Bigl(\tfrac{3n+7}{3n+6}\,\cO_X(1)\Bigr),
 \qquad
 \psi\le F_n\le\frac n{n+2}\,\psi .
\]
Test the estimate at three points.  Let \(v_\triv\) be the trivial
valuation; let \(v_\ell:=\operatorname{ord}_{\tilde\ell}\), the divisorial
valuation of the curve \(\tilde\ell=\{a=0\}\subset Y\) lying over the
nonnormal line \(\ell=\{x=y=0\}\subset X\); and let
\(v_D:=\operatorname{ord}_D\), where \(D\subset Y\) is the ruling of the
cone over the point \([\,1:-1\,]\in\mathbf P^1\).  Then
\(v_\ell(s_1)=\operatorname{ord}_a(ab^2+b^3)=0\), because the summand
\(b^3\) does not vanish along \(\{a=0\}\), and \(v_\ell(s_2)=0\), so
\(W_{v_\ell}(\mathfrak a)=\min\{0,0+1\}=0\) and
\(F_n(v_\ell)=\psi(v_\ell)=0\) exactly.  Since
\(s_1=ab^2+b^3=b^2(a+b)\), the zero divisor of \(s_1\) on \(Y\) is
\(2R_{[0:1]}+D\), where \(R_P\) is the ruling over \(P\); hence
\(v_D(s_1)=1\),
\(v_D(s_2)=0\), so
\[
 W_{v_D}(\mathfrak a)=\min\{3\cdot1,\ 3\cdot0+1\}=1,
 \qquad
 \psi(v_D)=-\tfrac13,
 \qquad
 -\tfrac13\le F_n(v_D)\le-\tfrac13\cdot\tfrac n{n+2},
\]
so \(F_2(v_D)\in[-\tfrac13,-\tfrac16]\) and
\(F_{50}(v_D)\in[-\tfrac13,-\tfrac{25}{78}]\).  At \(v_\triv\),
\(W(\mathfrak a)=0\) and \(F_n(v_\triv)=0\) for every \(n\), which is the
check that excludes the identically \(-\infty\) alternative.

\emph{(g) Conclusion.}  \Cref{prop:bj-char-zero-stage-reflection} gives
\(\psi\in\PSH_X(c_1(L))\), and
\Cref{thm:bj-char-zero-universal-reflection} gives, for every
\(\theta\in N^1(X)_{\mathbf R}\), a bijection
\(\PSH_X(\theta)\cong\PSH_Y(\nu^*\theta)\).  The right-hand side is a psh
theory on a \emph{normal} projective surface, to which
\Cref{thm:bj-normal-envelope} applies; the left-hand side is a psh theory on
a surface that is neither normal nor seminormal and on which no direct
argument is available.  Transporting the former to the latter is the
assembly of \Cref{chap:bj-envelope-consequences}, and is not performed here
(\Cref{rem:ch44-what-remains}).
\end{example}

\section{Extended-real and circularity audit}
\label{sec:ch44-audit}

\subsection{Where extended-real values occur}
\label{sec:ch44-extended-real}

The value \(-\infty\) occurs throughout, and the value \(+\infty\) occurs on
the valuative side.  We list every occurrence and say what is done with it,
in the manner of \Cref{sec:ch43-extended-real}.

\begin{enumerate}
\item\label{itm:ch44-er:homeo}
      \emph{The identification.}  \(\nu^{\An}\) is a homeomorphism, so
      \(\psi\) and \(\nu^*\psi\) are literally the same
      \(\Rbar\)-valued function of the same points.  Nothing is truncated,
      and points with \(\psi=-\infty\) are not discarded
      (\Cref{lem:bj-universal-analytic-homeomorphism}).
\item\label{itm:ch44-er:gauss}
      \emph{The Gauss weight.}  \(W_v\) takes values in \([0,+\infty]\);
      the additions in \Cref{lem:ch44-gauss-valuation} are of nonnegative
      extended reals and are unambiguous.  No difference of infinities is
      formed (\Cref{rem:ch44-gauss-infinity}).
\item\label{itm:ch44-er:encoding}
      \emph{The encoding identity.}  \(\psi_0=-\rho W(\mathfrak a)\) is an
      identity in \([-\infty,0]\); the only subtraction performed in its
      proof is of the finite rational constant \(\max_\ell\lambda_\ell\)
      (\Cref{lem:ch44-weight-encoding}).
\item\label{itm:ch44-er:sandwich}
      \emph{The two-sided estimate.}  In
      \Cref{prop:ch44-global-sandwich} all three quantities are
      simultaneously \(+\infty\) exactly on the common zero locus; the
      products \(n\cdot(+\infty)\) and \((n+C)\cdot(+\infty)\) occur with
      positive integer factors and are unambiguous.  The convention
      \(0\cdot(+\infty)=0\) is never used.
\item\label{itm:ch44-er:monotone}
      \emph{The monotonization.}  The chain
      \eqref{eq:bj-char-zero-monotone-chain} multiplies inequalities between
      elements of \([-\infty,0]\) by positive rationals \(a_j\) and
      \(r_{n_j}\); all four terms may be \(-\infty\) simultaneously, and the
      chain then reads \(-\infty\le-\infty\le-\infty\le-\infty\), which is
      true and is all that is needed.
\item\label{itm:ch44-er:psh2}
      \emph{The closure step.}  \Cref{thm:bj-basic-closure}\,(5) requires
      the limit to be generically finite.  This is checked at the single
      point \(v_\triv\), where the presence of a generator with \(d_j=0\)
      forces the value \(0\) (\Cref{sec:ch44-decreasing}); the alternative
      \(\psi\equiv-\infty\) is thereby excluded and never has to be
      discussed again.
\end{enumerate}

\begin{warning}{Forbidden manoeuvres}
\label{warn:ch44-forbidden}
Three operations would be illegitimate and are never performed.
\emph{Cancelling infinities}: from \(\psi_0\le F_n\le r_n\psi_0\) one may
\emph{not} deduce \(F_n-\psi_0\le(r_n-1)\psi_0\) pointwise where
\(\psi_0=-\infty\), since the left-hand side would be
\((-\infty)-(-\infty)\); all the estimates of this chapter are formulated
monotonically, without subtraction, exactly for this reason.
\emph{Restricting to finite loci}: one may not prove the sandwich on
\(\{\psi_0>-\infty\}\) and extend by continuity, since \(\psi_0\) is only
upper semicontinuous and the common zero locus can be large.
\emph{Passing to suprema}: at no point is a family of psh functions
replaced by its (regularized) supremum; see
\Cref{warn:ch44-firewall} and \Cref{rem:ch44-firewall-not-vacuous}.
\end{warning}

\subsection{Checkpoints}
\label{sec:ch44-checkpoints}

The following seven checkpoints are the ones an auditor should verify; each
names the statement in which it is discharged.

\begin{enumerate}
\item The extension \(A_U\subset B_U\) is subintegral, and not merely
      radicial --- \Cref{lem:ch44-char-zero-subintegral}, using
      \(\operatorname{char}k=0\).
\item The additive-power identity is applied \emph{after} polynomial base
      change and to the \emph{full} generic combination ---
      \Cref{lem:ch44-polynomial-base-change},
      \Cref{lem:bj-local-additive-sandwich}, Steps~1--2;
      failure mode \Cref{ex:ch44-splitting-fails}.
\item The Roberts--Singh data \(p,N,c_i\) do not depend on \(m\), hence not
      on \(n\) --- \Cref{thm:appA-roberts-singh},
      \Cref{rem:ch44-exponent-zero}.
\item The multinomial coefficients are inverted --- Step~5 of
      \Cref{lem:bj-local-additive-sandwich}, using
      \(\mathbf Q\subseteq A\).
\item The three twists \(R,S,H_1\) are chosen before \(n\) ---
      \Cref{lem:ch44-uniform-cutoffs}; failure mode
      \Cref{warn:ch44-moving-twist}.
\item The upper bound of the sandwich is verified at the \emph{center} of
      the given semivaluation, with a chart in which \(\tau_\alpha\) and
      some aligner are units --- \Cref{prop:ch44-global-sandwich}, and
      \Cref{rem:ch44-two-halves}.
\item The final closure appeal is to decreasing nets with varying classes,
      never to suprema --- \Cref{thm:bj-basic-closure}\,(5),
      \Cref{rem:ch44-firewall-not-vacuous},
      \Cref{prop:ch44-ledger}.
\end{enumerate}

\section{Failure boundaries}
\label{sec:ch44-failure}

\subsection{Dropping characteristic zero}
\label{sec:ch44-failure-charp}

\Cref{thm:appA-roberts-singh} requires \(A\) to contain \(\mathbf Q\), and
the requirement is used twice: implicitly, in the proof of the theorem in
\cite{RobertsSingh}, whose exponent-zero criterion is obtained by inverting
integers; and explicitly, in Step~5 of
\Cref{lem:bj-local-additive-sandwich}, where the multinomial coefficient
\(\binom{n+2p}\beta\) must be a unit.  The second use is easy to see
failing.  Over a field of characteristic \(p\), take \(s=2\) and
\(n+2p=p\); then \(\binom p1=p=0\), and the identity
\(\binom p\beta g^\beta\in K_nB[u]\) is vacuous for
\(\beta=(1,p-1)\) and \(\beta=(p-1,1)\).  So even granting the additive-power
identity, the extraction of monomials would fail.  Positive characteristic
is therefore handled by Frobenius
(\Cref{chap:bj-universal-homeomorphisms}), not by reducing the
Roberts--Singh argument modulo \(p\).  Conversely, the
characteristic-\(p\) argument fails in characteristic zero at
\Cref{lem:ch43-frobenius-subalgebra}, since \(\{b:b^q\in A\}\) is not closed
under addition (\Cref{warn:ch44-no-gf-bijection}).  The two proofs are
genuinely disjoint, and there is at present no uniform argument.

\subsection{Dropping subintegrality while keeping finiteness}
\label{sec:ch44-failure-subintegral}

If \(\nu\) is finite and birational but not a universal homeomorphism ---
for instance the normalization of a nodal curve --- then
\Cref{def:bj-subintegral-extension} fails, the Roberts--Singh theorem does
not apply, \(\nu^{\An}\) is not injective, and the conclusion itself is
false for ample classes.  This is \Cref{ex:ch44-node-seminormal}, and the
necessity of unibranchness is \Cref{thm:ep-implies-unibranch}.  The correct
statement for a general finite morphism of \emph{normal} varieties is the
averaging theorem of \Cref{chap:bj-finite-descent}, which divides by a fibre
degree; here the fibre degree is one and no averaging is possible or needed
(\Cref{rem:ch43-no-averaging}).

\subsection{Dropping projectivity}
\label{sec:ch44-failure-projective}

Projectivity of \(X\) is used twice.  First, every semivaluation of
\(X^{\An}\) has a center, which is what allows the upper bound of
\Cref{prop:ch44-global-sandwich} to be verified chartwise; on an affine or
merely proper \(X\) this must be re-examined.  Second, ample bundles
\(A_0,A_1\) exist, which is what makes the three twists of
\Cref{lem:ch44-uniform-cutoffs} available.  On a proper nonprojective \(X\)
one may still have enough line bundles for the argument, but the statement
as given is not claimed.

\subsection{Dropping uniformity of \(C\)}
\label{sec:ch44-failure-uniform}

Suppose the sandwich were only available with an exponent \(\kappa n\),
\(\kappa>1\), in place of \(n+C\).  Then \(r_n\to\kappa^{-1}<1\) and
\eqref{eq:bj-char-zero-function-sandwich} would give
\(\psi_0\le F_n\le\kappa^{-1}\psi_0\) in the limit, so the constructed
decreasing family would converge to something between \(\psi_0\) and
\(\kappa^{-1}\psi_0\) and \emph{not} to \(\psi_0\).  Nothing else in the
argument would visibly break.  This is the failure mode of
\Cref{ex:ch44-splitting-fails}\,(3), and it is why
\Cref{rem:ch44-sandwich-shape} insists on the word ``additive''.

\subsection{What is not claimed}
\label{sec:ch44-failure-notclaimed}

The chapter does not claim: that \(\nu^*\) is surjective on GF functions
(\Cref{warn:ch44-no-gf-bijection}, \Cref{rem:ch44-approximants-are-gf});
that the constants \(C\), \(H\) can be bounded in terms of discrete
invariants of \(X\) alone, independently of the presentation
\eqref{eq:bj-char-zero-gf} --- they cannot, since \(p\) depends on the
generic element (\Cref{warn:ch44-p-not-uniform}); that the Roberts--Singh
data are unique or canonical; that \(\mathcal Q\) is annihilated by the
conductor in any uniform way beyond
\Cref{lem:ch44-uniform-cutoffs}\,\cref{itm:ch44-cut:S}; or anything at all
about envelopes (\Cref{warn:ch44-firewall}).

\section{Chapter dictionary}
\label{sec:ch44-dictionary}

\begin{center}
\begin{tabular}{p{.29\textwidth}p{.31\textwidth}p{.29\textwidth}}
\toprule
geometric datum & algebraic datum & hypothesis\\
\midrule
\(\nu\) a finite universal homeomorphism, \(\operatorname{char}k=0\) &
\(A_U\subset B_U\) subintegral: bijective spectra, trivial residue
extensions &
\(X,Y\) integral, \(\nu\) finite surjective
(\Cref{def:bj-subintegral-extension},
\Cref{lem:ch44-char-zero-subintegral})\\
\(X\) unibranch, nonnormal &
\(A\subset B\) subintegral with \(A\) not seminormal in \(B\);
\(\mathcal Q=\nu_*\cO_Y/\cO_X\ne0\) &
\(k=\bar k\), \(\operatorname{char}k=0\)
(\Cref{prop:ch44-seminormal-unibranch-normal})\\
weights of a GF function &
powers of a formal variable \(u\) with \(W_v(u)=1\); weight \(-d\)
\(\leftrightarrow\) factor \(u^d\) &
weights rational, cleared by \(q\)
(\Cref{lem:ch44-gauss-valuation},
\Cref{lem:ch44-weight-encoding})\\
the GF function itself &
\(\psi_0=-\rho\,W(\mathfrak a)\), a Gauss weight of a homogeneous ideal &
\(\psi_0\le0\) after normalizing by \(c\)
(\eqref{eq:bj-char-zero-weighted-function})\\
generic linear combination &
\(b=\sum_jz_jg_j\in B[u,\mathbf z]\), \(u\)-inhomogeneous by design &
subintegrality survives polynomial base change
(\Cref{lem:ch44-polynomial-base-change}; Tags 0CND, 0EUH)\\
corrected powers &
\(a_m=b^m+\sum_i\binom mic_ib^{m-i}\in A[u,\mathbf z]\), \(c_i\)
independent of \(m\) &
\(\mathbf Q\subseteq A\)
(\Cref{thm:appA-roberts-singh}, \eqref{eq:bj-rs-additive-powers})\\
one power replaced by a block &
\(\mathfrak a^{n+2p}\subseteq K_nB[u]\subseteq\mathfrak a^n\),
\(K_n\subseteq A[u]\) &
\(g_j\) \(u\)-homogeneous; multinomials invertible
(\Cref{lem:bj-local-additive-sandwich})\\
globalization with one twist &
\(\tau_\alpha^R\) extends the \(c_{\alpha i}\); \(\tau_\alpha^S\) kills
\(\Gamma_{V(\tau_\alpha)}\mathcal Q\); \(H_1\) aligns degrees &
\(A_0\) ample, \(\mathcal Q\) coherent, \(C=2P\) finite
(\Cref{lem:ch44-uniform-cutoffs})\\
two-sided valuative bound &
\(nW(\mathfrak a)\le W(\nu^{-1}\mathfrak K_n)\le(n+C)W(\mathfrak a)\) &
every semivaluation centered; \(X\) projective
(\Cref{prop:ch44-global-sandwich})\\
psh reflection &
\(\nu^*:\PSH_X(\theta)\to\PSH_Y(\nu^*\theta)\) bijective, order preserving &
\(\operatorname{char}k=0\); birationality and normality not needed
(\Cref{thm:bj-char-zero-universal-reflection},
\Cref{prop:ch44-no-birational})\\
\bottomrule
\end{tabular}
\end{center}

Four separations should be carried forward.

\emph{Identity versus estimate.}  In characteristic \(p\) the descended
object \emph{is} the given function; here it is only a two-sided bracket
around it, and the given function is recovered as a limit.  The bracket has
additive width \(C\), which is what makes the limit correct; a
multiplicative width would not
(\Cref{rem:ch44-sandwich-shape}, \Cref{sec:ch44-failure-uniform}).

\emph{Exponent versus length.}  The number \(q=p^e\) of
\Cref{chap:bj-universal-homeomorphisms} is an exponent applied once.  The
number \(p\) here is the \emph{length} of a correction, and it is never
applied to anything; what is applied is the block of \(p+1\) consecutive
corrected powers.  The two numbers are not analogues, and the letter \(p\)
here has no relation to a characteristic, which is zero.

\emph{Order of grading extraction.}  Form \(a_m\); then take
\(\mathbf z\)-components; then take \(u\)-components.  Each arrow is
justified by a grading of \(A[u,\mathbf z]\) or of \(A[u]\), and no arrow
may be applied earlier (\Cref{sec:ch44-two-gradings},
\Cref{warn:ch44-no-splitting}, \Cref{ex:ch44-splitting-fails}).

\emph{Local data versus one global twist.}  Every chart has its own
\(p_\alpha,N_\alpha,c_{\alpha i}\); the globalization must not inherit that
dependence in the twist.  \(R\), \(S\), \(H_1\) and \(C\) are all maxima
over the finitely many charts of data attached to fixed generic elements,
and are chosen before \(n\)
(\Cref{lem:ch44-uniform-cutoffs}, \Cref{warn:ch44-moving-twist}).

\section{Source and provenance}
\label{sec:ch44-source-note}

The local sandwich fails as written without the \(\mathbf Q\)-algebra
hypothesis, because the multinomial coefficients extracted from
\(b^{n+2p}\) need not be units.
Positive characteristic is therefore handled by Frobenius, not by reducing
the Roberts--Singh proof modulo \(p\).

The cutoff exponent \(R\), the quotient-killing exponent \(S\), and the
alignment line bundle \(H_1\) are fixed before \(n\) varies.
Allowing any of them to grow linearly with \(n\) would produce an uncontrolled
class error instead of the vanishing term
\(\rho H/(n+C)\).
This uniformity is the geometric heart of the globalization.

The local additive-power lemma and its globalization follow the
characteristic-zero portion of the supplied envelope manuscript.
The elementwise exponent-zero theorem is
\cite[Theorem~4.17, Proposition~4.2, and Corollary~4.3]{RobertsSingh}.
Polynomial base change and the square--cube description are exactly Tags
0CND and 0EUH of \cite{stacks-project}.
GF regularization, pullback, and PSH2 are the BJ inputs from \cite{BJ22}.
No envelope property is used in the construction, so
\Cref{cor:bj-universal-reflection-all-characteristics} is noncircular.

\subsection{Correspondence with the supplied manuscript}
\label{sec:ch44-correspondence}

The manuscript is unpublished and has been made available to the author; its
results are quoted here as coming from a supplied manuscript, and none of
them is presented as a previously published theorem.  In its numbering, the
statements of this chapter correspond as follows.  Its subsection on
subintegral extensions and additive powers is the source of
\Cref{sec:ch44-subintegral} and \Cref{sec:ch44-roberts-singh}, and its
displayed additive-power identity is
\eqref{eq:bj-rs-additive-powers}; its ``additive-power sandwich'' lemma is
\Cref{lem:bj-local-additive-sandwich}; its characteristic-zero proposition
on descent of a Fubini--Study function is
\Cref{prop:bj-char-zero-stage-reflection}; and the characteristic-zero half
of its theorem on descent through the normalization is
\Cref{thm:bj-char-zero-universal-reflection}.  The
characteristic-\(p\) half of that theorem, together with the analytic
homeomorphism and the class-promotion lemma on which both halves rest, is
\Cref{chap:bj-universal-homeomorphisms} and is not duplicated here; the two
halves are assembled in
\Cref{cor:bj-universal-reflection-all-characteristics}.

Several items below are ours rather than the manuscript's, and are flagged
as such.
\begin{enumerate}
\item The finite-difference calculus is isolated as
      \Cref{lem:ch44-finite-differences}, with hypotheses weaker than those
      of the surrounding lemma (no subintegrality, no
      \(\mathbf Q\)-algebra, no finiteness).  The manuscript performs the
      same computation inline.  Separating it makes visible that
      \emph{only} the \(m\)-independence of the \(c_i\) is used, which is
      \Cref{rem:ch44-exponent-zero}.
\item The three uniformities of the globalization are collected into a
      single statement, \Cref{lem:ch44-uniform-cutoffs}, and the
      consequence of losing them is computed in
      \Cref{warn:ch44-moving-twist}.  The manuscript asserts the
      uniformities in the course of its proof.  Since ``one twist,
      independent of \(n\)'' is the point on which the class of the limit
      depends, we have made it a lemma with a proof.
\item The global two-sided estimate is stated as
      \Cref{prop:ch44-global-sandwich} rather than displayed inside a
      proof, so that \Cref{sec:ch44-decreasing} and
      \Cref{sec:ch44-reflection} can cite it.
\item \Cref{prop:ch44-seminormal-unibranch-normal} --- that a seminormal
      unibranch variety in characteristic zero is normal --- and its
      corollary that \(\mathcal Q\ne0\) in the situation of interest, are
      recorded here to show that the hypotheses are not vacuous.  The
      manuscript does not state them.
\item \Cref{prop:ch44-no-birational}, that neither birationality of
      \(\nu\) nor normality of \(Y\) is used, is an inspection of the
      proofs and is not asserted in the manuscript, which carries both as
      standing hypotheses because its application is to a normalization.
      The observation is harmless and is recorded because it clarifies
      which hypothesis does what; the same observation in characteristic
      \(p\) is \Cref{prop:ch43-frobenius-no-birational}.
\item The worked computations \Cref{ex:ch44-cusp-rs},
      \Cref{ex:ch44-two-gaps}, \Cref{ex:ch44-splitting-fails},
      \Cref{ex:ch44-cusp-end-to-end} and
      \Cref{ex:ch44-surface-end-to-end} are ours.  The second shows that
      the Roberts--Singh length can exceed one; the third is the
      counterexample to the prohibited splitting, and its part~(2) --- the
      explicit monomial \(g_1^2g_2^n\) missing from the grouped ideal ---
      is what turns a warning into a proof that the order of operations is
      forced.
\end{enumerate}

\subsection{External attributions}
\label{sec:ch44-external}

The notion of subintegrality arose in Swan's work on seminormality, and the
additive-power characterization used here is due to Roberts and Singh, with
a later refinement by Reid, Roberts and Singh; we cite the first of these by
key as \cite{RobertsSingh} --- the source of
\Cref{thm:appA-roberts-singh} --- and the other two in prose, no key for
them being present in the book's bibliography.  Traverso's seminormalization
and Swan's correction of it are the background for
\Cref{def:bj-seminormal-overring}, and are likewise cited in prose.

The Stacks Project inputs are exactly Tags 0CND and 0EUH, audited with their
exact statements, the hypotheses to be checked, and the first place of
application in \Cref{sec:stacks-universal-homeomorphisms}; no other tag is
used in this chapter, and none is invented.
\Cref{warn:ch44-tags-scope} records what those two tags do \emph{not}
supply.

The psh formalism --- the definition of
\(\cH^{\mathrm{gf}}_{\mathbf Q}\), the closure of \(\PSH\) under decreasing
nets with varying classes, and the regularization of psh functions by GF
functions --- is Boucksom--Jonsson \cite{BJ22}, with the corrections and
additions of \cite{BJ22b}; the book's versions are
\Cref{def:gf-fs-functions}, \Cref{def:bj-psh},
\Cref{thm:bj-basic-closure} and \Cref{thm:ch34-fs-regularization}, proved in
\Cref{chap:bj-fs-pl-psh}.  The necessity of unibranchness for an ample class
is \Cref{thm:ep-implies-unibranch}, from
\Cref{chap:bj-envelope-history}, and is independent of this chapter.

\subsection{Historical note}
\label{sec:ch44-historical}

The phenomenon exploited here has a long prehistory in commutative algebra
and almost none in pluripotential theory.  Seminormality was introduced in
the 1960s and 1970s to name the largest quotient of the normalization that
does not change the topology or the residue fields; the complementary
notion, subintegrality, names extensions that change neither.  The
striking discovery of the 1990s --- that over a \(\mathbf Q\)-algebra a
subintegral extension admits an additive characterization, in which a single
element is corrected by finitely many terms with coefficients independent of
the exponent --- was made in the course of computing Picard groups of
subintegral extensions, and had, so far as we are aware, no analytic
application before the argument recorded in this chapter.

The reason it applies is worth stating plainly, since it is the kind of
coincidence that is easy to miss.  A Fubini--Study function is a maximum of
logarithms of sections, so the natural algebraic object attached to it is an
\emph{ideal}, and the natural numerical object is the value of that ideal at
a semivaluation.  Values of ideals are insensitive to everything except the
ideal's integral closure, and are exactly additive in ideal powers.  The
Roberts--Singh identity, read as a statement about ideals rather than
elements, produces a downstairs ideal squeezed between the \(n\)-th and the
\((n+C)\)-th power of an upstairs ideal.  Because values are additive in
powers, the squeeze becomes a numerical estimate with relative error
\(C/(n+C)\), and because \(C\) is independent of \(n\) the error vanishes.
Nothing analogous happens for individual functions, which is why no
statement of this chapter is an identity.

The contrast with characteristic \(p\) is instructive in the opposite
direction.  There the descent is exact, requires no limit, and produces a
bijection at the level of Fubini--Study presentations
(\Cref{cor:ch43-gf-bijection}); but the mechanism --- Frobenius --- is
available only in characteristic \(p\), and its uniformity has to be won by
a separate localization argument
(\Cref{cor:ch43-frobenius-all-affines}).  Here the mechanism is available in
all characteristic zero at once, and its uniformity is automatic once the
additive-power identity is applied to the correct element; what has to be
won instead is the passage from an estimate to a limit, and with it the
manufactured monotonicity of \Cref{sec:ch44-monotonization} and the single
twist of \Cref{lem:ch44-uniform-cutoffs}.  Neither proof specializes to the
other, and we know of no argument covering both.

Finally, a word on what the chapter delivers to the rest of the book.  With
\Cref{cor:bj-universal-reflection-all-characteristics} in hand, the passage
from a normal variety to a unibranch one is a pointwise identification of
psh spaces rather than a transfer of a property, so every statement about
\(\PSH_{X^\nu}\) --- order, finite maxima, decreasing limits, and
boundedness --- transports pointwise, with the usual generic-finiteness
condition on decreasing limits.  Energy is not included in this formal
list: its transport requires a separate energy/intersection comparison.
That is more
than the envelope property alone would give, and it is the reason the
descent is formulated as a reflection theorem rather than as an implication
between envelope properties.
